% Sensibilisation sur la technique du génie génétique dans le cadre de l'amélioration des caractéristiques des organismes — SVT 1ère D — Cours signé OnBuch+
% Programme officiel MINESEC. Compiler avec : tectonic ND03-genie-genetique-amelioration-organismes.tex
\newcommand{\DOCMATIERE}{SVT}
\newcommand{\DOCNIVEAU}{1ère D}
\newcommand{\DOCMODULE}{Module 1 : Le Monde Vivant}
\newcommand{\DOCLECON}{ND03}
\newcommand{\DOCTITRE}{Sensibilisation sur la technique du génie génétique dans le cadre de l'amélioration des caractéristiques des organismes}
% OnBuch+ — préambule « cours signés » ENRICHI (compilé avec tectonic / XeLaTeX)
% Système visuel « Pop » d'OnBuch+ : fond crème (jamais blanc), bordures encre
% épaisses, ombres « dures » décalées, titres Archivo Black en capitales.
% Version MATHÉMATIQUES : graphes (pgfplots/tikz), boîtes pédagogiques
% (définition, propriété, méthode, attention, le savais-tu, exercice résolu,
% exercice, TP, bilan), page de garde et signature OnBuch+ ancrée en bas.
%
% Macros attendues dans main.tex AVANT \input{preamble} :
%   \DOCMATIERE  \DOCNIVEAU  \DOCMODULE  \DOCLECON (numéro)  \DOCTITRE
\documentclass[11pt]{article}

\usepackage{fontspec}
\usepackage[french]{babel}
\usepackage[a4paper,margin=2cm,top=3.4cm,bottom=2.8cm,headheight=1.4cm,footskip=1.3cm]{geometry}
\usepackage{amsmath,amssymb,amsfonts}
\usepackage{mathtools} % \xmapsto, \coloneqq... souvent utilisés par les modèles
\usepackage{cancel} % simplifications barrées (bilans de forces)
% \abs{z} (module, valeur absolue) et \norm{u} (norme), utilisés par les modèles
\providecommand{\abs}{}\let\abs\relax\DeclarePairedDelimiter{\abs}{\lvert}{\rvert}
\providecommand{\norm}{}\let\norm\relax\DeclarePairedDelimiter{\norm}{\lVert}{\rVert}
\usepackage[table]{xcolor} % [table] : fournit aussi \rowcolors (alternance de lignes)
\usepackage{pagecolor}
\usepackage{tikz}
\usetikzlibrary{arrows.meta,positioning,calc,shapes.geometric,shapes.misc,shapes.arrows,decorations.pathmorphing,decorations.pathreplacing,decorations.markings,patterns,shadows,mindmap,trees,fit,backgrounds,matrix,intersections,angles,quotes}
\usepackage{pgfplots}
\usepackage[version=4]{mhchem} % équations nucléaires \ce{^{235}_{92}U}
\usepackage[european,siunitx]{circuitikz} % schémas électriques
\pgfplotsset{compat=1.18}
\usepgfplotslibrary{fillbetween,groupplots} % aires entre courbes (calcul intégral)
\usepackage{enumitem}
\usepackage{fancyhdr}
% [most] charge toutes les bibliothèques tcolorbox (listings, minted, raster,
% documentation...) et gonfle inutilement la mémoire TeX sur un document long
% (~300 boîtes) : on ne charge que ce qui est réellement utilisé.
\usepackage[skins,breakable]{tcolorbox}
\usepackage{graphicx}
\usepackage{colortbl}
\usepackage{booktabs}
\usepackage{tabularx}
\usepackage{diagbox} % case d'en-tête diagonale des tableaux à double entrée
% Colonnes extensibles centrée / gauche / droite (C, L, R), souvent utilisées par les modèles
\newcolumntype{C}{>{\centering\arraybackslash}X}
\newcolumntype{L}{>{\raggedright\arraybackslash}X}
\newcolumntype{R}{>{\raggedleft\arraybackslash}X}
\usepackage{array}
\usepackage{multicol}
\usepackage{siunitx}
% parse-numbers=false : accepte aussi \SI{2,5 \times 10^{-3}}{...}
\sisetup{locale=FR,output-decimal-marker={,},group-minimum-digits=4,parse-numbers=false}
\DeclareSIUnit\ohm{\ensuremath{\symup{\Omega}}} % U+2126 absent de la police
\DeclareSIUnit\division{div} % réglages d'oscilloscope (ms/div, V/div)
\DeclareSIUnit\tr{tr} % tours (vitesse de rotation, ex. \SI{1500}{\tr\per\minute})
\DeclareSIUnit\ch{ch} % cheval-vapeur (puissance moteur, unité usuelle hors SI)
\DeclareSIUnit\franccfa{FCFA} % franc CFA (coûts, contexte camerounais)
\DeclareSIUnit\dioptre{\ensuremath{\delta}} % vergence d'une lentille (dioptrie)
\usepackage{qrcode}
% \degree : commande fréquemment utilisée par les modèles pour un angle ou une
% température en dehors de \SI{}{} (non fournie par les paquets chargés).
\providecommand{\degree}{^{\circ}}
% \qed (fin de démonstration), utilisé par les modèles sans amsthm
\providecommand{\qed}{\hfill$\square$}
% \barre : double barre des valeurs interdites dans un tableau de variations (tkz-tab)
\providecommand{\barre}{\ensuremath{\|}}
% environnement proof (amsthm non chargé) utilisé par les modèles
\ifdefined\proof\else\newenvironment{proof}[1][Démonstration]{\par\noindent\textit{#1.}\ }{\qed\par}\fi
\usepackage{lastpage}
\usepackage{hyperref}
\hypersetup{hidelinks}

% ============================================================
% Palette « Pop » OnBuch+ — mêmes hex que la classe P (plus_ui.dart)
% ============================================================
\definecolor{poporange}{HTML}{F59321}
\definecolor{popdark}{HTML}{E07A0C}
\definecolor{popcream}{HTML}{FFF6E8}
\definecolor{popink}{HTML}{1C1714}
\definecolor{poppurple}{HTML}{7A5AE0}
\definecolor{popgreen}{HTML}{2E9E6A}
\definecolor{poppink}{HTML}{E0457B}
\definecolor{popblue}{HTML}{2F7DE1}
\definecolor{popgold}{HTML}{C98A12}
\definecolor{popmuted}{HTML}{8A7F76}
\definecolor{popline}{HTML}{EADFD2}
\definecolor{popgray}{HTML}{8A7F76}
\colorlet{popgrey}{popgray}
% Alias tolérants (le modèle nomme parfois les couleurs différemment)
\colorlet{popwhite}{white}
\colorlet{popblack}{popink}
\colorlet{popred}{poppink}
\colorlet{popyellow}{popgold}
% Teintes claires pour fonds de tableaux / zones
\colorlet{popblueL}{popblue!10!white}
\colorlet{poporangeL}{poporange!14!white}
\colorlet{popgreenL}{popgreen!12!white}
\colorlet{poppurpleL}{poppurple!12!white}
\colorlet{poppinkL}{poppink!12!white}
\colorlet{popgoldL}{popgold!14!white}

\pagecolor{popcream}

% ============================================================
% Typographie
% ============================================================
\defaultfontfeatures{Ligatures=TeX}
\setmainfont{PlusJakartaSans-Medium.ttf}[Path=../../../fonts/,BoldFont=PlusJakartaSans-Bold.ttf]
% Maths en sans-serif (Fira Math) avec chiffres et lettres droites de Plus
% Jakarta, pour que les formules se fondent dans le texte.
\usepackage[math-style=ISO,bold-style=ISO]{unicode-math}
\setmathfont{FiraMath-Regular.otf}[Path=../../../fonts/]
\setmathfont{PlusJakartaSans-Medium.ttf}[Path=../../../fonts/,range={up/num,up/{Latin,latin},\mathup}]
\usepackage{icomma} % virgule décimale sans espace en mode math
% Caractères mathématiques Unicode absents des polices de texte (Plus Jakarta,
% Archivo Black) : rendus par leur équivalent mathématique, en texte comme en
% formule — sinon ils s'affichent en carré vide (ex. « Arithmétique dans ℤ »).
\usepackage{newunicodechar}
\newunicodechar{ℝ}{\ensuremath{\mathbb{R}}}
\newunicodechar{ℤ}{\ensuremath{\mathbb{Z}}}
\newunicodechar{ℂ}{\ensuremath{\mathbb{C}}}
\newunicodechar{ℕ}{\ensuremath{\mathbb{N}}}
\newunicodechar{ℚ}{\ensuremath{\mathbb{Q}}}
\newunicodechar{𝔻}{\ensuremath{\mathbb{D}}}
\newunicodechar{α}{\ensuremath{\alpha}}
\newunicodechar{β}{\ensuremath{\beta}}
\newunicodechar{γ}{\ensuremath{\gamma}}
\newunicodechar{δ}{\ensuremath{\delta}}
\newunicodechar{ε}{\ensuremath{\varepsilon}}
\newunicodechar{θ}{\ensuremath{\theta}}
\newunicodechar{λ}{\ensuremath{\lambda}}
\newunicodechar{μ}{\ensuremath{\mu}}
\newunicodechar{σ}{\ensuremath{\sigma}}
\newunicodechar{τ}{\ensuremath{\tau}}
\newunicodechar{φ}{\ensuremath{\varphi}}
\newunicodechar{ψ}{\ensuremath{\psi}}
\newunicodechar{ω}{\ensuremath{\omega}}
\newunicodechar{Σ}{\ensuremath{\Sigma}}
% majuscules produites par \MakeUppercase dans les titres (α-aminés -> Α)
\newunicodechar{Α}{\ensuremath{\alpha}}
\newunicodechar{Β}{\ensuremath{\beta}}
\newunicodechar{Γ}{\ensuremath{\Gamma}}
\newunicodechar{Δ}{\ensuremath{\Delta}}
\newunicodechar{Ε}{\ensuremath{\varepsilon}}
\newunicodechar{Θ}{\ensuremath{\Theta}}
\newunicodechar{Λ}{\ensuremath{\Lambda}}
\newunicodechar{Μ}{\ensuremath{\mu}}
\newunicodechar{Τ}{\ensuremath{\tau}}
\newunicodechar{Φ}{\ensuremath{\Phi}}
\newunicodechar{Ψ}{\ensuremath{\Psi}}
\newunicodechar{Ω}{\ensuremath{\Omega}}
\newunicodechar{↦}{\ensuremath{\mapsto}}
\newunicodechar{∈}{\ensuremath{\in}}
\newunicodechar{∉}{\ensuremath{\notin}}
\newunicodechar{∧}{\ensuremath{\wedge}}
\newunicodechar{⇔}{\ensuremath{\Leftrightarrow}}
\newunicodechar{⟺}{\ensuremath{\Longleftrightarrow}}
\newunicodechar{⇒}{\ensuremath{\Rightarrow}}
\newunicodechar{⟹}{\ensuremath{\Longrightarrow}}
\newunicodechar{∘}{\ensuremath{\circ}}
\newunicodechar{∪}{\ensuremath{\cup}}
\newunicodechar{∩}{\ensuremath{\cap}}
\newunicodechar{≡}{\ensuremath{\equiv}}
\newunicodechar{⊂}{\ensuremath{\subset}}
\newunicodechar{⊥}{\ensuremath{\perp}}
\newunicodechar{∃}{\ensuremath{\exists}}
\newunicodechar{∀}{\ensuremath{\forall}}
\newunicodechar{∛}{\ensuremath{\sqrt[3]{\,}}}
\newunicodechar{∜}{\ensuremath{\sqrt[4]{\,}}}
\newunicodechar{⃗}{\ensuremath{{}^{\to}}}
\newunicodechar{ⁿ}{\ifmmode^{n}\else\textsuperscript{\ensuremath{n}}\fi}
\newunicodechar{ˣ}{\ifmmode^{x}\else\textsuperscript{\ensuremath{x}}\fi}
\newunicodechar{ᵇ}{\ifmmode^{b}\else\textsuperscript{\ensuremath{b}}\fi}
\newunicodechar{ʳ}{\ifmmode^{r}\else\textsuperscript{\ensuremath{r}}\fi}
\newunicodechar{ᵘ}{\ifmmode^{u}\else\textsuperscript{\ensuremath{u}}\fi}
\newunicodechar{ʸ}{\ifmmode^{y}\else\textsuperscript{\ensuremath{y}}\fi}
\newunicodechar{ᵃ}{\ifmmode^{a}\else\textsuperscript{\ensuremath{a}}\fi}
\newunicodechar{ᵉ}{\ifmmode^{e}\else\textsuperscript{\ensuremath{e}}\fi}
\newunicodechar{ᵏ}{\ifmmode^{k}\else\textsuperscript{\ensuremath{k}}\fi}
\newunicodechar{ᵖ}{\ifmmode^{p}\else\textsuperscript{\ensuremath{p}}\fi}
\newunicodechar{ᴺ}{\ifmmode^{N}\else\textsuperscript{\ensuremath{N}}\fi}
\newunicodechar{ᶜ}{\ifmmode^{c}\else\textsuperscript{\ensuremath{c}}\fi}
\newunicodechar{ʰ}{\ifmmode^{h}\else\textsuperscript{\ensuremath{h}}\fi}
\newunicodechar{ᵠ}{\ifmmode^{q}\else\textsuperscript{\ensuremath{q}}\fi}
\newunicodechar{ₙ}{\ifmmode_{n}\else\textsubscript{\ensuremath{n}}\fi}
\newunicodechar{ₐ}{\ifmmode_{a}\else\textsubscript{\ensuremath{a}}\fi}
\newunicodechar{ᵢ}{\ifmmode_{i}\else\textsubscript{\ensuremath{i}}\fi}
\newunicodechar{ᵣ}{\ifmmode_{r}\else\textsubscript{\ensuremath{r}}\fi}
\newunicodechar{ₖ}{\ifmmode_{k}\else\textsubscript{\ensuremath{k}}\fi}
\newunicodechar{ᵦ}{\ifmmode_{\beta}\else\textsubscript{\ensuremath{\beta}}\fi}
\newunicodechar{ₓ}{\ifmmode_{x}\else\textsubscript{\ensuremath{x}}\fi}
\newfontfamily\popdisplay{ArchivoBlack-Regular.ttf}[Path=../../../fonts/]
\newfontfamily\poplogo{SpaceGrotesk-Bold.ttf}[Path=../../../fonts/]
\newfontfamily\popsemi{PlusJakartaSans-SemiBold.ttf}[Path=../../../fonts/]
\newfontfamily\popxbold{PlusJakartaSans-ExtraBold.ttf}[Path=../../../fonts/]
\newfontfamily\popxboldls{PlusJakartaSans-ExtraBold.ttf}[Path=../../../fonts/,LetterSpace=3.0]

\setlist[enumerate]{leftmargin=1.7em,itemsep=5pt,topsep=4pt}
\setlist[itemize]{leftmargin=1.7em,itemsep=4pt,topsep=4pt,label={\color{poporange}\textbullet}}
\setlength{\parindent}{0pt}
\setlength{\parskip}{5pt}
\renewcommand{\arraystretch}{1.25}

% pgfplots : style Pop par défaut
\pgfplotsset{
  every axis/.append style={axis line style={popink,line width=1pt},tick style={popink},
    grid style={popline,line width=0.5pt},label style={font=\small},tick label style={font=\footnotesize}},
  popcurve/.style={poporange,line width=1.8pt,no markers,smooth},
}
% Même style disponible dans un \draw TikZ simple (hors pgfplots)
\tikzset{popcurve/.style={poporange,line width=1.8pt}}
\tikzset{popgrid/.style={popline,very thin}}
\colorlet{popcurve}{poporange} % le nom du style est parfois utilisé comme couleur

% ============================================================
% Pill / squiggle
% ============================================================
\newcommand{\poppill}[3]{%
  \tikz[baseline=-0.55ex]\node[draw=popink,line width=1.2pt,rounded corners=2.6pt,%
    fill=#2,inner xsep=6.5pt,inner ysep=3pt]{%
    \popxboldls\fontsize{7}{7}\selectfont\color{#3}\MakeUppercase{#1}};%
}
\newcommand{\popsquiggle}[1]{%
  \tikz[baseline=-0.4ex]{
    \draw[#1,line width=2.6pt,line cap=round]
      plot [smooth] coordinates {(0,0.09) (0.34,0.01) (0.68,0.21) (1.02,0.01) (1.36,0.10)};
  }%
}
% Mot-clé surligné (marqueur orange sous le texte)
\newcommand{\cle}[1]{{\popxbold\color{popdark}#1}}

% ============================================================
% En-tête / pied
% ============================================================
\pagestyle{fancy}
\fancyhf{}
\renewcommand{\headrulewidth}{0pt}
\renewcommand{\footrulewidth}{0pt}
\lhead{\raisebox{-0.15ex}{\poplogo\fontsize{13}{13}\selectfont\color{popink}OnBuch\color{poporange}+}}
\rhead{\poppill{\DOCMATIERE\ \textbullet\ \DOCNIVEAU\ \textbullet\ Leçon \DOCLECON}{white}{popink}}
\lfoot{\popxbold\fontsize{7.6}{7.6}\selectfont\color{popmuted}\MakeUppercase{Cours signé OnBuch+}\ \textbullet\ {\popsemi\DOCTITRE}}
\rfoot{\tikz[baseline=-0.6ex]\node[rounded corners=8.5pt,fill=popink,minimum height=17pt,minimum width=17pt,inner xsep=5pt,inner sep=0]{\popxbold\fontsize{8}{8}\selectfont\color{popcream}\thepage\,/\,\pageref*{LastPage}};}

% ============================================================
% Titres
% ============================================================
\newcommand{\chaptitre}[2]{%
  \begin{tcolorbox}[enhanced,colback=poporange,colframe=popink,boxrule=2.6pt,arc=11pt,
      drop shadow=popink,left=15pt,right=15pt,top=12pt,bottom=11pt,before skip=3pt,after skip=18pt]
    \poppill{#2}{popink}{popcream}\par\vspace{9pt}
    {\raggedright\hyphenpenalty=10000\exhyphenpenalty=10000\popdisplay\fontsize{22}{24}\selectfont\color{popcream}\MakeUppercase{#1}\par}
  \end{tcolorbox}
}
% Section numérotée : pastille orange avec le numéro + titre Archivo
\newcounter{popsec}
\newcommand{\coursec}[1]{%
  \stepcounter{popsec}\setcounter{popsub}{0}%
  \par\vspace{16pt}\noindent
  \begin{minipage}{\linewidth}
  \tikz[baseline=-0.7ex]\node[draw=popink,line width=1.6pt,fill=poporange,rounded corners=4pt,minimum size=22pt,inner sep=0]{\popdisplay\fontsize{12}{12}\selectfont\color{popcream}\thepopsec};%
  \hspace{8pt}{\popdisplay\fontsize{15}{17}\selectfont\color{popink}\MakeUppercase{#1}}\par
  \vspace{4pt}\noindent{\color{poporange}\rule{36pt}{3.6pt}}
  \end{minipage}\par\nopagebreak\vspace{8pt}
}
\newcounter{popsub}
\newcommand{\courssub}[1]{%
  \stepcounter{popsub}%
  \par\vspace{9pt}\noindent{\popxbold\fontsize{11.5}{13}\selectfont\color{popink}{\color{poporange}\thepopsec.\thepopsub}\hspace{6pt}#1}\par\nopagebreak\vspace{4pt}
}

% ============================================================
% Boîtes pédagogiques — même recette : fond blanc, bordure épaisse colorée,
% ombre dure encre, badge plein. Titre optionnel [#1].
% ============================================================
\tcbset{popbase/.style={breakable,enhanced,colback=white,boxrule=2.3pt,arc=9pt,
  shadow={3pt}{-3pt}{0pt}{popink},left=14pt,right=14pt,top=11pt,bottom=11pt,before skip=11pt,after skip=15pt,
  fontupper=\popsemi\fontsize{10.6}{14.5}\selectfont\color{popink},
  fontlower=\popsemi\fontsize{10.6}{14.5}\selectfont\color{popink}}}
% Coche dessinée (le glyphe ✓ n'existe pas dans les polices du document)
\newcommand{\popcheck}[1]{\tikz[baseline=-0.5ex]\draw[#1,line width=1.6pt,line cap=round,line join=round](-0.13,0.0)--(-0.04,-0.09)--(0.13,0.1);}
\providecommand{\coche}{\popcheck{popgreen}}
\providecommand{\resultat}[1]{\par\smallskip\noindent\fcolorbox{popgreen}{popgreenL}{\parbox{\dimexpr\linewidth-2\fboxsep-2\fboxrule\relax}{\textbf{Résultat.} #1}}\par\smallskip}
\newcommand{\popboxhead}[3]{\poppill{#1}{#2}{white}\ifx\relax#3\relax\else\hspace{6pt}{\popxbold\fontsize{10.6}{12}\selectfont\color{popink}#3}\fi\par\vspace{7pt}}

\newtcolorbox{aretenir}[1][]{popbase,colframe=popgreen,before upper={\popboxhead{À retenir}{popgreen}{#1}}}
\newtcolorbox{exemplebox}[1][]{popbase,colframe=poporange,before upper={\popboxhead{Exemple}{poporange}{#1}}}
\newtcolorbox{definition}[1][]{popbase,colframe=popblue,before upper={\popboxhead{Définition}{popblue}{#1}}}
\newtcolorbox{propriete}[1][]{popbase,colframe=poppurple,before upper={\popboxhead{Propriété}{poppurple}{#1}}}
\newtcolorbox{methode}[1][]{popbase,colframe=popgold,before upper={\popboxhead{Méthode}{popgold}{#1}}}
\newtcolorbox{attention}[1][]{popbase,colframe=poppink,before upper={\popboxhead{Attention}{poppink}{#1}}}
\newtcolorbox{experience}[1][]{popbase,colframe=popblue,colback=popblueL,before upper={\popboxhead{Expérience}{popblue}{#1}}}
% « Le savais-tu ? » : carte violette pleine (contexte, culture, Cameroun)
\newtcolorbox{savaistu}[1][]{popbase,colback=poppurple,colframe=popink,
  fontupper=\popsemi\fontsize{10.4}{14.2}\selectfont\color{white},
  before upper={\poppill{Le savais-tu\,?}{white}{poppurple}\ifx\relax#1\relax\else\hspace{6pt}{\popxbold\color{white}#1}\fi\par\vspace{7pt}}}
% Exercice résolu : énoncé en haut, \tcblower puis la solution sur fond crème
\newcounter{popexo}\newcounter{popexr}
\newtcolorbox{exoresolu}[1][]{popbase,colframe=poporange,colbacklower=popcream,
  segmentation style={popink,line width=1.2pt,dashed},
  before upper={\stepcounter{popexr}\popboxhead{Exercice résolu \thepopexr}{poporange}{#1}},
  before lower={\poppill{Solution}{popink}{popcream}\par\vspace{6pt}}}
% Exercice (non résolu en place) — la correction est dans la partie Corrigés
\newtcolorbox{exercice}[1][]{popbase,colframe=popink,
  before upper={\stepcounter{popexo}\popboxhead{Exercice \thepopexo}{popink}{#1}}}
% Corrigé
\newtcolorbox{corrige}[1][]{popbase,colframe=popgreen,colback=popgreenL,
  before upper={\popboxhead{Corrigé}{popgreen}{#1}}}
% Objectifs / prérequis / bilan
\newtcolorbox{objectifs}{popbase,colframe=popink,colback=poporangeL,
  before upper={\popboxhead{Objectifs de la leçon}{popink}{}}}
\newtcolorbox{prerequis}{popbase,colframe=popink,colback=popblueL,
  before upper={\popboxhead{Prérequis}{popblue}{}}}
\newtcolorbox{bilan}{popbase,colframe=popink,colback=white,boxrule=2.8pt,
  before upper={\popboxhead{Fiche bilan}{poporange}{L'essentiel à connaître pour l'examen}}}
% Figure encadrée avec légende
\newtcolorbox{popfigure}[1][]{enhanced,breakable=false,colback=white,colframe=popline,boxrule=1.4pt,arc=8pt,
  left=10pt,right=10pt,top=10pt,bottom=8pt,before skip=10pt,after skip=12pt,halign=center,
  lower separated=false,halign lower=center,
  fontlower=\popsemi\fontsize{9}{11.5}\selectfont\color{popmuted},
  #1}
\newcommand{\legende}[1]{\par\vspace{6pt}{\popsemi\fontsize{9}{11.5}\selectfont\color{popmuted}#1}}

% ============================================================
% Signature OnBuch+
% ============================================================
\newcommand{\poplogoinline}[1]{{\poplogo\fontsize{#1}{#1}\selectfont\color{popink}OnBuch\color{poporange}+}}
\newcommand{\signatureonbuch}{%
  \begin{tcolorbox}[enhanced,colback=popink,colframe=popink,boxrule=2.2pt,arc=10pt,
      drop shadow=poporange,left=14pt,right=14pt,top=10pt,bottom=10pt,before skip=0pt,after skip=0pt]
    \tikz[baseline=-0.7ex]\node[circle,fill=popgreen,draw=popcream,line width=1.3pt,minimum size=22pt,inner sep=0]{\popcheck{white}};%
    \hspace{9pt}\begin{minipage}[c]{0.62\linewidth}
      {\popxbold\fontsize{12}{13}\selectfont\color{popcream}Signé\ \textbullet\ {\poplogo OnBuch\color{poporange}+}}\par\vspace{2pt}
      {\raggedright\popsemi\fontsize{8}{10.5}\selectfont\color{popline}\MakeUppercase{Équipe pédagogique OnBuch+}\\\MakeUppercase{Conforme au programme officiel MINESEC}\par}
    \end{minipage}\hfill
    {\poplogo\fontsize{22}{22}\selectfont\color{popcream}OnBuch\color{poporange}+}
  \end{tcolorbox}%
}

% ============================================================
% Page de garde
% #1 titre · #2 matière · #3 niveau · #4 module · #5 n° leçon · #6 durée ·
% #7 description / objectifs (texte ou itemize)
% ============================================================
\newcommand{\pagegarde}[7]{%
  \thispagestyle{empty}
  \newgeometry{a4paper,margin=1.7cm,top=1.6cm,bottom=1.4cm}%
  \begin{tikzpicture}[remember picture,overlay]
    \fill[poporange!18] ([xshift=-3.2cm,yshift=-3.6cm]current page.north east) circle (4.6cm);
    \fill[poppurple!14] ([xshift=2.4cm,yshift=7.6cm]current page.south west) circle (3.8cm);
  \end{tikzpicture}%
  {\poplogo\fontsize{30}{30}\selectfont\color{popink}OnBuch\color{poporange}+}\hfill
  \raisebox{0.6ex}{\poppill{Cours signé \textbullet\ Premium}{popink}{popcream}}\par
  \vspace{1pt}\popsquiggle{poppurple}\par
  \vspace{8pt}
  \begin{tcolorbox}[enhanced,colback=poporange,colframe=popink,boxrule=2.8pt,arc=13pt,
      drop shadow=popink,left=18pt,right=18pt,top=12pt,bottom=13pt,after skip=10pt]
    \poppill{#2 \textbullet\ #3}{popink}{popcream}\hspace{5pt}\poppill{#4}{white}{popink}\par\vspace{8pt}
    {\popdisplay\fontsize{14}{14}\selectfont\color{popink}LEÇON #5}\par\vspace{4pt}
    {\raggedright\hyphenpenalty=10000\exhyphenpenalty=10000\popdisplay\fontsize{25}{27}\selectfont\color{popcream}\MakeUppercase{#1}\par}
    \vspace{8pt}
    {\popxbold\fontsize{9.5}{9.5}\selectfont\color{popink}\MakeUppercase{Durée conseillée : #6}}
  \end{tcolorbox}
  \begin{tcolorbox}[enhanced,colback=white,colframe=popink,boxrule=2.2pt,arc=10pt,
      drop shadow=popink,left=16pt,right=16pt,top=10pt,bottom=10pt,after skip=8pt]
    \poppill{Dans cette leçon}{poppurple}{white}\par\vspace{5pt}
    \popsemi\fontsize{9.3}{12.6}\selectfont\color{popink}#7
  \end{tcolorbox}
  \noindent\begin{minipage}[t]{0.487\linewidth}
    \begin{tcolorbox}[enhanced,colback=popgreen,colframe=popink,boxrule=2.2pt,arc=10pt,
        drop shadow=popink,left=11pt,right=11pt,top=10pt,bottom=10pt,height=3.1cm,valign=center]
      \tcbox[colback=white,colframe=popink,boxrule=1.3pt,arc=5pt,boxsep=0pt,nobeforeafter,
        left=3pt,right=3pt,top=3pt,bottom=3pt,valign=center]{\qrcode[height=1.9cm]{https://play.google.com/store/apps/details?id=com.luvvix.onbuch&hl=fr}}%
      \hspace{8pt}\begin{minipage}[c]{0.5\linewidth}
        {\popdisplay\fontsize{11}{12.5}\selectfont\color{white}\MakeUppercase{Télécharge OnBuch}}\par\vspace{4pt}
        {\raggedright\popsemi\fontsize{8.4}{11}\selectfont\color{white}Gratuit \textbullet\ Google Play\\Tuteur IA, annales, concours, résultats\par}
      \end{minipage}
    \end{tcolorbox}
  \end{minipage}\hfill
  \begin{minipage}[t]{0.487\linewidth}
    \begin{tcolorbox}[enhanced,colback=poppurple,colframe=popink,boxrule=2.2pt,arc=10pt,
        drop shadow=popink,left=11pt,right=11pt,top=10pt,bottom=10pt,height=3.1cm,valign=center]
      \tikz[baseline=-0.7ex]\node[circle,fill=white,draw=popink,line width=1.3pt,minimum size=1.6cm,inner sep=0]{\popdisplay\fontsize{24}{24}\selectfont\color{poppurple}+};%
      \hspace{8pt}\begin{minipage}[c]{0.6\linewidth}
        {\popdisplay\fontsize{11}{12.5}\selectfont\color{white}\MakeUppercase{Passe à OnBuch+}}\par\vspace{4pt}
        {\raggedright\popsemi\fontsize{8.4}{11}\selectfont\color{white}Tous les cours signés, Tuteur IA illimité, fiches PDF — onglet OnBuch+ de l'app\par}
      \end{minipage}
    \end{tcolorbox}
  \end{minipage}
  \vspace{8pt}
  \popcontenu
  \vfill
  \signatureonbuch
  \restoregeometry
  \newpage
}

% Rangée « Ce cours contient » (page de garde)
\newcommand{\poptile}[3]{% #1 couleur #2 gros texte #3 légende
  \begin{tcolorbox}[enhanced,colback=#1,colframe=popink,boxrule=2pt,arc=9pt,shadow={2.5pt}{-2.5pt}{0pt}{popink},
      left=6pt,right=6pt,top=9pt,bottom=9pt,halign=center,valign=center,height=1.75cm,width=\linewidth,nobeforeafter]
    {\popdisplay\fontsize{12.5}{13}\selectfont\color{white}\MakeUppercase{#2}}\par\vspace{3pt}
    {\centering\popsemi\fontsize{7.8}{9.5}\selectfont\color{white}#3\par}
  \end{tcolorbox}}
\newcommand{\popcontenu}{%
  \par\noindent{\popxboldls\fontsize{7.5}{7.5}\selectfont\color{popmuted}CE COURS SIGNÉ CONTIENT}\par\vspace{7pt}
  \noindent\begin{minipage}[t]{0.235\linewidth}\poptile{popblue}{Cours}{complet, illustré, pas à pas}\end{minipage}\hfill
  \begin{minipage}[t]{0.235\linewidth}\poptile{poporange}{Méthodes}{et exercices résolus}\end{minipage}\hfill
  \begin{minipage}[t]{0.235\linewidth}\poptile{poppink}{Exercices}{type examen + corrigés}\end{minipage}\hfill
  \begin{minipage}[t]{0.235\linewidth}\poptile{popgreen}{Fiche bilan}{l'essentiel à retenir}\end{minipage}\par}

% Sommaire
\newtcolorbox{sommaire}{enhanced,colback=white,colframe=popink,boxrule=2.2pt,arc=10pt,
  drop shadow=popink,left=18pt,right=18pt,top=13pt,bottom=13pt,before skip=6pt,after skip=20pt,
  before upper={\poppill{Sommaire}{poppurple}{white}\par\vspace{9pt}\popsemi\fontsize{10.6}{14.5}\selectfont\color{popink}}}

% Signature de fin de cours — poussée tout en bas de la dernière page
\newcommand{\findecours}{%
  \par\vfill\nopagebreak
  \begin{center}\popsquiggle{poporange}\end{center}\vspace{4pt}
  \signatureonbuch
}
\AtBeginDocument{\def\llbracket{\mathopen{[\mkern-3.5mu[}}\def\rrbracket{\mathclose{]\mkern-3.5mu]}}} % intervalles d'entiers (glyphe ⟦ absent de la police)
% Glyphes absents de Fira Math : redéfinis à partir de glyphes disponibles
\AtBeginDocument{%
  \def\setminus{\mathbin{\backslash}}\let\smallsetminus\setminus
  \def\vdots{\mathord{\vbox{\baselineskip3.2pt\lineskiplimit0pt\kern4pt\hbox{.}\hbox{.}\hbox{.}}}}%
  \def\ddots{\mathinner{\mkern1mu\raise6.5pt\hbox{.}\mkern2mu\raise3.6pt\hbox{.}\mkern2mu\raise0.7pt\hbox{.}\mkern1mu}}%
  \def\triangle{\mathord{\scalebox{0.9}{$\symup{\Delta}$}}}%
  \def\nabla{\mathord{\raisebox{\depth}{\scalebox{1}[-1]{$\symup{\Delta}$}}}}%
  \def\checkmark{\popcheck{popdark}}%
  \def\star{\mathbin{\ast}}\def\bigstar{\mathord{\ast}}%
  \def\diamond{\mathbin{\scalebox{0.6}{$\Diamond$}}}%
  \def\bigcup{\mathop{\mathchoice{\raisebox{-0.3em}{\scalebox{1.7}{$\cup$}}}{\scalebox{1.25}{$\cup$}}{\cup}{\cup}}\displaylimits}%
  \def\bigcap{\mathop{\mathchoice{\raisebox{-0.3em}{\scalebox{1.7}{$\cap$}}}{\scalebox{1.25}{$\cap$}}{\cap}{\cap}}\displaylimits}%
  \def\lesseqqgtr{\mathrel{\lessgtr}}\def\bigoplus{\mathop{\oplus}\displaylimits}%
}
\providecommand{\ssi}{si et seulement si} % abréviation fréquente
\AtBeginDocument{\providecommand{\wideparen}{\overparen}} % arc de cercle

%% FIGFIX-BEGIN v3 — correctifs globaux des figures (docs/onbuch-generation/tools/figfix)
\usepackage{adjustbox}\usepackage{etoolbox}
% toute figure TikZ/pgfplots plus large que la ligne est réduite à la largeur disponible
\BeforeBeginEnvironment{tikzpicture}{\begin{adjustbox}{max width=\linewidth}}
\AfterEndEnvironment{tikzpicture}{\end{adjustbox}}
% légendes pgfplots lisibles par-dessus les courbes
\pgfplotsset{every axis/.append style={legend style={fill=white,fill opacity=0.92,text opacity=1,draw=black!15,inner sep=3pt}}}
% étiquettes d'arêtes (syntaxe edge["..."]) avec fond blanc
\tikzset{every edge quotes/.append style={fill=white,inner sep=1.2pt}}
% bulles de cartes mentales : texte à la ligne dans la bulle, bulle un peu plus grande
\tikzset{every concept/.append style={text width=2.0cm,align=center,minimum size=2.3cm,inner sep=2pt}}
%% FIGFIX-END
\begin{document}

\pagegarde{Sensibilisation sur la technique du génie génétique dans le cadre de l'amélioration des caractéristiques des organismes}{SVT}{1ère D}{Module 1 : Le Monde Vivant}{ND03}{10 h}{Cette leçon présente les principes et les étapes du génie génétique, technique permettant de modifier le patrimoine génétique d'un organisme pour améliorer ses caractéristiques. À travers des exemples concrets (insuline humaine, amélioration variétale), elle sensibilise aux applications, aux enjeux et aux débats suscités par les organismes génétiquement modifiés (OGM).
\vspace{4pt}
\begin{multicols}{2}\small
\begin{itemize}
\item À la fin de cette leçon, je sais définir le génie génétique et un organisme génétiquement modifié (OGM).
\item À la fin de cette leçon, je sais décrire le rôle des enzymes de restriction, des ligases et des vecteurs de transfert.
\item À la fin de cette leçon, je sais décrire, dans l'ordre, les étapes d'obtention d'un OGM à partir d'un schéma.
\item À la fin de cette leçon, je sais expliquer comment des bactéries transgéniques produisent de l'insuline humaine.
\item À la fin de cette leçon, je sais citer des exemples d'applications du génie génétique en médecine et en agriculture (amélioration variétale).
\item À la fin de cette leçon, je sais discuter des avantages et des risques liés à l'utilisation des OGM.
\end{itemize}\end{multicols}}

\begin{sommaire}
\begin{multicols}{2}
\begin{enumerate}
\item Génie génétique et OGM : définitions et enjeux
\item Les outils du génie génétique
\item Étapes de l'obtention d'un OGM
\item Applications du génie génétique
\item Enjeux et débats autour des OGM
\item Activité d'intégration
\item Méthodes et exercices résolus
\item Exercices
\item Corrigés des exercices
\item Fiche bilan
\end{enumerate}
\end{multicols}
\end{sommaire}

\begin{objectifs}
\begin{itemize}
\item À la fin de cette leçon, je sais définir le génie génétique et un organisme génétiquement modifié (OGM).
\item À la fin de cette leçon, je sais décrire le rôle des enzymes de restriction, des ligases et des vecteurs de transfert.
\item À la fin de cette leçon, je sais décrire, dans l'ordre, les étapes d'obtention d'un OGM à partir d'un schéma.
\item À la fin de cette leçon, je sais expliquer comment des bactéries transgéniques produisent de l'insuline humaine.
\item À la fin de cette leçon, je sais citer des exemples d'applications du génie génétique en médecine et en agriculture (amélioration variétale).
\item À la fin de cette leçon, je sais discuter des avantages et des risques liés à l'utilisation des OGM.
\item À la fin de cette leçon, je sais construire un argumentaire scientifique équilibré sur les enjeux des biotechnologies génétiques.
\end{itemize}
\end{objectifs}

\begin{prerequis}
\begin{itemize}
\item Structure de l'ADN : nucléotides, double hélice, complémentarité des bases (Seconde).
\item Notion de gène et relation gène – protéine – caractère (Seconde).
\item Notion d'enzyme : protéine catalysant une réaction spécifique (Seconde).
\item Organisation du matériel génétique chez les bactéries : chromosome circulaire et plasmides (début d'année de Première).
\item Notion de mutation et de variabilité génétique (Première).
\end{itemize}
\end{prerequis}

\newpage

\coursec{Génie génétique et OGM : définitions et enjeux}

Au marché Mokolo de Yaoundé, une vendeuse de maïs affirme fièrement que son maïs « amélioré » résiste aux insectes sans aucun traitement insecticide. Quelques kilomètres plus loin, à l'hôpital central de Yaoundé, un jeune diabétique reçoit chaque jour une injection d'insuline. Or cette insuline, hormone humaine, n'a jamais été prélevée sur un être humain : elle est fabriquée... par des bactéries !

Comment de simples bactéries peuvent-elles produire une hormone humaine ? Comment obtient-on des plantes capables de résister à leurs ravageurs ? Ces deux prouesses reposent sur une même technologie : le \cle{génie génétique}. Mais ces organismes « modifiés » sont-ils sans danger pour notre santé et pour notre environnement ?

\textbf{Problématique :} Qu'est-ce que le génie génétique ? Qu'appelle-t-on un organisme génétiquement modifié (OGM) ? En quoi cette approche diffère-t-elle des méthodes classiques d'amélioration des organismes ?

\courssub{Qu'est-ce que le génie génétique ?}

\courssub{Une intuition d'abord : réparer ou transformer un « texte » vivant}

Imagine le génome d'un organisme comme un immense livre de recettes écrit dans une langue universelle : la langue de l'ADN, faite de quatre lettres (A, T, G, C). Chaque recette est un \cle{gène}, c'est-à-dire un fragment d'ADN qui porte l'information nécessaire à la fabrication d'une protéine. Le gène de l'insuline humaine est la « recette » de l'insuline ; un gène de résistance aux insectes est la « recette » d'une protéine toxique pour certains ravageurs.

Or, depuis les années 1970, les biologistes ont découvert des outils moléculaires capables de \emph{découper}, \emph{copier} et \emph{coller} des fragments d'ADN, comme on coupe et colle un paragraphe dans un texte. Il devient alors possible de prélever une « recette » chez un organisme et de l'insérer dans le livre d'un autre organisme, qui se mettra à fabriquer la protéine correspondante. C'est exactement ce que font les bactéries productrices d'insuline : on leur a « collé » la recette humaine de l'insuline, et elles l'exécutent fidèlement.

\begin{definition}[Le génie génétique]
Le \cle{génie génétique} est l'ensemble des techniques de manipulation \emph{in vitro} de l'ADN permettant de modifier volontairement le patrimoine génétique (le génome) d'un organisme : isolement d'un gène, transfert de ce gène, insertion et expression de ce gène dans un organisme receveur.
\end{definition}

Trois verbes résument l'essentiel du génie génétique :
\begin{itemize}
\item \textbf{isoler} : extraire et identifier le gène d'intérêt dans l'ADN de l'organisme donneur ;
\item \textbf{transférer} : introduire ce gène dans les cellules de l'organisme receveur, à l'aide d'un vecteur (un « véhicule » à ADN) ;
\item \textbf{insérer et faire exprimer} : intégrer le gène au génome du receveur de façon stable, afin qu'il soit lu, traduit en protéine et transmis aux descendants.
\end{itemize}

\begin{attention}[Ce que le génie génétique n'est pas]
Le génie génétique n'est pas une simple sélection ni un croisement entre variétés. C'est une \textbf{manipulation directe de la molécule d'ADN}, réalisée en laboratoire, hors de tout rapport sexuel. Dire « ce maïs hybride obtenu par croisement est un OGM » est une erreur classique au Probatoire : un hybride résulte d'une reproduction sexuée classique, pas d'une modification artificielle du génome.
\end{attention}

\courssub{Qu'est-ce qu'un OGM ?}

\courssub{Définition rigoureuse}

\begin{definition}[Organisme génétiquement modifié (OGM)]
Un \cle{OGM} (organisme génétiquement modifié) est un organisme vivant (bactérie, levure, plante, animal) dont le génome a été artificiellement modifié par les techniques du génie génétique, soit par \textbf{insertion d'un ou plusieurs gènes étrangers} (appelés \cle{transgènes}), soit par \textbf{modification d'un de ses propres gènes}.
\end{definition}

Lorsque l'OGM a reçu un gène provenant d'une \emph{autre espèce}, on parle d'\textbf{organisme transgénique}. C'est le cas le plus fréquent : la bactérie \emph{Escherichia coli} portant le gène humain de l'insuline, le maïs Bt portant un gène de la bactérie \emph{Bacillus thuringiensis}, le riz doré portant des gènes de synthèse du bêta-carotène.

\begin{attention}[Erreur fréquente à éviter absolument]
Un OGM n'est \textbf{jamais} un organisme issu de croisements ou de sélection classique, même si ces méthodes « modifient » les caractéristiques des variétés au fil des générations. Un organisme n'est un OGM que s'il résulte \textbf{obligatoirement d'une manipulation directe de l'ADN} (insertion d'un transgène ou modification ciblée d'un gène). Le maïs hybride de tes champs, le zébu amélioré par sélection, le manioc sélectionné pour son rendement : aucun de ces organismes n'est un OGM.
\end{attention}

\courssub{La transgénèse : franchir les barrières d'espèces}

\begin{definition}[Transgénèse]
La \cle{transgénèse} est le transfert d'un gène d'une espèce vers une autre espèce, y compris entre espèces très éloignées sur le plan de l'évolution (par exemple un gène humain introduit dans une bactérie, ou un gène bactérien introduit dans une plante).
\end{definition}

Pourquoi un tel transfert est-il possible ? Parce que le \textbf{code génétique est quasi universel} : les mêmes « mots » d'ADN (les codons) désignent les mêmes acides aminés chez la bactérie, la plante et l'Homme. Une bactérie qui reçoit le gène humain de l'insuline le « lit » exactement comme le lirait une cellule humaine, et fabrique donc la même protéine. Cette universalité du code génétique, découverte dans les années 1960, est le fondement théorique de toute la transgénèse.

\begin{popfigure}
\begin{center}
\begin{tikzpicture}[font=\small, >=Stealth]
% Donneur
\draw[fill=popblueL, rounded corners=4pt] (0,0) rectangle (4.2,2.6);
\node[font=\small\bfseries, popblue] at (2.1,2.25) {Organisme donneur};
\node[font=\footnotesize, popmuted] at (2.1,1.9) {(ex. cellule humaine)};
\draw[very thick, popdark] (0.5,1.2) -- (3.7,1.2);
\fill[poporange] (1.6,1.05) rectangle (2.6,1.35);
\node[font=\footnotesize, white] at (2.1,1.2) {gène d'intérêt};
\node[font=\footnotesize, popmuted] at (2.1,0.6) {ADN (génome)};
% Flèche prélèvement
\draw[->, very thick, poporange] (2.1,1.0) -- (2.1,-0.6) -- (6.3,-0.6) -- (6.3,0.0);
\node[font=\footnotesize, poporange, align=center] at (4.2,-0.95) {prélèvement puis\\ transfert du gène};
% Receveur
\draw[fill=popgreenL, rounded corners=4pt] (4.6,0) rectangle (8.8,2.6);
\node[font=\small\bfseries, popgreen!60!black] at (6.7,2.25) {Organisme receveur};
\node[font=\footnotesize, popmuted] at (6.7,1.9) {(ex. bactérie)};
\draw[very thick, popdark] (5.1,1.2) -- (8.3,1.2);
\fill[poporange] (6.2,1.05) rectangle (7.2,1.35);
\node[font=\footnotesize, white] at (6.7,1.2) {transgène};
\node[font=\footnotesize, popmuted] at (6.7,0.6) {génome modifié : OGM};
\end{tikzpicture}
\end{center}
\legende{Principe de la transgénèse : le gène d'intérêt (orange) est prélevé dans le génome de l'organisme donneur, puis inséré dans le génome de l'organisme receveur, qui devient un OGM capable d'exprimer la protéine correspondante.}
\end{popfigure}

\courssub{Amélioration classique ou génie génétique : quelles différences ?}

Depuis la domestication des plantes et des animaux, l'Homme améliore les espèces utiles. Le paysan camerounais qui conserve chaque année les grains des plus beaux épis de maïs pour les ressemer pratique la \textbf{sélection} ; l'agronome qui croise deux variétés de manioc pour réunir leurs qualités pratique l'\textbf{hybridation}. Le génie génétique poursuit le même objectif général — améliorer les caractéristiques des organismes — mais par une voie radicalement différente. Comparons rigoureusement les deux approches.

\begin{center}
\begin{tabularx}{\linewidth}{|l|X|X|}
\hline
\rowcolor{popblueL} \textbf{Critère} & \textbf{Amélioration classique} & \textbf{Génie génétique} \\
\hline
Principe & Sélection et croisements (reproduction sexuée) entre individus de la même espèce ou d'espèces proches & Manipulation directe de l'ADN : isolement, transfert et insertion d'un gène identifié \\
\hline
Précision & Faible : on transfère des milliers de gènes à la fois, sans les connaître tous & Élevée : on transfère un gène précis, de fonction connue \\
\hline
Barrières d'espèces & Infranchissables : impossible de croiser une bactérie avec une plante & Franchies : un gène peut passer d'une espèce à une autre, même très éloignée \\
\hline
Durée & Longue : plusieurs générations (souvent 10 à 15 ans pour une variété végétale) & Courte : quelques mois à quelques années en laboratoire \\
\hline
Résultat & Variété ou race améliorée (non OGM) & Organisme génétiquement modifié (OGM) \\
\hline
\end{tabularx}
\end{center}

\begin{popfigure}
\begin{center}
\begin{tikzpicture}[font=\footnotesize, >=Stealth]
% Voie classique
\node[font=\small\bfseries, popblue] at (3.2,3.1) {Amélioration classique (croisements)};
\draw[fill=popblueL, rounded corners=3pt] (0,2.2) rectangle (2.2,2.9);
\node at (1.1,2.55) {Variété A};
\draw[fill=popblueL, rounded corners=3pt] (4.2,2.2) rectangle (6.4,2.9);
\node at (5.3,2.55) {Variété B};
\draw[->, thick, popblue] (2.2,2.55) -- (4.2,2.55) node[midway, above]{croisement};
\draw[->, thick, popblue] (3.2,2.2) -- (3.2,1.5);
\draw[fill=popblue!15, rounded corners=3pt] (1.2,0.7) rectangle (5.2,1.5);
\node[align=center] at (3.2,1.1) {sélection des descendants\\ sur plusieurs générations};
\draw[->, thick, popblue] (3.2,0.7) -- (3.2,0.0);
\draw[fill=popblue!30, rounded corners=3pt] (1.7,-0.7) rectangle (4.7,0.0);
\node at (3.2,-0.35) {variété améliorée};
% Voie transgénèse
\node[font=\small\bfseries, poporange] at (10.2,3.1) {Transgénèse (génie génétique)};
\draw[fill=poporangeL, rounded corners=3pt] (8.0,2.2) rectangle (10.0,2.9);
\node[align=center] at (9.0,2.55) {gène cible\\ identifié};
\draw[fill=popgreenL, rounded corners=3pt] (10.6,2.2) rectangle (12.6,2.9);
\node[align=center] at (11.6,2.55) {organisme\\ receveur};
\draw[->, very thick, poporange] (10.0,2.55) -- (10.6,2.55) node[midway, above, font=\scriptsize]{insertion};
\draw[->, very thick, poporange] (11.6,2.2) -- (11.6,0.0) node[midway, right, align=left]{une seule\\ étape de transfert};
\draw[fill=poporange!35, rounded corners=3pt] (10.4,-0.7) rectangle (12.8,0.0);
\node at (11.6,-0.35) {OGM};
\end{tikzpicture}
\end{center}
\legende{Comparaison des deux voies d'amélioration : les croisements classiques exigent la sélection des descendants sur plusieurs générations, tandis que la transgénèse transfère directement un gène cible identifié dans le génome du receveur.}
\end{popfigure}

\courssub{Pourquoi recourir au génie génétique ? Objectifs et enjeux}

L'objectif général du génie génétique est d'\textbf{améliorer les caractéristiques des organismes} au service de la santé, de l'agriculture et de l'industrie. Trois grandes familles d'objectifs se dégagent :

\begin{itemize}
\item \textbf{Produire des substances utiles} : faire fabriquer par des micro-organismes des protéines médicalement précieuses (insuline, hormone de croissance, vaccins, enzymes industrielles). C'est l'exemple de l'insuline de l'hôpital central de Yaoundé.
\item \textbf{Améliorer les productions agricoles} : obtenir des plantes résistantes aux insectes ravageurs (maïs Bt), tolérantes à la sécheresse ou aux maladies, ou à meilleure valeur nutritionnelle (riz doré enrichi en provitamine A) — un enjeu majeur pour la sécurité alimentaire en Afrique.
\item \textbf{Augmenter les rendements et réduire les pertes} : limiter l'usage d'insecticides, diminuer les récoltes détruites par les ravageurs, adapter les cultures aux conditions locales.
\end{itemize}

\begin{savaistu}[1982 : la première insuline « bactérienne »]
Le premier médicament issu du génie génétique est l'\textbf{insuline humaine produite par la bactérie} \emph{Escherichia coli}, mise sur le marché en \textbf{1982} par l'entreprise Genentech. Auparavant, les diabétiques recevaient de l'insuline extraite de pancréas de porcs ou de bovins abattus, qui provoquait parfois des allergies et dont l'approvisionnement était limité. Grâce à la transgénèse, on dispose aujourd'hui d'une insuline strictement identique à l'insuline humaine, produite en quantité illimitée et à moindre coût. Des millions de diabétiques, y compris au Cameroun, en bénéficient chaque jour.
\end{savaistu}

Mais ces prouesses soulèvent aussi des questions légitimes : les protéines nouvelles sont-elles allergènes ? Les transgènes peuvent-ils se disséminer dans l'environnement ? Qui contrôle ces technologies ? Ces questions, qui alimentent un débat de société au Cameroun comme ailleurs, seront examinées en détail dans la dernière partie de la leçon. Retenons dès maintenant que le débat porte sur les \emph{usages} et les \emph{risques} des OGM, non sur la définition scientifique de la technique.

\begin{exoresolu}[Classer des organismes : OGM ou non ?]
Pour chacun des organismes suivants, indique s'il s'agit d'un OGM et justifie ta réponse : (a) un maïs hybride obtenu en croisant deux variétés locales à haut rendement ; (b) une bactérie \emph{E. coli} ayant reçu le gène humain de l'insuline ; (c) un maïs Bt portant un gène de \emph{Bacillus thuringiensis} qui le rend toxique pour les chenilles ; (d) un zébu sélectionné pendant vingt ans pour sa production laitière.
\tcblower
\textbf{Solution détaillée.}

Le critère décisif est l'existence d'une \textbf{manipulation directe de l'ADN} (insertion d'un transgène ou modification ciblée d'un gène).

(a) Le maïs hybride n'est \textbf{pas} un OGM : il résulte d'un croisement, c'est-à-dire d'une reproduction sexuée classique entre deux variétés de la même espèce. Aucune manipulation de l'ADN n'a eu lieu.

(b) La bactérie \emph{E. coli} portant le gène humain de l'insuline \textbf{est} un OGM (plus précisément un organisme transgénique) : un gène étranger, le transgène, a été inséré dans son génome par génie génétique.

(c) Le maïs Bt \textbf{est} un OGM : il a reçu par transgénèse un gène bactérien codant une protéine insecticide. Le franchissement de la barrière d'espèces (bactérie $\rightarrow$ plante) est la signature du génie génétique.

(d) Le zébu sélectionné n'est \textbf{pas} un OGM : la sélection sur vingt ans repose sur la reproduction et le choix des reproducteurs, sans aucune modification artificielle du génome.
\end{exoresolu}

\begin{aretenir}[L'essentiel de la section]
\begin{itemize}
\item Le \cle{génie génétique} est l'ensemble des techniques de manipulation \emph{in vitro} de l'ADN permettant de modifier volontairement le génome d'un organisme (isolement, transfert, insertion de gènes).
\item Un \cle{OGM} est un organisme dont le génome a été artificiellement modifié par génie génétique, par insertion d'un transgène ou modification d'un gène propre.
\item La \cle{transgénèse} permet de transférer un gène d'une espèce à une autre, même très éloignée, grâce à l'universalité du code génétique.
\item Par rapport à l'amélioration classique (sélection, croisements), le génie génétique est plus \textbf{précis}, plus \textbf{rapide} et \textbf{franchit les barrières d'espèces}.
\item Objectif général : améliorer les caractéristiques des organismes (rendement, résistance, production de substances utiles comme l'insuline).
\end{itemize}
\end{aretenir}

\begin{exercice}[À toi de jouer]
Le riz doré est un riz dont le génome contient deux gènes étrangers (l'un issu d'un narcissе, l'autre d'une bactérie) qui lui permettent de synthétiser du bêta-carotène, précurseur de la vitamine A. 1) Le riz doré est-il un OGM ? Justifie. 2) Quel avantage la transgénèse offre-t-elle ici par rapport à un croisement classique avec une plante riche en bêta-carotène ? 3) Cite un enjeu sanitaire de cette innovation pour les populations d'Afrique subsaharienne.
\end{exercice}

\begin{corrige}[Exercice : le riz doré]
1) Oui, le riz doré est un OGM transgénique : son génome a été artificiellement modifié par insertion de deux transgènes provenant d'espèces différentes.

2) Un croisement classique est impossible entre un riz, un narcisse et une bactérie : la barrière d'espèces est infranchissable par reproduction sexuée. La transgénèse permet de transférer directement et précisément les seuls gènes utiles, rapidement, sans introduire les milliers d'autres gènes des espèces donneuses.

3) La carence en vitamine A est une cause majeure de cécité infantile et de mortalité en Afrique subsaharienne ; un riz enrichi en bêta-carotène pourrait contribuer à réduire cette carence dans les populations dont l'alimentation de base est le riz.
\end{corrige}

\coursec{Les outils du génie génétique}

Dans la section précédente, nous avons défini le génie génétique comme l'ensemble des techniques permettant de modifier volontairement le patrimoine génétique d'un organisme pour obtenir un organisme génétiquement modifié (OGM). Nous avons vu que l'enjeu est immense : produire de l'insuline humaine à grande échelle, créer des variétés de maïs résistantes aux insectes, soigner certaines maladies génétiques. Mais une question pratique se pose immédiatement : \emph{comment fait-on, concrètement, pour prélever un gène dans une cellule et l'installer dans une autre ?} Le gène est une molécule microscopique, enfouie dans un chromosome, au cœur du noyau. On ne peut pas le manipuler avec une pince ou une aiguille. Il faut des outils moléculaires. Cette section présente les trois grandes familles d'outils du génie génétique : les \cle{enzymes de restriction} qui coupent l'ADN, les \cle{ligases} qui le soudent, et les \cle{vecteurs de transfert} qui le transportent.

\courssub{Les enzymes de restriction : les ciseaux moléculaires}

\textbf{Une observation pour commencer.} Au laboratoire de biotechnologie de l'Université de Yaoundé I, un chercheur veut isoler le gène humain de l'insuline pour l'introduire dans une bactérie. Ce gène est noyé parmi les dizaines de milliers de gènes du génome humain. Comment le découper précisément, sans détruire le gène lui-même ni le reste de l'ADN ? La solution ne vient pas de l'homme : elle vient des bactéries elles-mêmes.

\textbf{Démarche d'investigation.} Dans les années 1960-1970, les biologistes observent un phénomène étrange : certains virus bactériens (les bactériophages) infectent facilement certaines souches de bactéries, mais sont détruits par d'autres.
\begin{itemize}
\item \emph{Observation} : l'ADN du virus injecté dans la bactérie résistante est retrouvé découpé en fragments, alors que l'ADN de la bactérie elle-même reste intact.
\item \emph{Hypothèse} : la bactérie possède une enzyme capable de couper l'ADN étranger à des endroits précis, tout en protégeant son propre ADN.
\item \emph{Vérification expérimentale} : on extrait des protéines bactériennes et on les met en contact avec de l'ADN viral en éprouvette. L'ADN est coupé en fragments de tailles bien définies, toujours les mêmes pour une enzyme donnée.
\item \emph{Interprétation} : l'enzyme reconnaît une séquence précise de nucléotides et coupe l'ADN à cet endroit. La bactérie protège son propre ADN en modifiant chimiquement (méthylation) ses propres sites de coupure.
\item \emph{Conclusion} : il existe des enzymes, appelées \cle{enzymes de restriction}, qui coupent l'ADN double brin à des sites spécifiques. Elles constituent une arme de défense naturelle des bactéries contre les virus.
\end{itemize}

\begin{definition}[Enzyme de restriction]
Une \cle{enzyme de restriction} est une enzyme d'origine bactérienne qui coupe une molécule d'ADN double brin à des sites précis appelés \cle{sites de restriction} (ou séquences de reconnaissance), constitués en général de 4 à 8 nucléotides. On les surnomme les \og ciseaux moléculaires \fg{} du génie génétique. Chaque enzyme porte un nom issu de la bactérie dont elle provient : \emph{EcoRI} (d'\emph{Escherichia coli}), \emph{BamHI}, \emph{HindIII}, etc.
\end{definition}

\textbf{Le mécanisme de coupure.} La séquence de reconnaissance est souvent \emph{palindromique} : elle se lit de la même façon sur les deux brins, en sens opposés (5' $\to$ 3' sur un brin, 5' $\to$ 3' sur l'autre). Par exemple, l'enzyme \emph{EcoRI} reconnaît la séquence 5'-GAATTC-3'. La coupure est le plus souvent \emph{décalée} (en zigzag) : les deux brins ne sont pas coupés au même niveau. Il en résulte des fragments d'ADN portant de courtes extrémités à simple brin, complémentaires l'une de l'autre : ce sont les \cle{extrémités cohésives} (ou extrémités « collantes »). Certaines enzymes (ex. \emph{EcoRV}) coupent en face l'une de l'autre, produisant des \emph{extrémités plates} (ou « émoussées »), utilisables aussi mais moins faciles à liguer.

\begin{definition}[Fragment d'ADN et extrémités cohésives]
Un \cle{fragment d'ADN} est un morceau de molécule d'ADN obtenu après coupure par une enzyme de restriction. Lorsque la coupure est décalée, le fragment possède des \cle{extrémités cohésives} : de courtes séquences à simple brin dont les bases sont complémentaires (A face à T, G face à C). Deux fragments portant des extrémités cohésives complémentaires peuvent s'associer spontanément par appariement des bases, ce qui facilite leur soudure ultérieure.
\end{definition}

\begin{popfigure}
\begin{center}
\begin{tikzpicture}[scale=1.0, every node/.style={font=\small}]
% --- ADN avant coupure ---
\node[font=\bfseries\small, popdark] at (3,3.4) {1. ADN avant coupure};
\draw[line width=1.1pt, popblue] (0,2.6) -- (6,2.6);
\draw[line width=1.1pt, popblue] (0,1.8) -- (6,1.8);
\node[popdark] at (0.4,2.6) {\footnotesize G};
\node[popdark] at (0.8,2.6) {\footnotesize A};
\node[popdark] at (1.2,2.6) {\footnotesize A};
\node[popdark] at (1.6,2.6) {\footnotesize T};
\node[popdark] at (2.0,2.6) {\footnotesize T};
\node[popdark] at (2.4,2.6) {\footnotesize C};
\node[popdark] at (0.4,1.8) {\footnotesize C};
\node[popdark] at (0.8,1.8) {\footnotesize T};
\node[popdark] at (1.2,1.8) {\footnotesize T};
\node[popdark] at (1.6,1.8) {\footnotesize A};
\node[popdark] at (2.0,1.8) {\footnotesize A};
\node[popdark] at (2.4,1.8) {\footnotesize G};
\draw[poporange, thick, decorate, decoration={brace, amplitude=4pt, mirror}] (0.2,1.55) -- (2.6,1.55) node[midway, below=6pt, poporange, font=\footnotesize] {site de reconnaissance GAATTC};
% enzyme
\node[draw=popgreen, fill=popgreenL, rounded corners, font=\footnotesize\bfseries, align=center] (enz) at (4.6,3.1) {enzyme de\\restriction};
\draw[-{Stealth}, popgreen, thick] (enz.south) -- (1.4,2.75);
% --- Après coupure ---
\node[font=\bfseries\small, popdark] at (3,-0.6) {2. Après coupure : deux fragments à extrémités cohésives};
% fragment gauche
\draw[line width=1.1pt, popblue] (0,0.6) -- (1.0,0.6);
\draw[line width=1.1pt, popblue] (0,-0.2) -- (2.0,-0.2);
\node[popdark] at (0.4,0.6) {\footnotesize G};
\node[popdark] at (0.4,-0.2) {\footnotesize C};
\node[popdark] at (0.8,-0.2) {\footnotesize T};
\node[popdark] at (1.2,-0.2) {\footnotesize T};
\node[popdark] at (1.6,-0.2) {\footnotesize A};
\node[popdark] at (2.0,-0.2) {\footnotesize A};
\draw[popgold, thick, fill=popgoldL, opacity=0.6] (0.15,0.85) rectangle (1.05,0.35);
% fragment droit
\draw[line width=1.1pt, popblue] (1.4,0.6) -- (6,0.6);
\draw[line width=1.1pt, popblue] (2.4,-0.2) -- (6,-0.2);
\node[popdark] at (1.8,0.6) {\footnotesize A};
\node[popdark] at (2.2,0.6) {\footnotesize A};
\node[popdark] at (2.6,0.6) {\footnotesize T};
\node[popdark] at (3.0,0.6) {\footnotesize T};
\node[popdark] at (3.4,0.6) {\footnotesize C};
\node[popdark] at (2.4,-0.2) {\footnotesize G};
\draw[popgold, thick, fill=popgoldL, opacity=0.6] (1.55,0.85) rectangle (3.45,0.35);
% légendes pointées
\draw[thin, popmuted] (2.0,0.95) -- (4.2,1.5) node[right, font=\footnotesize, align=left, popdark] {extrémités cohésives\\(simples brins complémentaires AATT)};
\draw[thin, popmuted] (1.2,0.2) -- (1.2,-1.3) node[below, font=\footnotesize, popdark, align=center] {point de coupure\\(décalé, en zigzag)};
\end{tikzpicture}
\end{center}
\legende{Action d'une enzyme de restriction (type \emph{EcoRI}) sur une molécule d'ADN double brin. L'enzyme reconnaît la séquence GAATTC et coupe les deux brins de façon décalée (entre G et A), ce qui génère des fragments portant des extrémités cohésives complémentaires (AATT / TTAA).}
\end{popfigure}

\begin{aretenir}[L'essentiel sur les enzymes de restriction]
\begin{itemize}
\item Elles sont d'origine bactérienne et coupent l'ADN à des sites spécifiques (séquences palindromiques le plus souvent).
\item Une même enzyme coupe toujours à la même séquence : la coupure est donc précise et reproductible.
\item La coupure décalée produit des extrémités cohésives complémentaires (ex. AATT), qui permettront le recollement de fragments d'origines différentes. La coupure droite produit des extrémités plates.
\item \textbf{Clé du génie génétique} : si l'on coupe l'ADN humain et un plasmide avec la \emph{même} enzyme, leurs extrémités cohésives sont compatibles.
\end{itemize}
\end{aretenir}

\begin{savaistu}[Des ciseaux inventés par les bactéries]
Les enzymes de restriction n'ont pas été inventées par l'homme : elles sont naturellement produites par les bactéries pour se défendre contre les virus (bactériophages) qui les infectent. En découpant l'ADN viral injecté, la bactérie « restreint » l'infection — d'où le nom d'enzyme de restriction. Son propre ADN est protégé par une modification chimique (méthylation) de ses sites de coupure. La découverte de ces enzymes a valu le prix Nobel de physiologie ou médecine en 1978 à Werner Arber, Daniel Nathans et Hamilton Smith. Aujourd'hui, des centaines d'enzymes de restriction différentes sont commercialisées et utilisées dans les laboratoires du monde entier, y compris au Cameroun.
\end{savaistu}

\courssub{Les ligases : la colle moléculaire}

\textbf{Situation concrète.} Revenons à notre chercheur : il a découpé le gène de l'insuline et ouvert un plasmide avec la même enzyme de restriction. Les extrémités cohésives du gène et du plasmide s'apparient spontanément par liaisons hydrogène, mais cette association est fragile : les deux molécules ne sont pas réellement soudées, la liaison phosphodiester entre les nucléotides n'est pas refermée. Il manque un outil : la soudeuse.

\begin{definition}[Ligase]
Une \cle{ligase} est une enzyme qui soude des fragments d'ADN entre eux en reformant les liaisons chimiques (liaisons phosphodiester) entre les nucléotides des brins, consommant de l'ATP (ou NAD$^+$). On la surnomme la \og colle moléculaire \fg{}. La ligase la plus utilisée en laboratoire est la ligase T4, issue d'un virus bactérien (bactériophage T4). Dans la cellule vivante, les ligases interviennent naturellement lors de la réplication (jonction des fragments d'Okazaki) et de la réparation de l'ADN.
\end{definition}

\textbf{Mécanisme.} Après appariement des extrémités cohésives complémentaires (ou mise en contact d'extrémités plates), la ligase catalyse la formation de la liaison phosphodiester manquante sur chacun des deux brins. La molécule d'ADN obtenue est alors continue, stable et covalente : le gène étranger est définitivement inséré dans le vecteur.

\begin{attention}[Ne pas confondre ligase et enzyme de restriction]
C'est l'erreur la plus fréquente au Probatoire :
\begin{itemize}
\item l'\cle{enzyme de restriction} \textbf{coupe} l'ADN (ciseaux moléculaires) ;
\item la \cle{ligase} \textbf{soude} les fragments d'ADN (colle moléculaire).
\end{itemize}
La ligase ne coupe jamais ; l'enzyme de restriction ne soude jamais. Dans un schéma de manipulation, repère le sens de l'action : si deux molécules séparées deviennent une seule molécule, c'est une ligase qui est intervenue ; si une molécule devient plusieurs fragments, c'est une enzyme de restriction.
\end{attention}

\courssub{Les vecteurs de transfert : les transporteurs de gènes}

\textbf{Le problème à résoudre.} Même soudé correctement, un gène isolé ne peut pas pénétrer seul dans une cellule hôte, ni s'y multiplier, ni s'y exprimer. Il faut un « véhicule » capable de franchir la membrane cellulaire, de se répliquer dans la cellule (ou de s'intégrer au génome) et d'entraîner le gène étranger avec lui. Ce véhicule est le vecteur de transfert.

\begin{definition}[Vecteur de transfert]
Un \cle{vecteur de transfert} est une molécule d'ADN capable d'introduire un gène étranger dans une cellule hôte et de s'y répliquer (ou de s'y intégrer). Le vecteur le plus utilisé est le \cle{plasmide} bactérien ; on utilise aussi des virus modifiés (désarmés) pour transférer des gènes dans des cellules végétales ou animales.
\end{definition}

\textbf{Le plasmide bactérien, vecteur de choix.} Le \cle{plasmide} est une petite molécule d'ADN circulaire, double brin, présente dans le cytoplasme des bactéries, en plus du chromosome bactérien. Il possède trois propriétés essentielles qui en font un excellent vecteur :
\begin{enumerate}
\item il est \textbf{petit} (quelques kpb) et facile à extraire, purifier et manipuler en éprouvette ;
\item il possède une \textbf{origine de réplication} (\emph{ori}) : il se réplique de façon autonome dans la bactérie, indépendamment du chromosome, si bien que le gène inséré sera copié en grand nombre à chaque division bactérienne (amplification) ;
\item il porte un \textbf{site de clonage multiple} (SCM ou \emph{polylinker}), région courte contenant de nombreux sites de restriction uniques pour y insérer le gène d'intérêt, et des \textbf{gènes marqueurs} (souvent un gène de résistance à un antibiotique comme l'ampicilline) qui permettent de \emph{sélectionner} les bactéries ayant bien reçu le plasmide recombinant.
\end{enumerate}

\textbf{Autres vecteurs (complément utile au Probatoire).}
\begin{itemize}
\item \textbf{Pour les cellules végétales} : on utilise fréquemment le plasmide Ti (Tumor-inducing) de la bactérie du sol \emph{Agrobacterium tumefaciens}. Cette bactérie est un « ingénieur génétique naturel » : elle transfère spontanément un fragment de son plasmide Ti (l'ADN-T) dans le génome nucléaire de la plante hôte. On remplace les gènes tumoraux de l'ADN-T par le gène d'intérêt.
\item \textbf{Pour les cellules animales ou humaines (thérapie génique)} : on utilise des \textbf{vecteurs viraux} (rétrovirus, adénovirus, AAV). On retire leurs gènes pathogènes (on les « désarme ») et on les remplace par le gène d'intérêt ; le virus conserve sa capacité naturelle à pénétrer dans la cellule cible, mais ne provoque plus de maladie.
\end{itemize}

\courssub{L'ADN recombinant : le produit de la soudure}

Lorsqu'un gène étranger (par exemple le gène humain de l'insuline) est soudé par une ligase à un vecteur (par exemple un plasmide bactérien ouvert), on obtient une molécule d'ADN hybride, composée d'ADN d'origines différentes : c'est l'\cle{ADN recombinant} (ou ADN chimérique).

\begin{definition}[ADN recombinant]
Un \cle{ADN recombinant} est une molécule d'ADN artificielle obtenue par la soudure, grâce à une ligase, d'un fragment d'ADN étranger (le gène d'intérêt) à une molécule d'ADN vecteur (plasmide ou virus modifié). Introduit dans une cellule hôte, l'ADN recombinant se réplique (si le vecteur est un plasmide) ou s'intègre (si vecteur viral) et permet l'expression du gène étranger (transcription $\to$ traduction $\to$ protéine).
\end{definition}

\begin{popfigure}
\begin{center}
\begin{tikzpicture}[scale=1.0, every node/.style={font=\small}]
% Plasmide ouvert
\node[font=\bfseries\small, popdark] at (2.2,3.2) {1. Plasmide ouvert par une enzyme de restriction};
\draw[line width=1.2pt, popblue] (0.6,2.2) arc[start angle=180, end angle=0, x radius=1.6, y radius=1.3];
\draw[line width=1.2pt, popblue] (0.6,1.6) arc[start angle=180, end angle=360, x radius=1.6, y radius=1.3];
\draw[line width=1.2pt, popblue] (0.6,2.2) -- (0.2,2.2);
\draw[line width=1.2pt, popblue] (3.8,1.6) -- (4.2,1.6);
\node[popdark, font=\footnotesize] at (0.1,2.2) {AATT};
\node[popdark, font=\footnotesize] at (4.55,1.6) {TTAA};
\draw[thin, popmuted] (2.2,0.1) -- (2.2,-0.5) node[below, font=\footnotesize, popdark] {ADN plasmidique linéarisé (bactérie)};
% Gène étranger
\node[font=\bfseries\small, popdark] at (8.6,3.2) {2. Gène étranger découpé (même enzyme)};
\draw[line width=1.2pt, poporange] (6.6,2.2) -- (10.6,2.2);
\draw[line width=1.2pt, poporange] (6.6,1.6) -- (10.6,1.6);
\node[popdark, font=\footnotesize] at (6.35,2.2) {TTAA};
\node[popdark, font=\footnotesize] at (10.9,1.6) {AATT};
\draw[thin, popmuted] (8.6,1.4) -- (8.6,0.6) node[below, font=\footnotesize, popdark, align=center] {gène d'intérêt\\(ex. : insuline humaine)};
% flèche + ligase
\draw[-{Stealth}, very thick, popgreen] (5.2,-1.3) -- (5.2,-2.6) node[midway, right, font=\footnotesize\bfseries, popgreen, align=left] {ligase T4 + ATP\\(soudure des\\extrémités cohésives)};
% Plasmide recombinant
\node[font=\bfseries\small, popdark] at (5.2,-3.3) {3. Plasmide recombinant (ADN recombinant)};
\draw[line width=1.2pt, popblue] (3.6,-4.6) arc[start angle=180, end angle=90, radius=1.6];
\draw[line width=1.2pt, popblue] (6.8,-4.6) arc[start angle=0, end angle=90, radius=1.6];
\draw[line width=1.2pt, poporange] (3.6,-4.6) arc[start angle=180, end angle=270, radius=1.6];
\draw[line width=1.2pt, poporange] (6.8,-4.6) arc[start angle=360, end angle=270, radius=1.6];
\draw[thin, popmuted] (5.2,-6.4) -- (7.6,-6.9) node[right, font=\footnotesize, popdark, align=left] {gène étranger inséré\\(en orange)};
\draw[thin, popmuted] (5.2,-3.0) -- (7.6,-2.6) node[right, font=\footnotesize, popdark, align=left] {ADN du plasmide\\(en bleu)};
\end{tikzpicture}
\end{center}
\legende{Formation d'un plasmide recombinant. Le plasmide bactérien (bleu) et le gène étranger (orange) sont coupés par la même enzyme de restriction (\emph{EcoRI}) : leurs extrémités cohésives sont complémentaires (AATT / TTAA). La ligase soude le gène dans le plasmide : on obtient un ADN recombinant circulaire, qui sera introduit dans une bactérie hôte par transformation.}
\end{popfigure}

\begin{methode}[Identifier les outils dans un schéma de manipulation]
Face à un schéma de génie génétique au Probatoire, pose-toi trois questions dans l'ordre :
\begin{enumerate}
\item \textbf{Qu'est-ce qui coupe ?} Repère l'étape où une molécule d'ADN est divisée en fragments : l'outil est une \cle{enzyme de restriction} (ciseaux moléculaires). Précise si possible le nom (ex. \emph{EcoRI}) et le type d'extrémités (cohésives ou plates).
\item \textbf{Qu'est-ce qui soude ?} Repère l'étape où deux fragments d'ADN ne font plus qu'un (formation d'une molécule continue) : l'outil est une \cle{ligase} (colle moléculaire). Cite la ligase T4 et le cofacteur (ATP).
\item \textbf{Qu'est-ce qui transporte ?} Repère la molécule qui porte le gène étranger et pénètre dans la cellule hôte : c'est le \cle{vecteur de transfert} (le plus souvent un plasmide bactérien, parfois un virus modifié ou le plasmide Ti d'\emph{Agrobacterium}).
\end{enumerate}
Vérifie enfin que le produit final (gène + vecteur soudés) est bien qualifié d'\cle{ADN recombinant} et que le vecteur possède un gène marqueur pour la sélection.
\end{methode}

\begin{exoresolu}[Identifier les outils sur un document]
\emph{Énoncé.} Sur un schéma de production d'insuline, on lit les étapes suivantes : (a) l'ADN humain est découpé en nombreux fragments ; (b) un plasmide de bactérie est ouvert ; (c) le fragment portant le gène de l'insuline est inséré dans le plasmide et la molécule est refermée ; (d) le plasmide est introduit dans une bactérie qui produit ensuite de l'insuline. Pour chaque étape, indique l'outil du génie génétique mis en jeu et justifie.
\tcblower
\emph{Solution détaillée.}
\begin{itemize}
\item \textbf{Étape (a)} : l'ADN humain est découpé en fragments $\to$ il s'agit d'une \cle{enzyme de restriction}, car elle seule coupe l'ADN à des sites spécifiques (rôle de ciseaux moléculaires).
\item \textbf{Étape (b)} : le plasmide est ouvert $\to$ la \emph{même} enzyme de restriction, afin que le plasmide porte des extrémités cohésives complémentaires de celles du gène.
\item \textbf{Étape (c)} : le gène est inséré et la molécule refermée $\to$ c'est la \cle{ligase} (T4) qui soude le gène au plasmide (colle moléculaire). La molécule obtenue est un \cle{ADN recombinant}.
\item \textbf{Étape (d)} : le plasmide transporte le gène dans la bactérie $\to$ le plasmide joue le rôle de \cle{vecteur de transfert}. Sa réplication autonome (origine \emph{ori}) permet d'amplifier le gène et son gène marqueur (résistance antibiotique) permet de sélectionner les bactéries transformées. L'expression du gène (transcription/traduction) produit l'insuline.
\end{itemize}
\emph{Conclusion} : toute manipulation de génie génétique mobilise la même trilogie d'outils : enzyme de restriction (couper), ligase (souder), vecteur (transporter/sélectionner/amplifier).
\end{exoresolu}

\begin{exercice}[À toi de jouer]
On veut créer une bactérie capable de produire une enzyme qui dégrade les hydrocarbures, afin de dépolluer des sols souillés par des fuites de pétrole (problème fréquent dans la région de Limbé/Kribi). 
\begin{enumerate}
\item Nomme les trois outils indispensables à cette manipulation et précise le rôle de chacun.
\item Explique pourquoi il faut utiliser la \emph{même} enzyme de restriction pour découper le gène d'intérêt et ouvrir le plasmide.
\item Comment appelle-t-on la molécule d'ADN obtenue à la fin de la soudure ? Quelle propriété du plasmide permettra d'amplifier cette molécule ?
\item Comment pourra-t-on isoler les bactéries qui ont bien reçu le plasmide recombinant parmi toutes les bactéries du milieu de culture ?
\end{enumerate}
\end{exercice}

\begin{corrige}[Exercice]
\begin{enumerate}
\item Les trois outils sont : 
      \begin{itemize}
      \item l'\cle{enzyme de restriction} : elle coupe l'ADN portant le gène d'intérêt (ex. gène de dégradation d'hydrocarbures) et ouvre le plasmide au niveau de son site de clonage multiple ;
      \item la \cle{ligase} (T4) : elle soude le gène au plasmide en reformant les liaisons phosphodiester (colle moléculaire) ;
      \item le \cle{vecteur de transfert} (plasmide) : il introduit le gène dans la bactérie hôte, s'y réplique (amplification) et porte un gène marqueur pour la sélection.
      \end{itemize}
\item Il faut utiliser la même enzyme de restriction afin que le gène et le plasmide portent des extrémités cohésives \textbf{complémentaires} (ex. AATT / TTAA) : seules ces extrémités peuvent s'apparier par liaisons hydrogène puis être soudées covalemment par la ligase.
\item La molécule obtenue est un \cle{ADN recombinant} (ou ADN chimérique), car elle associe de l'ADN d'origines différentes (gène étranger + plasmide bactérien). La propriété du plasmide qui permet son amplification est son \textbf{origine de réplication (\emph{ori})}, qui assure sa réplication autonome lors de la division bactérienne.
\item On utilise le \textbf{gène marqueur} porté par le plasmide (ex. résistance à l'ampicilline). On ensemence les bactéries sur un milieu contenant l'antibiotique : seules celles qui ont intégré le plasmide recombinant survivent et forment des colonies.
\end{enumerate}
\end{corrige}

\begin{aretenir}[Bilan de la section]
Le génie génétique repose sur trois outils moléculaires complémentaires : les \cle{enzymes de restriction}, ciseaux bactériens qui coupent l'ADN à des sites spécifiques (séquences palindromiques) en générant le plus souvent des extrémités cohésives ; les \cle{ligases}, colle moléculaire qui soude les fragments d'ADN en reformant les liaisons phosphodiester (avec ATP) ; et les \cle{vecteurs de transfert}, principalement les plasmides bactériens (ADN circulaires, origine de réplication \emph{ori}, site de clonage multiple, gène marqueur de sélection), parfois des virus modifiés ou le plasmide Ti d'\emph{Agrobacterium} pour les plantes. La soudure d'un gène à un vecteur produit un \cle{ADN recombinant}. La section suivante montrera comment ces outils s'enchaînent, étape par étape, pour obtenir un organisme génétiquement modifié (OGM) complet.
\end{aretenir}

\coursec{Étapes de l'obtention d'un OGM}

Dans la section précédente, nous avons découvert la « boîte à outils » du génie génétique : les \cle{enzymes de restriction} qui coupent l'ADN à des sites précis, les \cle{ligases} qui soudent les fragments, et les \cle{vecteurs} (plasmides) qui transportent le gène d'intérêt jusqu'à la cellule hôte. Mais disposer d'outils ne suffit pas : encore faut-il savoir dans quel ordre les utiliser. C'est précisément l'objet de cette section : décrire, étape par étape, le protocole complet qui conduit d'un gène d'intérêt à un \cle{organisme génétiquement modifié} (OGM) capable d'exprimer ce gène.

Pour rendre la démarche concrète, nous suivrons tout au long de la section l'exemple historique et emblématique de la \cle{production d'insuline humaine par des bactéries transgéniques}, réalisée pour la première fois en 1978--1979 par les équipes de Genentech (États-Unis). Cet exemple est un classique du Probatoire : il faut savoir le décrire dans l'ordre, sans confusion.

\courssub{Situation-problème : comment soigner des millions de diabétiques ?}

Le diabète de type 1 est une maladie dans laquelle les cellules $\beta$ des îlots de Langerhans du pancréas ne produisent plus d'\cle{insuline}, l'hormone qui régule la glycémie. Le traitement consiste à injecter quotidiennement de l'insuline au malade. Au Cameroun comme partout dans le monde, le nombre de diabétiques ne cesse d'augmenter.

Avant 1982, l'insuline était extraite du pancréas de porcs ou de bovins abattus à l'abattoir. Cette méthode posait de graves problèmes :
\begin{itemize}
\item il fallait d'énormes quantités de pancréas animaux pour obtenir peu d'insuline (environ \SI{8000}{kg} de pancréas pour \SI{1}{kg} d'insuline purifiée) ;
\item l'insuline animale, légèrement différente de l'insuline humaine, provoquait parfois des \cle{réactions immunitaires} (allergies, résistances) chez les malades ;
\item l'approvisionnement dépendait entièrement de l'abattage des animaux.
\end{itemize}

\begin{center}
\textbf{Problème scientifique :} comment obtenir de l'insuline \emph{strictement humaine}, en grande quantité, à moindre coût ?
\end{center}

\textbf{Hypothèse :} puisque toute protéine est codée par un gène, si l'on parvenait à introduire le \emph{gène humain de l'insuline} dans une bactérie qui se multiplie très vite (comme \emph{Escherichia coli}, une division toutes les 20 minutes environ), la bactérie deviendrait une véritable « usine » à insuline.

La vérification de cette hypothèse a donné naissance à la première biotechnologie moderne. Décortiquons-en les six étapes.

\courssub{Étape 1 : Identification et isolement du gène d'intérêt}

\begin{definition}[Gène d'intérêt]
Le \cle{gène d'intérêt} (ou \cle{transgène}) est le fragment d'ADN que l'on souhaite transférer dans un autre organisme parce qu'il code une protéine utile (ici l'insuline humaine) ou confère un caractère recherché (résistance à un insecte, tolérance à la sécheresse...).
\end{definition}

La première difficulté est de \emph{retrouver} ce gène parmi les quelque 20\,000 gènes du génome humain. Deux stratégies existent :
\begin{enumerate}
\item \textbf{Extraire l'ADN} de cellules humaines et découper le génome avec des enzymes de restriction, puis identifier le fragment portant le gène de l'insuline (grâce à une \emph{sonde moléculaire} complémentaire du gène recherché) ;
\item \textbf{Fabriquer le gène artificiellement} : à partir de l'ARN messager de l'insuline (abondant dans les cellules $\beta$ du pancréas), on synthétise \emph{in vitro} un ADN complémentaire (ADNc) grâce à une enzyme appelée \emph{transcriptase inverse}. C'est la stratégie historiquement utilisée pour l'insuline, car l'ADNc ne contient pas les introns que les bactéries ne savent pas éliminer.
\end{enumerate}

\begin{attention}[Piège classique]
Une bactérie ne possède pas de \emph{spliceosome} (épissosome) : elle est \textbf{incapable d'enlever les introns} d'un gène humain. Si l'on insère le gène humain complet (avec ses introns), la bactérie produira une protéine fausse. C'est pourquoi on utilise l'ADNc (copie de l'ARN messager mature, donc sans introns) ou un gène synthétisé chimiquement. Cette précision est souvent attendue dans les questions de synthèse.
\end{attention}

\courssub{Étape 2 : Découpage du gène et du vecteur par la même enzyme de restriction}

On dispose maintenant de deux molécules d'ADN : le fragment contenant le gène de l'insuline et un \cle{plasmide} (petite molécule d'ADN circulaire bactérien) qui servira de vecteur. On les traite \textbf{tous les deux par la même enzyme de restriction}, par exemple \emph{EcoRI}.

L'enzyme coupe l'ADN au niveau de sa séquence de reconnaissance spécifique, en générant des \cle{extrémités cohésives} (ou « collantes ») : de courts fragments d'ADN simple brin complémentaires les uns des autres.

\begin{propriete}[Extrémités complémentaires (admise)]
Lorsque deux molécules d'ADN sont coupées par la \textbf{même enzyme de restriction}, elles portent des extrémités cohésives \textbf{complémentaires}. Ces extrémités s'apparient spontanément par complémentarité des bases (A avec T, G avec C), ce qui permet ensuite la soudure définitive par l'ADN ligase. Cette propriété est admise ; aucune démonstration n'est exigible au Probatoire, mais tu dois savoir l'exploiter dans la description d'un schéma.
\end{propriete}

\courssub{Étape 3 : Insertion du gène dans le plasmide et soudure par la ligase}

On mélange le gène d'intérêt et les plasmides ouverts. Par simple agitation thermique, les extrémités cohésives complémentaires se reconnaissent et s'apparient. On ajoute alors l'\cle{ADN ligase}, qui rétablit les liaisons phosphodiester du squelette sucre-phosphate : le gène est définitivement « cousu » dans le plasmide.

\begin{definition}[ADN recombinant]
On obtient un \cle{ADN recombinant} : une molécule d'ADN chimérique constituée de fragments d'ADN d'origines différentes (ici, ADN plasmidique bactérien + gène humain de l'insuline).
\end{definition}

\begin{attention}[Le mélange n'est pas homogène]
Dans le tube, toutes les molécules ne donnent pas le résultat voulu : certains plasmides se referment sur eux-mêmes sans avoir capturé le gène, d'autres intègrent le gène dans le mauvais sens. Le génie génétique est une technique \emph{statistique} : on obtient un mélange dont il faudra \textbf{trier} les bons éléments (voir étape 5).
\end{attention}

\courssub{Étape 4 : Introduction du plasmide recombinant dans la cellule hôte (transformation)}

Il faut maintenant faire entrer le plasmide recombinant dans la bactérie hôte, \emph{E. coli}. Cette étape s'appelle la \cle{transformation}.

\begin{definition}[Transformation bactérienne]
La \cle{transformation} est l'introduction d'un ADN étranger (ici le plasmide recombinant) dans une cellule hôte, qui acquiert ainsi de nouvelles caractéristiques héréditaires.
\end{definition}

En pratique, on rend la membrane bactérienne temporairement perméable (traitement par le chlorure de calcium \ce{CaCl2} suivi d'un choc thermique, ou électroporation : brève impulsion électrique). Une petite fraction des bactéries internalise alors un plasmide. Ce point est crucial : \textbf{la grande majorité des bactéries n'intègre aucun plasmide}.

\courssub{Étape 5 : Sélection des cellules transformées}

C'est l'étape que les élèves oublient le plus souvent, alors qu'elle est \emph{indispensable}. Comment distinguer les rares bactéries transformées parmi des millions de bactéries non transformées ? On utilise un \cle{marqueur de sélection}.

Le plasmide vecteur a été choisi (ou construit) de façon à porter, en plus du site d'insertion du gène d'intérêt, un \textbf{gène de résistance à un antibiotique}, par exemple l'ampicilline. On étale alors la culture sur un milieu gélosé contenant cet antibiotique :
\begin{itemize}
\item les bactéries \textbf{non transformées} (sans plasmide) sont sensibles à l'antibiotique : elles \emph{meurent} ;
\item les bactéries \textbf{transformées} (portant le plasmide recombinant) résistent à l'antibiotique : elles \emph{survivent} et forment des colonies.
\end{itemize}

\begin{attention}[Erreur fréquente au Probatoire]
Ne jamais écrire que « toutes les bactéries reçoivent le gène ». C'est faux : la transformation est un événement rare. Sans l'étape de \textbf{sélection} sur milieu avec antibiotique, on ne pourrait jamais isoler les bactéries productrices d'insuline. Chaque fois qu'un schéma de transgénèse montre un milieu de culture avec un antibiotique, tu dois l'identifier comme l'étape de sélection et l'expliquer.
\end{attention}

\begin{popfigure}
\begin{center}
\begin{tikzpicture}[scale=0.9]
% Boîte de Pétri gauche : avant sélection
\draw[thick, popblue] (0,0) circle (1.6);
\draw[thick, popblue] (0,0) circle (1.75);
\foreach \x/\y in {-0.9/0.6, -0.3/-0.8, 0.7/0.4, 1.0/-0.5, -1.1/-0.4, 0.2/1.0, 0.5/-1.1, -0.5/0.1, 1.2/0.9, -1.3/0.9, 0.9/1.2, -0.2/-1.3}{
  \fill[popmuted] (\x,\y) circle (0.06);}
\foreach \x/\y in {0.1/0.5, -0.7/-0.9}{
  \fill[poporange] (\x,\y) circle (0.09);}
\node[below, align=center] at (0,-1.9) {\small Milieu \textbf{avant} sélection\\ \small (toutes les bactéries)};
\draw[popmuted, thin] (-0.5,0.1) -- (-2.6,1.3) node[left, align=right, text width=3.2cm] {\scriptsize Bactéries non transformées (sensibles)};
\draw[poporange, thin] (0.1,0.5) -- (2.6,1.5) node[right, align=left, text width=3cm] {\scriptsize Bactéries transformées (résistantes)};
% Flèche
\draw[-{Stealth}, very thick, popdark] (2.1,0) -- (3.4,0) node[midway, above, align=center, text width=2.6cm] {\scriptsize Ajout d'ampicilline};
% Boîte de Pétri droite : après sélection
\draw[thick, popgreen] (5.5,0) circle (1.6);
\draw[thick, popgreen] (5.5,0) circle (1.75);
\foreach \x/\y in {5.6/0.5, 4.8/-0.9}{
  \fill[poporange] (\x,\y) circle (0.09);}
\node[below, align=center] at (5.5,-1.9) {\small Milieu \textbf{après} sélection\\ \small (seules les transformées survivent)};
\end{tikzpicture}
\end{center}
\legende{Sélection des bactéries transformées sur milieu contenant un antibiotique (ampicilline). Seules les bactéries ayant intégré le plasmide recombinant, porteur du gène de résistance, forment des colonies.}
\end{popfigure}

\courssub{Étape 6 : Multiplication (clonage) et expression du gène}

Chaque bactérie transformée survivante est une cellule génétiquement modifiée. Placée dans un \cle{fermenteur} (grand récipient de culture contrôlé en température, pH et oxygène), elle se divise par mitose (scissiparité) environ toutes les 20 minutes : on obtient en quelques heures des milliards de bactéries \textbf{génétiquement identiques}, toutes porteuses du gène de l'insuline. C'est le \cle{clonage moléculaire}.

Chaque bactérie \emph{exprime} le transgène : elle transcrit le gène de l'insuline en ARN messager, puis le traduit en protéine. L'insuline humaine s'accumule dans la culture ; il ne reste plus qu'à l'\textbf{extraire et la purifier}.

\begin{exemplebox}[Une « usine » microscopique : ordre de grandeur]
À partir d'\textbf{une seule} bactérie transformée, une culture de quelques litres en fermenteur produit en \textbf{quelques jours} des quantités d'insuline suffisantes pour préparer des \textbf{millions de doses} injectables. Depuis 1982 (insuline « Humuline »), toute l'insuline utilisée dans le monde, y compris dans les hôpitaux camerounais, est produite par cette biotechnologie : elle est humaine à 100\,\%, abondante et bien moins chère que l'insuline animale.
\end{exemplebox}

\begin{popfigure}
\begin{center}
\begin{tikzpicture}[scale=0.86, every node/.style={transform shape},
  etape/.style={draw, thick, rounded corners=4pt, align=center, text width=3.5cm, minimum height=1.5cm, font=\small},
  num/.style={circle, draw=popdark, fill=popdark, text=white, font=\bfseries\small, inner sep=2pt}]
% Étape 1
\node[etape, fill=popblueL, draw=popblue] (e1) at (0,0) {Gène humain de l'insuline isolé (ADNc)};
\node[num] at ([shift={(-0.25,0.55)}]e1.north west) {1};
% Étape 2
\node[etape, fill=popblueL, draw=popblue] (e2) at (4.6,0) {Découpage du gène et du plasmide par la \textbf{même} enzyme de restriction};
\node[num] at ([shift={(-0.25,0.55)}]e2.north west) {2};
% Étape 3
\node[etape, fill=poppurpleL, draw=poppurple] (e3) at (9.2,0) {Insertion + soudure par la ligase : \textbf{ADN recombinant}};
\node[num] at ([shift={(-0.25,0.55)}]e3.north west) {3};
% Étape 4
\node[etape, fill=poporangeL, draw=poporange] (e4) at (9.2,-2.9) {Transformation : introduction du plasmide dans \emph{E. coli}};
\node[num] at ([shift={(-0.25,0.55)}]e4.north west) {4};
% Étape 5
\node[etape, fill=popgoldL, draw=popgold] (e5) at (4.6,-2.9) {Sélection sur milieu + antibiotique (marqueur de résistance)};
\node[num] at ([shift={(-0.25,0.55)}]e5.north west) {5};
% Étape 6
\node[etape, fill=popgreenL, draw=popgreen] (e6) at (0,-2.9) {Clonage en fermenteur + expression : production d'\textbf{insuline humaine}};
\node[num] at ([shift={(-0.25,0.55)}]e6.north west) {6};
% Flèches
\draw[-{Stealth}, very thick, popdark] (e1) -- (e2);
\draw[-{Stealth}, very thick, popdark] (e2) -- (e3);
\draw[-{Stealth}, very thick, popdark] (e3) -- (e4);
\draw[-{Stealth}, very thick, popdark] (e4) -- (e5);
\draw[-{Stealth}, very thick, popdark] (e5) -- (e6);
\end{tikzpicture}
\end{center}
\legende{Les six étapes de l'obtention d'une bactérie transgénique productrice d'insuline humaine. L'ordre des étapes est exigible : isolement du gène $\rightarrow$ découpage $\rightarrow$ insertion dans le vecteur $\rightarrow$ transformation $\rightarrow$ sélection $\rightarrow$ multiplication et expression.}
\end{popfigure}

\courssub{Cas des plantes transgéniques : de la cellule transformée à la plante entière}

Chez une bactérie, il suffit de laisser la cellule transformée se diviser. Mais chez une plante (maïs, coton, riz...), on part généralement d'\textbf{une seule cellule végétale transformée} (par exemple après infection par la bactérie \emph{Agrobacterium tumefaciens}, vecteur naturel de gènes). Comment obtenir une \emph{plante entière} ?

\begin{propriete}[Totipotence cellulaire (admise)]
Chez les végétaux, toute cellule somatique (non sexuelle) possède le \textbf{génome complet} de l'individu et conserve la capacité de régénérer une plante entière : c'est la \cle{totipotence}. Placée sur un milieu de culture adapté (hormones de croissance : auxines et cytokinines), une cellule végétale transformée se divise, forme un \emph{cal} (amas de cellules indifférenciées), puis se différencie en bourgeons, racines, et finalement en une plante entière dont \textbf{toutes les cellules portent le transgène}.
\end{propriete}

C'est ainsi que l'on obtient, par exemple, le maïs Bt (porteur d'un gène bactérien de résistance aux insectes foreurs) ou le coton Bt cultivé dans certains pays africains. Chaque graine récoltée transmet le transgène à la génération suivante : le caractère est \emph{héréditaire}.

\courssub{Savoir-faire : décrire un schéma de transgénèse}

\begin{methode}[Décrire une manipulation de génie génétique à partir d'un schéma]
Face à un schéma de transgénèse au Probatoire, procède toujours dans cet ordre :
\begin{enumerate}
\item \textbf{Identifier} le gène d'intérêt et l'organisme donneur (ex. gène de l'insuline, humain) ;
\item \textbf{Repérer} l'enzyme de restriction et vérifier qu'elle coupe le gène \emph{et} le vecteur (même enzyme $\rightarrow$ extrémités complémentaires) ;
\item \textbf{Nommer} la molécule obtenue après action de la ligase : ADN recombinant ;
\item \textbf{Identifier} la cellule hôte et nommer l'étape : transformation ;
\item \textbf{Repérer} le marqueur de sélection (résistance à un antibiotique) et expliquer son rôle : éliminer les cellules non transformées ;
\item \textbf{Conclure} par la multiplication (clonage) et l'expression du gène : production de la protéine d'intérêt.
\end{enumerate}
Rédige avec le vocabulaire exact : transgène, vecteur, ADN recombinant, transformation, sélection, clonage, expression.
\end{methode}

\begin{exoresolu}[Analyse d'un schéma de transgénèse]
\textbf{Énoncé.} Un document montre les opérations suivantes : (a) extraction d'ADN de cellules pancréatiques humaines ; (b) action de l'enzyme \emph{EcoRI} sur cet ADN et sur des plasmides ; (c) mélange des fragments avec de l'ADN ligase ; (d) contact des plasmides avec des bactéries \emph{E. coli} ; (e) étalement sur milieu gélosé contenant de l'ampicilline ; (f) culture en fermenteur et récolte d'une protéine injectée ensuite à des diabétiques. Décris cette manipulation en identifiant chaque étape et en précisant le rôle de l'ampicilline.
\tcblower
\textbf{Solution détaillée.}
\begin{itemize}
\item \textbf{(a) et (b)} : identification et isolement du gène d'intérêt (gène de l'insuline humaine), puis découpage du gène et du vecteur plasmidique par la \emph{même} enzyme de restriction \emph{EcoRI}, ce qui génère des extrémités cohésives complémentaires.
\item \textbf{(c)} : insertion du gène dans le plasmide et soudure par l'ADN ligase : on obtient un \textbf{ADN recombinant}.
\item \textbf{(d)} : \textbf{transformation} : introduction du plasmide recombinant dans la bactérie hôte \emph{E. coli}. Seule une faible proportion des bactéries intègre le plasmide.
\item \textbf{(e)} : \textbf{sélection}. L'ampicilline tue toutes les bactéries non transformées ; seules survivent les bactéries portant le plasmide, qui contient un gène de résistance à l'ampicilline (marqueur de sélection). L'ampicilline sert donc à \emph{éliminer les bactéries n'ayant pas intégré le transgène}.
\item \textbf{(f)} : \textbf{clonage et expression} : les bactéries transformées se multiplient en fermenteur et expriment le gène humain ; la protéine récoltée est l'\textbf{insuline humaine}, utilisée pour traiter les diabétiques.
\end{itemize}
\textbf{Conclusion :} cette manipulation illustre les six étapes de l'obtention d'un OGM ; l'organisme obtenu est une bactérie transgénique productrice d'insuline humaine.
\end{exoresolu}

\begin{aretenir}[L'essentiel de la section]
\begin{itemize}
\item L'obtention d'un OGM suit \textbf{six étapes ordonnées} : isolement du gène d'intérêt $\rightarrow$ découpage du gène et du vecteur par la même enzyme de restriction $\rightarrow$ insertion et soudure par la ligase (ADN recombinant) $\rightarrow$ transformation de la cellule hôte $\rightarrow$ sélection des cellules transformées (marqueur de résistance à un antibiotique) $\rightarrow$ clonage et expression du gène.
\item La même enzyme de restriction génère des \textbf{extrémités complémentaires} qui s'apparient spontanément.
\item L'étape de \textbf{sélection} est indispensable car la transformation ne concerne qu'une minorité de cellules.
\item Chez les plantes, la \textbf{totipotence} permet de régénérer une plante entière transgénique à partir d'une seule cellule transformée.
\end{itemize}
\end{aretenir}

\begin{exercice}[Pour t'entraîner]
On souhaite obtenir des plants de maïs résistants à un insecte foreur en transférant le gène Bt d'une bactérie. 1) Liste dans l'ordre les étapes de cette transgénèse. 2) Explique pourquoi l'étape de sélection est nécessaire et propose un marqueur utilisable. 3) Quelle propriété des cellules végétales permet d'obtenir une plante entière à partir d'une cellule transformée ?
\end{exercice}

\begin{corrige}[Exercice]
1) Isolement du gène Bt chez la bactérie donneuse ; découpage du gène et d'un vecteur (plasmide) par la même enzyme de restriction ; insertion et soudure par la ligase (ADN recombinant) ; introduction dans des cellules de maïs (transformation, souvent via \emph{Agrobacterium tumefaciens}) ; sélection des cellules transformées ; régénération de plantes entières par culture \emph{in vitro}, puis multiplication.
2) La transformation ne réussit que dans une faible proportion de cellules : sans sélection, les cellules transformées seraient noyées parmi les cellules ordinaires. On utilise un gène de résistance à un antibiotique (ou à un herbicide) présent sur le vecteur : seules les cellules transformées survivent sur le milieu sélectif.
3) C'est la \textbf{totipotence} : toute cellule végétale somatique possède le génome complet et peut régénérer une plante entière sur un milieu de culture contenant des hormones de croissance (auxines et cytokinines).
\end{corrige}

Maintenant que le protocole de fabrication d'un OGM n'a plus de secret pour toi, il est temps de découvrir, dans la section suivante, la diversité des \emph{applications} concrètes du génie génétique, en médecine comme en agriculture.

\coursec{Applications du génie génétique}

Dans la section précédente, nous avons détaillé les \cle{étapes de l'obtention d'un organisme génétiquement modifié} : isolement du gène d'intérêt, coupure par les enzymes de restriction, insertion dans un vecteur grâce à la ligase, transformation de la cellule hôte, puis sélection et multiplication des organismes modifiés. Une question se pose alors naturellement : \emph{à quoi cela sert-il concrètement ?} Pourquoi les laboratoires du monde entier ont-ils investi dans ces techniques ? C'est ce que nous allons découvrir maintenant, en examinant les grandes applications du génie génétique, de la médecine à l'agriculture tropicale.

\courssub{Application médicale : la production d'insuline humaine par des bactéries transgéniques}

\textbf{Une situation concrète pour commencer.}\par
Au Centre Hospitalier d'Essos à Yaoundé, comme dans tous les hôpitaux du Cameroun, des patients diabétiques reçoivent chaque jour des injections d'\cle{insuline}. Le diabète de type 1 est une maladie dans laquelle les cellules \cle{bêta} des îlots de Langerhans du pancréas ne produisent plus d'insuline, l'hormone qui permet aux cellules d'absorber le glucose sanguin. Sans apport extérieur d'insuline, le patient meurt. Mais d'où vient cette insuline injectée ? Avant 1982, on l'extrayait du pancréas de porcs ou de bovins abattus dans les abattoirs. Aujourd'hui, elle est fabriquée par des \emph{bactéries}. Comment est-ce possible ?

\begin{definition}[Insuline et bactéries transgéniques]
L'\cle{insuline} est une hormone protéique (deux chaînes polypeptidiques liées par des ponts disulfure) sécrétée par le pancréas, indispensable à la régulation de la glycémie (taux de glucose sanguin). Une \cle{bactérie transgénique} est une bactérie (ici \emph{Escherichia coli}) dans le génome de laquelle on a inséré un gène étranger, ici le \textbf{gène humain de l'insuline} (souvent sous forme de cDNA codant les chaînes A et B séparément, puis assemblées \emph{in vitro}, ou sous forme d'un précurseur unique, la proinsuline). La bactérie devient alors une véritable « usine microscopique » : en se multipliant, elle produit massivement la protéine humaine.
\end{definition}

\textbf{Principe de la fabrication.}\par
On applique exactement les étapes vues dans la section précédente :
\begin{enumerate}
\item On \textbf{isole le gène humain de l'insuline} (ou on le synthétise chimiquement à partir de la séquence protéique connue) ;
\item on l'insère dans un \textbf{plasmide} (petit ADN circulaire bactérien) ouvert par une enzyme de restriction, puis refermé par la ligase : c'est l'ADN recombinant ; le plasmide porte un gène de résistance à un antibiotique (marqueur de sélection) ;
\item on introduit ce plasmide dans des bactéries \emph{E. coli} compétentes (\textbf{transformation}) ;
\item on cultive les bactéries modifiées dans de grands \textbf{fermenteurs} (cuves de plusieurs milliers de litres, milieu nutritif contrôlé, aération, température) où elles se multiplient très vite et expriment le gène recombinant pour produire l'insuline (souvent sous forme de corps d'inclusion qu'il faut solubiliser et replier, ou sécrétée dans le périplasme) ;
\item on \textbf{extrait et purifie} l'insuline (chromatographies successives), qui est ensuite conditionnée pour l'usage médical.
\end{enumerate}

\begin{popfigure}
\begin{center}
\begin{tikzpicture}[
  etape/.style={rectangle, rounded corners=3pt, draw=popblue, fill=popblueL, thick, minimum height=1.1cm, text width=3.1cm, align=center, font=\small},
  fleche/.style={-{Stealth[length=3mm]}, very thick, poporange}
]
\node[etape] (e1) at (0,0) {\textbf{1.} Gène humain de l'insuline (cDNA) inséré dans un plasmide};
\node[etape] (e2) at (4.6,0) {\textbf{2.} Transformation d'\emph{E. coli} (bactérie transgénique)};
\node[etape] (e3) at (9.2,0) {\textbf{3.} Fermenteur : multiplication et production massive};
\node[etape] (e4) at (4.6,-2.3) {\textbf{4.} Extraction, purification, repliement de l'insuline};
\node[etape, draw=popgreen, fill=popgreenL] (e5) at (0,-2.3) {\textbf{5.} Insuline humaine pure conditionnée (ampoules, stylos)};
\draw[fleche] (e1) -- (e2);
\draw[fleche] (e2) -- (e3);
\draw[fleche] (e3.south) -- ++(0,-0.6) -| (e4.east);
\draw[fleche] (e4) -- (e5);
\end{tikzpicture}
\end{center}
\legende{Chaîne de production de l'insuline humaine par des bactéries transgéniques : le gène humain (cDNA) inséré dans un plasmide transforme la bactérie en usine à insuline, cultivée en fermenteur ; l'hormone est ensuite extraite, purifiée et repliée.}
\end{popfigure}

\textbf{Quel est l'intérêt ? Comparons avec l'ancienne méthode.}\par
L'insuline de porc diffère de l'insuline humaine par un seul acide aminé (Thr en position B30 au lieu d'Ala), celle de bœuf par trois acides aminés. Cette petite différence suffit parfois à déclencher chez le patient des \textbf{réactions immunitaires} (production d'anticorps anti-insuline, allergies, lipodystrophies au point d'injection, diminution d'efficacité). De plus, l'extraction animale était coûteuse, peu rendement et limitée par l'approvisionnement en pancréas d'abattoirs.

\begin{exemplebox}[Un exemple chiffré éloquent]
Avant 1982, il fallait environ \textbf{\SI{8000}{kg} de pancréas de porc} pour produire seulement \textbf{\SI{1}{kg} d'insuline} ! Aujourd'hui, quelques fermenteurs peuplés de bactéries transgéniques suffisent à produire des quantités bien supérieures, en continu, partout dans le monde. En 1982, l'insuline humaine produite par \emph{E. coli} (\textbf{Humulin}) a été le \emph{premier médicament issu du génie génétique} autorisé sur le marché (FDA).
\end{exemplebox}

\begin{popfigure}
\begin{center}
\begin{tikzpicture}
\node[rectangle, draw=popdark, fill=popcream, thick, rounded corners=3pt, text width=0.94\linewidth, inner sep=8pt, font=\small] {%
\textbf{\textcolor{popdark}{Insuline bactérienne (transgénique)}}\quad vs\quad \textbf{\textcolor{popdark}{Insuline animale (porc/bœuf)}}\\[4pt]
\begin{tabular}{@{}p{0.44\linewidth}p{0.48\linewidth}@{}}
\textcolor{popgreen}{\rule{2.2mm}{2.2mm}} Structure \textbf{strictement identique} à l'insuline humaine & \textcolor{poporange}{\rule{2.2mm}{2.2mm}} Diffère par 1 à 3 acides aminés \\[2pt]
\textcolor{popgreen}{\rule{2.2mm}{2.2mm}} \textbf{Très peu de rejets} ni d'allergies & \textcolor{poporange}{\rule{2.2mm}{2.2mm}} Réactions immunitaires possibles (anticorps, allergies, lipodystrophies) \\[2pt]
\textcolor{popgreen}{\rule{2.2mm}{2.2mm}} Production \textbf{massive et continue} en fermenteurs & \textcolor{poporange}{\rule{2.2mm}{2.2mm}} Production limitée : \SI{8000}{kg} de pancréas pour \SI{1}{kg} d'insuline \\[2pt]
\textcolor{popgreen}{\rule{2.2mm}{2.2mm}} Coût réduit, disponibilité mondiale & \textcolor{poporange}{\rule{2.2mm}{2.2mm}} Dépendance aux abattoirs, coût élevé \\
\end{tabular}};
\end{tikzpicture}
\end{center}
\legende{Comparaison des avantages de l'insuline produite par bactéries transgéniques (colonne de gauche) avec l'insuline extraite de pancréas de porc ou de bœuf (colonne de droite).}
\end{popfigure}

\courssub{Autres productions pharmaceutiques}

L'insuline a ouvert la voie à de nombreuses autres \cle{protéines médicaments} (biomédicaments) produites par génie génétique :
\begin{itemize}
\item l'\textbf{hormone de croissance} humaine (somatotropine, \emph{somatropine}), utilisée pour traiter les retards de croissance d'origine hormonale ; auparavant extraite d'hypophyses de cadavres, avec un risque de transmission de maladies à prions (maladie de Creutzfeldt-Jakob) ;
\item des \textbf{vaccins recombinants}, comme le vaccin contre l'hépatite B (protéine de surface HBsAg produite par levures \emph{Saccharomyces cerevisiae} transgéniques) — sans aucun risque d'infection, puisqu'aucun virus complet n'est utilisé ;
\item des \textbf{facteurs de coagulation} (facteur VIII et IX recombinants) pour les hémophiles, autrefois extraits de concentrés de sang de donneurs avec un risque majeur de contamination virale (VIH, VHC dans les années 1980) ;
\item des \textbf{anticorps monoclonaux} humanisés et des enzymes thérapeutiques (ex. : tissue plasminogen activator -- tPA pour l'infarctus).
\end{itemize}

\begin{aretenir}[L'idée-force des applications médicales]
Le génie génétique permet de faire produire par des micro-organismes (bactéries \emph{E. coli}, levures \emph{S. cerevisiae}, cellules CHO pour les protéines complexes nécessitant des glycosylations) des \textbf{protéines humaines pures}, en grande quantité, à moindre coût et sans risque de contamination pathogène ni de rejet immunitaire lié à des protéines hétérologues. C'est le principe des \cle{biomédicaments}.
\end{aretenir}

\courssub{Application agricole : l'amélioration variétale}

\textbf{Problème concret.}\par
Dans les champs de maïs du Nord-Cameroun (régions de l'Adamaoua, du Nord, de l'Extrême-Nord), un agriculteur peut perdre une part significative de sa récolte à cause d'insectes ravageurs (chenilles foreuses de tiges : \emph{Busseola fusca}, \emph{Sesamia calamistis}, \emph{Chilo partellus}). Il doit alors acheter des insecticides chimiques : coûteux, parfois dangereux pour sa santé (manipulation sans EPI) et polluants pour les sols, les nappes phréatiques et la biodiversité auxiliaire. Le génie génétique propose une autre solution : et si la plante se défendait elle-même ?

\begin{definition}[Plante transgénique]
Une \cle{plante transgénique} est une plante dans le génome nucléaire de laquelle on a inséré un ou plusieurs gènes étrangers (transgènes) lui conférant un caractère nouveau et utile : résistance à un insecte (gène \emph{cry}), tolérance à un herbicide (gène \emph{epsps} muté ou \emph{bar}), meilleure valeur nutritionnelle, tolérance à la sécheresse, etc. L'insertion se fait généralement via \emph{Agrobacterium tumefaciens} ou par biolistique.
\end{definition}

\textbf{Exemple 1 : le maïs Bt, résistant aux insectes.}\par
\emph{Bacillus thuringiensis} (Bt) est une bactérie du sol qui produit naturellement des protéines cristallines (Cry) toxiques pour certains ordres d'insectes (Lépidoptères, Coléoptères, Diptères) mais inoffensives pour l'homme, les mammifères et la plupart des auxiliaires (absence de récepteurs spécifiques dans l'intestin). Les chercheurs ont \textbf{transféré un gène \emph{cry} (ex. \emph{cry1Ab}, \emph{cry1F}) dans le génome du maïs}. Le maïs Bt obtenu fabrique lui-même la protéine insecticide dans ses tissus (feuilles, tiges, épis) : la chenille qui le mange ingère la protoxine, qui est activée dans son tube digestif alcalin, se lie à des récepteurs spécifiques de la membrane des cellules intestinales, forme des pores et provoque la lyse cellulaire et la mort de l'insecte, sans qu'il soit nécessaire de pulvériser d'insecticide.

\begin{savaistu}[Le saviez-tu ?]
Le maïs Bt produit lui-même une protéine insecticide issue d'une bactérie du sol, \emph{Bacillus thuringiensis}. Cette bactérie était d'ailleurs utilisée depuis longtemps en agriculture biologique sous forme de pulvérisations de spores et cristaux ! Grâce au maïs Bt, l'usage d'insecticides chimiques de synthèse a fortement diminué dans les cultures concernées (ex. : maïs, coton), ce qui réduit les coûts pour l'agriculteur, l'exposition des applicateurs et la pollution de l'environnement. Au Cameroun, des essais de plein champ (confined field trials) ont été menés sur le coton Bt et le maïs Bt par l'IRAD (Institut de Recherche Agricole pour le Développement).
\end{savaistu}

\textbf{Exemple 2 : les plantes tolérantes aux herbicides.}\par
Certaines variétés de soja, maïs, colza ou coton ont reçu un gène les rendant \textbf{tolérantes à un herbicide à large spectre} (principalement le glyphosate : gène \emph{cp4 epsps} d'\emph{Agrobacterium} sp. souche CP4 ; ou le glufosinate : gène \emph{bar} de \emph{Streptomyces}). L'agriculteur peut ainsi désherber son champ \emph{après} la levée de la culture en détruisant les mauvaises herbes sans abîmer la culture (désherbage post-levée simplifié). Ce type d'OGM est très répandu (majorité des surfaces OGM mondiales) mais aussi très débattu, car il peut encourager un usage massif et répété du même herbicide, favorisant l'apparition de mauvaises herbes résistantes (voir section 5).

\textbf{Exemple 3 : le riz doré, à meilleure valeur nutritionnelle.}\par
La carence en vitamine A (xérophtalmie, affaiblissement immunitaire, mortalité infantile) est un grave problème de santé publique dans de nombreux pays où l'alimentation repose sur le riz blanc (pauvre en bêta-carotène). Le \cle{riz doré} (Golden Rice) est un riz transgénique dans lequel on a inséré deux gènes codant des enzymes de la voie de biosynthèse des caroténoïdes : la \emph{phytoène synthase} (originellement du maïs, puis du riz lui-même) et la \emph{carotène desaturase} (crtI, de la bactérie \emph{Erwinia uredovora} ou \emph{Pantoea ananatis}). Le grain de riz accumule ainsi du \textbf{bêta-carotène} (provitamine A), d'où sa couleur jaune-orangé. Consommé régulièrement, il peut aider à lutter contre la carence en vitamine A. La variété GR2E a reçu des autorisations de sécurité alimentaire (Australie, Nouvelle-Zélande, USA, Philippines).

\textbf{Intérêt particulier pour les pays tropicaux (ex. Cameroun).}\par
Pour des pays comme le Cameroun, l'amélioration variétale par génie génétique (ou par sélection assistée par marqueurs, plus acceptée réglementairement) présente des enjeux majeurs :
\begin{itemize}
\item obtenir des variétés \textbf{résistantes aux ravageurs} locaux (foreurs du maïs, mouches des fruits \emph{Bactrocera} sur manguiers/agrumes, charançons du bananier) et aux \textbf{maladies} virales (mosaïque du manioc -- CMD, streak du maïs -- MSV, bunchy top du bananier -- BBTV) ;
\item créer des variétés \textbf{tolérantes à la sécheresse} et à la chaleur, cruciales dans les zones soudano-sahéliennes du Grand Nord où les pluies sont irrégulières et le changement climatique s'accentue ;
\item \textbf{réduire les pertes de récoltes} (avant et après récolte), améliorer les rendements et donc la sécurité alimentaire d'une population en forte croissance ;
\item réduire la dépendance aux intrants chimiques coûteux et importés (devises).
\end{itemize}

\courssub{Autres domaines d'application}

\begin{itemize}
\item \textbf{Dépollution (bioremédiation)} : des bactéries génétiquement modifiées capables de \textbf{dégrader des polluants récalcitrants} (hydrocarbures aromatiques polycycliques après une marée noire, BTEX, métaux lourds par biosorption/bioaccumulation, pesticides organochlorés) sont étudiées pour assainir les sols et les eaux contaminées (ex. : \emph{Pseudomonas putida} modifiée).
\item \textbf{Recherche fondamentale} : insérer (transgénèse), inactiver (\emph{knock-out}), modifier (\emph{knock-in}) ou marquer (GFP, luciferase) des gènes chez des organismes modèles (souris, drosophile, zébrifoin, \emph{Arabidopsis}) permet de comprendre le rôle de chaque gène, les réseaux de régulation et les mécanismes physiopathologiques des maladies humaines.
\item \textbf{Thérapie génique (perspective thérapeutique)} : il s'agit d'introduire un gène fonctionnel (ou d'éditer le génome via CRISPR-Cas9) dans les cellules somatiques d'un patient pour \textbf{corriger une maladie génétique} (mucoviscidose -- gène \emph{CFTR}, amyotrophie spinale -- gène \emph{SMN1}, certaines immunodéficiences combinées sévères -- SCID, drépanocytose, hémophilie B). Des médicaments de thérapie génique sont désormais autorisés (ex. : Zolgensma, Luxturna, Casgevy), mais la technique reste coûteuse, techniquement délicate (vecteurs viraux AAV, lentivirus ; \emph{ex vivo} vs \emph{in vivo}) et éthiquement encadrée (interdiction de la lignée germinale).
\item \textbf{Industrie (biotechnologies blanches)} : production d'enzymes industrielles (protéases, amylases, lipases pour les lessives ; chymosine pour la fabrication du fromage -- remplacement de la présure animale ; cellulases pour bioéthanol 2G) par des micro-organismes modifiés (souches surproductrices, sécrétion optimisée).
\end{itemize}

\courssub{Savoir-faire : expliquer l'intérêt d'un organisme génétiquement modifié}

Au Probatoire, on te donnera souvent un \textbf{document} (schéma, texte, tableau de données, graphique) décrivant un OGM, et on te demandera d'en \emph{expliquer l'intérêt}. Voici la démarche rigoureuse à suivre.

\begin{methode}[Expliquer l'intérêt d'un OGM à partir d'un document]
\begin{enumerate}
\item \textbf{Identifier le problème posé} : maladie humaine ou animale, ravageur agricole, carence nutritionnelle, pollution environnementale, coût de production élevé, limite de la sélection classique...
\item \textbf{Repérer le(s) gène(s) transféré(s)} : quelle est son origine (organisme donneur) ? Quelle protéine ou quel ARN code-t-il ? Quel est son mode d'action connu ?
\item \textbf{Décrire le caractère nouveau obtenu} (phénotype) chez l'organisme modifié grâce à l'expression de ce transgène (production d'une protéine d'intérêt, résistance à un stress biotique/abiotique, tolérance à un herbicide, modification métabolique...).
\item \textbf{Conclure par le(s) bénéfice(s)} pour l'homme, l'agriculture, la santé ou l'environnement : santé améliorée (qualité, sécurité, coût), rendements accrus et stabilisés, réduction des intrants chimiques (pesticides, engrais), valeur nutritionnelle accrue, dépollution, connaissance fondamentale...
\end{enumerate}
Articule toujours ta réponse selon la chaîne logique : \emph{problème → gène transféré → caractère obtenu → bénéfice(s)}. Utilise le vocabulaire précis : \cle{transgène}, \cle{protéine recombinante}, \cle{phénotype}, \cle{résistance}, \cle{tolérance}.
\end{methode}

\begin{exoresolu}[Le maïs Bt : analyse d'un document]
\textbf{Énoncé.} Un document présente le maïs Bt, obtenu en insérant dans le génome du maïs un gène \emph{cry1Ab} de la bactérie du sol \emph{Bacillus thuringiensis}, codant une protéine Cry1Ab toxique pour les chenilles foreuses (Lépidoptères). Des essais au champ comparatifs (maïs Bt vs maïs isogénique non-Bt traité ou non) montrent une baisse de 60 \% du nombre de traitements insecticides sur les parcelles Bt et une augmentation de rendement de 15 \%. Explique l'intérêt de cet OGM.
\tcblower
\textbf{Solution détaillée.} On suit la démarche en quatre points :
\begin{enumerate}
\item \emph{Problème posé} : les chenilles foreuses (\emph{Busseola fusca}, \emph{Sesamia} spp.) détruisent une partie importante des récoltes de maïs en Afrique subsaharienne ; la lutte chimique par insecticides de synthèse est coûteuse pour le petit paysan, techniquement difficile (chenilles dans la tige), dangereuse pour la santé (manipulation) et polluante pour l'environnement (eaux, sols, auxiliaires).
\item \emph{Gène transféré} : le gène \emph{cry1Ab} (originaire de \emph{Bacillus thuringiensis} souche HD-1) a été inséré dans le génome nucléaire du maïs sous le contrôle d'un promoteur constitutif (ex. : promoteur 35S du CaMV ou promoteur ubiquitine du maïs) ou spécifique des tissus verts.
\item \emph{Caractère obtenu} : le maïs Bt exprime la protéine Cry1Ab dans ses tissus (feuilles, tiges, épis) ; les chenilles foreuses qui s'en nourrissent ingèrent la protéine, qui est activée dans leur tube digestif alcalin et provoque leur mort par lyse des cellules intestinales. La plante est donc \textbf{résistante aux foreurs} (protection intégrée, continue, sans application externe).
\item \emph{Bénéfices} : (1) \textbf{Agronomique/Economique} : les pertes de rendement dues aux foreurs diminuent fortement (gain de 15 \% ici), sécurité de la récolte accrue ; (2) \textbf{Environnemental/Sanitaire} : l'usage d'insecticides chimiques chute de 60 \%, réduisant l'exposition de l'agriculteur, la contamination des eaux/sols et l'impact sur la faune auxiliaire (pollinisateurs, prédateurs) ; (3) \textbf{Social} : réduction de la pénibilité et du coût des intrants.
\end{enumerate}
\emph{Conclusion} : le maïs Bt est un OGM dont l'intérêt est multiple — économique (rendement, coût), environnemental (moins de pesticides), sanitaire (moins d'exposition) et social (accessibilité pour petits producteurs).
\end{exoresolu}

\begin{attention}[Erreurs fréquentes à éviter au Probatoire]
\begin{itemize}
\item Ne dis jamais qu'un OGM est un organisme « monstrueux », « totalement transformé » ou « hybride animal-végétal ». Le gène introduit ajoute généralement \textbf{un seul caractère précis} (par exemple : produire de l'insuline, résister à un insecte, tolérer un herbicide). Le maïs Bt reste du maïs (même morphologie, même grain, même valeur nutritive de base), la bactérie productrice d'insuline reste \emph{E. coli} K12 (souche de laboratoire non pathogène).
\item Ne confonds pas \emph{gène transféré} (ADN, information héréditaire) et \emph{protéine produite} (produit d'expression, effecteur du caractère). Le gène est le support de l'information, la protéine est le produit fonctionnel qui confère le phénotype.
\item Ne confonds pas \emph{transgénèse} (ajout gène étranger) et \emph{édition de génome} (CRISPR-Cas9 : modification ciblée d'un gène endogène sans forcément ajouter de gène étranger). Les deux relèvent du génie génétique mais les réglementations diffèrent.
\item Attention au vocabulaire : on dit « \textbf{plante transgénique} » (ou \textbf{OGM végétal}), pas « plante mutante » (mutation = aléa, non dirigé).
\end{itemize}
\end{attention}

\begin{exercice}[À toi de jouer]
Le riz doré (Golden Rice, événement GR2E) est un riz transgénique contenant deux gènes étrangers (\emph{psy} du maïs et \emph{crtI} de \emph{Pantoea ananatis}) permettant au grain (endosperme) de synthétiser du bêta-carotène (provitamine A). Dans certaines régions d'Asie et d'Afrique, des centaines de milliers d'enfants souffrent de carence en vitamine A (xérophtalmie, mortalité). En suivant la méthode en quatre points, explique l'intérêt de cet OGM.
\end{exercice}

\begin{corrige}[Exercice : le riz doré (Golden Rice)]
\begin{enumerate}
\item \emph{Problème} : la carence en vitamine A (hypovitaminose A) est un problème majeur de santé publique dans les populations dont l'alimentation de base est le riz blanc (pauvre en micronutriments). Elle provoque la cécité nocturne puis totale (xérophtalmie), affaiblit gravement les défenses immunitaires (augmentation mortalité rougeole, diarrhées, paludisme) et touche des centaines de millions d'enfants et de femmes enceintes.
\item \emph{Gènes transférés} : deux gènes étrangers codant les enzymes clés de la voie de biosynthèse des caroténoïdes absente de l'endosperme du riz sauvage : la \emph{phytoène synthase} (\emph{psy}, gène du maïs ou du riz) et la \emph{phytoène désaturase} (\emph{crtI}, gène de la bactérie \emph{Pantoea ananatis} -- ex \emph{Erwinia uredovora}), placés sous promoteurs spécifiques de l'endosperme.
\item \emph{Caractère obtenu} : l'endosperme du riz transgénique accumule du \textbf{bêta-carotène} (provitamine A) et d'autres caroténoïdes, ce qui lui donne une couleur jaune-orangé caractéristique (« doré »). La teneur en bêta-carotène est suffisante pour couvrir les besoins quotidiens en vitamine A avec une ration normale de riz cuit (biodisponibilité démontrée).
\item \emph{Bénéfice} : la consommation quotidienne de ce riz par les populations à risque apporte de la provitamine A, transformée en vitamine A (rétinol) par l'organisme selon ses besoins, ce qui aide à \textbf{prévenir la carence et ses conséquences graves} (cécité, mortalité infantile, immunodépression) de manière durable, peu coûteuse et culturellement acceptée (aliment de base), en complément des programmes de supplémentation et de diversification alimentaire.
\end{enumerate}
\end{corrige}

Ainsi, le génie génétique offre des applications remarquables en médecine (biomédicaments, vaccins, thérapie génique), en agriculture (résistance ravageurs/maladies, tolérance herbicides/sécheresse, biofortification) et dans l'environnement (bioremédiation). Mais ces techniques puissantes soulèvent aussi des questions légitimes : les OGM sont-ils sans risque pour la santé et l'environnement ? Comment coexistent-ils avec l'agriculture conventionnelle et biologique ? Qui les contrôle, qui en bénéficie ? Quelles sont les questions éthiques ? C'est l'objet de la dernière section, consacrée aux \emph{enjeux et débats} autour des OGM.

\coursec{Enjeux et débats autour des OGM}

Dans les sections précédentes, nous avons découvert les applications concrètes du génie génétique : production d'insuline humaine par des bactéries transgéniques, création de variétés végétales améliorées (coton Bt résistant aux insectes, riz doré enrichi en provitamine A). Ces réussites technologiques sont indéniables. Pourtant, dès qu'un OGM est libéré dans l'environnement ou introduit dans l'alimentation, des questions légitimes se posent : est-ce sans danger pour la santé ? Pour les écosystèmes ? Pour les paysans ? C'est tout l'objet de cette dernière section : apprendre à \cle{peser les avantages et les risques} du génie génétique avec rigueur, sans enthousiasme aveugle ni peur irraisonnée.

\courssub{Les avantages attendus des OGM}

\textbf{Situation concrète.} Dans la plaine des Mbô, un planteur de maïs perd chaque année une partie de sa récolte à cause de la pyrale, un insecte foreur de tiges. Il doit acheter des insecticides coûteux, qu'il pulvérise parfois sans protection. À quelques kilomètres, un essai de maïs Bt montre des plants intacts, sans aucune pulvérisation. Que s'est-il passé ?

Le maïs \cle{Bt} possède un gène de la bactérie \emph{Bacillus thuringiensis} qui code une protéine (toxine Cry) toxique pour les larves de certains insectes ravageurs (lépidoptères), mais inoffensive pour l'Homme et les vertébrés (absence de récepteurs spécifiques, dégradation rapide dans l'estomac acide). L'observation de ce type de résultats expérimentaux permet de dégager les principaux avantages attendus des OGM :

\begin{itemize}
\item \textbf{Augmentation des rendements agricoles} : les variétés résistantes aux ravageurs, aux maladies ou tolérantes à la sécheresse limitent les pertes de récolte, ce qui accroît la production sur une même surface.
\item \textbf{Réduction de l'usage des pesticides} : une plante qui produit elle-même sa défense (comme le maïs Bt ou le coton Bt) nécessite moins de traitements insecticides, ce qui diminue les coûts, la pollution des sols et des eaux, et les risques d'intoxication aiguë ou chronique des agriculteurs.
\item \textbf{Production de médicaments à moindre coût et plus sûrs} : l'insuline humaine produite par des bactéries (\emph{E. coli}) ou des levures transgéniques est aujourd'hui bien moins chère, plus pure et exemptée de risques de transmission de pathogènes animaux, contrairement à l'insuline extraite autrefois de pancréas porcins ou bovins. Hormones de croissance, vaccins (hépatite B), facteurs de coagulation sont obtenus par des procédés similaires.
\item \textbf{Lutte contre les carences alimentaires} : le \cle{riz doré} (Golden Rice), enrichi en bêta-carotène (provitamine A) grâce à l'introduction de gènes du pathway des caroténoïdes (phytoène synthase du maïs, désaturase de la bactérie \emph{Erwinia uredovora}), vise à combattre la carence en vitamine A, responsable de cécité et de mortalité infantile dans plusieurs pays — un enjeu de santé publique réel, y compris en Afrique.
\end{itemize}

\begin{exemplebox}[Le coton Bt, un cas africain instructif]
Au Burkina Faso, voisin du Cameroun, le coton Bt (variétés Bollgard II) a été cultivé à grande échelle à partir de 2008 : les rendements ont augmenté de 15 à 20 \% et le nombre de traitements insecticides a chuté de 6-8 à 2 par campagne. Mais des contestations sur la qualité de la fibre (longueur légèrement inférieure) et le prix élevé des semences (redevance technologique) ont conduit le pays à suspendre sa culture en 2016 pour revenir au coton conventionnel. Cet exemple montre qu'un OGM n'est jamais qu'un outil, dont le bilan global dépend du contexte agronomique, économique et social.
\end{exemplebox}

\courssub{Les risques sanitaires débattus}

\textbf{Question de santé publique.} Si une protéine nouvelle, issue d'un gène étranger, se retrouve dans notre assiette, peut-elle provoquer des allergies ou des effets à long terme ?

\begin{itemize}
\item \textbf{Allergies possibles} : le transgène code une protéine nouvelle dans l'aliment. Si cette protéine est \cle{allergène}, elle peut déclencher des réactions chez les personnes sensibilisées. C'est pourquoi les protocoles d'évaluation (Codex Alimentarius) imposent, avant toute mise sur le marché, de vérifier que la protéine produite ne présente pas d'homologie de séquence avec des allergènes connus, qu'elle est rapidement dégradée par la pepsine (simulation digestion gastrique) et qu'elle est thermolabile. Un cas célèbre : un projet de soja enrichi en méthionine avec un gène de noix du Brésil (\emph{Bertholletia excelsa}) a été abandonné dès les tests pré-cliniques, car la protéine 2S albumine introduite s'est révélée être l'allergène majeur de la noix du Brésil — preuve que le système d'évaluation \emph{ex ante} fonctionne.
\item \textbf{Effets à long terme insuffisamment évalués} : la plupart des études toxicologiques réglementaires portent sur 90 jours chez le rongeur (étude subchronique). Les effets éventuels d'une consommation sur des décennies (cancérogénicité, effets transgénérationnels) sont difficiles à mesurer, ce qui alimente le débat. À ce jour, les agences sanitaires internationales (OMS, FAO, EFSA, ANSES) considèrent, sur la base des données disponibles, que les OGM autorisés ne présentent pas de risque démontré supérieur à leurs équivalents conventionnels, mais la surveillance post-commercialisation (\emph{biovigilance}) reste nécessaire.
\end{itemize}

\begin{attention}[Ni danger automatique, ni produit miracle]
Erreur fréquente au Probatoire : affirmer « les OGM sont dangereux » ou, inversement, « les OGM vont nourrir le monde ». Ces deux slogans sont scientifiquement faux. Chaque OGM est un cas particulier : le risque dépend du \textbf{gène introduit} (nature de la protéine), de la \textbf{plante hôte} (effets de position, métabolisme), de l'\textbf{usage} (aliment, fourrage, médicament) et du \textbf{milieu} de dissémination. L'évaluation se fait donc \cle{au cas par cas}, jamais en bloc.
\end{attention}

\courssub{Les risques environnementaux}

\textbf{Observation.} Le sorgho (\emph{Sorghum bicolor}) est cultivé au Nord-Cameroun. Des sorghos sauvages apparentés (\emph{Sorghum halepense}, \emph{Sorghum verticilliflorum}) poussent en bordure des champs. Le pollen voyage avec le vent. Que risque-t-il de se passer si le sorgho cultivé est transgénique (ex. tolérance à un herbicide) ?

\textbf{Hypothèse} : le transgène pourrait passer de la culture vers les plantes sauvages apparentées par pollinisation croisée (hybridation).

\textbf{Données expérimentales} : des études de terrain (notamment sur le colza en Europe, le sorgho et le riz en Afrique) ont effectivement détecté, chez des plantes sauvages poussant près de cultures OGM, la présence du transgène (ex. résistance à l'herbicide glyphosate ou au glufosinate). Ce phénomène est appelé \cle{flux de gènes} (ou \emph{introgression}).

\textbf{Interprétation et conclusion} : le gène introduit par l'Homme ne reste pas confiné dans la culture ; il peut se disséminer dans les populations sauvages ou adventices. Trois grands risques environnementaux sont ainsi identifiés :

\begin{itemize}
\item \textbf{Dissémination des transgènes (flux de gènes)} : le transfert d'un gène de résistance (à un herbicide, à un insecte) vers des espèces sauvages ou adventices apparentées sexuellement compatibles peut créer des « super-adventices » difficiles à éliminer (ex. résistance à l'herbicide transférée au sorgho sauvage).
\item \textbf{Apparition de ravageurs résistants (évolution de résistance)} : exposés en permanence à la toxine Bt dans la plante, les insectes ravageurs subissent une \cle{pression de sélection} intense : les rares individus porteurs d'allèles de résistance (récepteurs modifiés, détoxification accrue) survivent et se reproduisent. Des populations de pyrale du maïs (\emph{Ostrinia nubilalis}) ou de ver de la capsule du coton (\emph{Helicoverpa armigera}) résistantes à la toxine Bt ont déjà été signalées dans plusieurs pays, ce qui rend l'OGM inefficace à terme. La stratégie des \emph{refuges} (zones de culture non-Bt) vise à retarder ce phénomène.
\item \textbf{Réduction de la biodiversité cultivée} : la généralisation de quelques variétés transgéniques uniformes, souvent hybrides F1 non ressemables, peut évincer les variétés locales traditionnelles (les nombreux cultivars de maïs, de sorgho, de manioc, d'igname, de haricot du Cameroun), appauvrissant le patrimoine génétique pourtant précieux pour l'avenir de la sélection face au changement climatique.
\end{itemize}

\begin{popfigure}
\begin{center}
\begin{tikzpicture}[scale=1, every node/.style={font=\small}]
% Champ OGM (Sorgho)
\fill[popgreenL] (-4.2,-1.6) rectangle (-0.4,1.6);
\node[popdark, font=\small\bfseries] at (-2.3,1.9) {Champ de sorgho OGM (tolérance herbicide)};
\foreach \x in {-3.8,-3.0,-2.2,-1.4,-0.6}{
  \draw[popgreen!70!black, thick] (\x,-1.2) -- (\x,0.8);
  \draw[popgreen!70!black] (\x,0.8) -- ++(-0.15,0.25);
  \draw[popgreen!70!black] (\x,0.8) -- ++(0.15,0.25);
  \fill[popgold] (\x,1.15) circle (0.08);
}
\node[popmuted, font=\footnotesize] at (-2.3,-1.35) {plantes cultivées transgéniques};
% Zone intermédiaire : pollen
\foreach \p in {(0.4,1.0),(1.0,0.4),(1.6,1.2),(2.2,0.6),(1.3,-0.3),(2.6,0.0)}{
  \fill[poporange] \p circle (0.06);
}
\draw[->, poporange, thick, decorate, decoration={snake, amplitude=0.5mm, segment length=2mm}] (-0.2,0.8) -- (3.0,0.8);
\node[poporange, font=\footnotesize, align=center] at (1.5,1.45) {pollen transporté\\ par le vent (anémogamie)};
% Plantes sauvages
\fill[popblueL] (3.4,-1.6) rectangle (6.6,1.6);
\node[popdark, font=\small\bfseries] at (5.0,1.9) {Sorgho sauvage apparenté};
\foreach \x in {3.9,4.7,5.5,6.2}{
  \draw[popblue!70!black, thick] (\x,-1.2) -- (\x,0.5);
  \draw[popblue!70!black] (\x,0.5) -- ++(-0.14,0.25);
  \draw[popblue!70!black] (\x,0.5) -- ++(0.14,0.25);
}
% Une plante sauvage a reçu le transgène (hachurée)
\fill[poppink!50!white, pattern=north east lines, pattern color=poppink] (4.7,-0.2) rectangle (4.95,0.7);
\draw[popblue!70!black, thick] (4.825,-0.2) -- (4.825,0.5);
\draw[popblue!70!black] (4.825,0.5) -- ++(-0.14,0.25);
\draw[popblue!70!black] (4.825,0.5) -- ++(0.14,0.25);
\fill[poppink] (4.825,0.85) circle (0.08);
\node[poppink, font=\footnotesize, align=center] at (5.0,-1.3) {hybride fertile :\\ transgène introgressé};
% Flèches légendes
\draw[thin, popmuted, ->] (4.825,0.85) -- (6.9,0.9) node[right, font=\footnotesize, text=popdark] {fécondation croisée + introgression};
\draw[thin, popmuted, ->] (-2.2,1.15) -- (-2.2,2.5) node[above, font=\footnotesize, text=popdark] {graine OGM};
\end{tikzpicture}
\end{center}
\legende{Flux de gènes possible entre une culture OGM (sorgho tolérant un herbicide) et une plante sauvage apparentée (sorgho halepense) : le pollen de la plante transgénique, transporté par le vent, peut féconder une plante sauvage voisine. L'hybride fertile issu de cette fécondation intègre le transgène dans le patrimoine génétique de la population sauvage (introgression).}
\end{popfigure}

\courssub{Les enjeux économiques, sociaux et éthiques}

\textbf{Enjeux économiques et sociaux.} Les semences OGM sont protégées par des \cle{brevets} (propriété intellectuelle) détenus par de grandes firmes semencières multinationales. Concrètement, l'agriculteur signe un contrat de licence : il doit racheter les semences chaque année (les variétés sont souvent des hybrides F1 dont la descendance ségrège, ou contractuellement interdites de ressemis), et s'interdit de ressemer une partie de sa récolte, pratique pourtant millénaire et base de la sécurité alimentaire paysanne. Pour les paysans camerounais, dont la majorité pratique une agriculture familiale de subsistance et conserve leurs propres semences de maïs, de sorgho, de haricot ou d'arachide d'une saison à l'autre, cette dépendance économique et la perte de l'autonomie semencière sont des questions sensibles : qui contrôle la semence contrôle la production alimentaire. À l'inverse, des semences plus performantes (rendement, tolérance stress) peuvent améliorer les revenus si leur prix reste accessible et si l'agriculteur conserve le choix.

\textbf{Enjeux éthiques.} Jusqu'où l'Homme peut-il manipuler le vivant ? Insérer un gène de bactérie dans une plante, c'est franchir la barrière des espèces (transgenèse), ce que la sélection classique (croisements interspécifiques) ne permet pas ou très difficilement. Certains y voient un progrès maîtrisé au service de l'humanité, d'autres une transgression des limites naturelles ou une marchandisation inacceptable du vivant. Ces questions ne relèvent pas seulement de la science : elles appellent un débat citoyen démocratique et une réglementation claire, incluant la \cle{biosécurité} et l'\cle{étiquetage} obligatoire des produits (seuil 0,9 \% en UE), pour que le consommateur et l'agriculteur puissent choisir en connaissance de cause.

\begin{definition}[Principe de précaution]
Le \cle{principe de précaution} (inscrit dans la Constitution camerounaise de 1996, loi n°96/12, et la Charte de l'environnement) impose d'évaluer et de limiter les risques d'une technologie avant sa diffusion, même en l'absence de certitude scientifique totale sur l'étendue du danger. Autrement dit : face à un risque potentiel grave ou irréversible pour la santé ou l'environnement, l'absence de preuve complète du danger ne doit pas servir de prétexte pour ne rien faire. Il justifie l'évaluation rigoureuse des OGM \emph{avant} leur mise sur le marché (autorisation cas par cas), et non après les premiers accidents.
\end{definition}

\courssub{Le cadre réglementaire : de Cartagena au Cameroun}

Face à ces enjeux, la communauté internationale a mis en place des règles de \cle{biosécurité} : ensemble de mesures juridiques et techniques visant à encadrer la manipulation confinée, l'échange transfrontalier et la dissémination volontaire des OGM afin de protéger la santé humaine et la biodiversité.

\begin{savaistu}[Le Protocole de Cartagena et le Cameroun]
Le \textbf{Protocole de Cartagena sur la prévention des risques biotechnologiques}, adopté en 2000 à Carthagène (Colombie) dans le cadre de la Convention sur la diversité biologique (CDB), encadre les échanges internationaux d'OGM destinés à une introduction dans l'environnement (semences, plantes vivantes) : tout exportateur doit notifier le pays importateur, qui dispose d'un délai pour évaluer les risques et accepter ou refuser la marchandise (procédure d'« accord préalable en connaissance de cause » - APRC). Le Cameroun a ratifié ce protocole en 2002 et s'est doté de la \textbf{loi n°2003/006 du 21 avril 2003} fixant le régime de la biosécurité, complétée par des décrets d'application. Concrètement, aucun OGM ne peut être importé, cultivé en plein champ ou commercialisé au Cameroun sans autorisation préalable du Ministère de l'Environnement, après avis du Comité National de Biosécurité.
\end{savaistu}

Le tableau suivant synthétise les arguments structurés du débat.

\begin{popfigure}
\begin{center}
\begin{tikzpicture}
\node[font=\small, align=center]{
\begin{tabular}{|p{2.8cm}|p{5.0cm}|p{5.0cm}|}
\hline
\rowcolor{popblueL} \textbf{Domaine} & \textbf{Avantages attendus (Exemples)} & \textbf{Risques ou limites (Exemples)} \\
\hline
Sanitaire & Médicaments recombinants purs, moins chers, illimités (insuline, vaccin Hb) ; aliments biofortifiés (riz doré : provitamine A) & Allergénicité potentielle de la nouvelle protéine (testée \emph{ex ante}) ; incertitudes toxicologiques long terme (biovigilance nécessaire) \\
\hline
Environnemental & Réduction drastique insecticides (coton/maïs Bt) ; moindre labour (tolérance herbicide) $\rightarrow$ moins érosion, moins CO$_2$ & Flux de gènes vers sauvages/adventices (sorgho, riz) ; résistance ravageurs (sélection Bt) ; érosion variétale (remplacement cultivars locaux) \\
\hline
Économique et social & Hausse rendement/revenu possible ; baisse coûts traitements ; pénibilité réduite & Brevets/semences hybrides $\rightarrow$ dépendance semencière ; coût semences élevé ; asymétrie pouvoir firmes/paysans \\
\hline
Éthique et juridique & Progrès technique maîtrisé ; réponse besoins santé/alimentation & Transgression barrières espèces ; brevetage du vivant ; droit du consommateur (étiquetage) ; justice Nord/Sud (accès techno) \\
\hline
\end{tabular}
};
\end{tikzpicture}
\end{center}
\legende{Tableau de synthèse à double entrée : avantages et risques des OGM selon quatre domaines (sanitaire, environnemental, économique/social, éthique/juridique).}
\end{popfigure}

\courssub{Construire un avis argumenté (savoir-faire exigible au Probatoire)}

Au Probatoire, on peut te demander de \emph{discuter des avantages et risques du génie génétique} ou d'\emph{analyser un document sur les OGM}. Il ne s'agit pas de réciter un avis personnel, mais de construire un raisonnement scientifique structuré, appuyé sur des données précises.

\begin{methode}[Discuter des avantages et risques des OGM en quatre temps]
\begin{enumerate}
\item \textbf{Rappeler les faits scientifiques} : définir l'OGM concerné et rappeler brièvement le principe de sa construction (identification du gène d'intérêt, coupure par enzymes de restriction, ligature dans un vecteur, transformation de la cellule hôte, régénération).
\item \textbf{Présenter les avantages} en les illustrant d'exemples précis (insuline humaine, coton Bt au Burkina, riz doré) et de données chiffrées si l'énoncé en fournit (rendement, nombre de traitements, coût).
\item \textbf{Présenter les risques et limites} en les classant par catégorie (sanitaires, environnementaux, économiques, éthiques), toujours avec des mécanismes explicatifs (flux de gènes, pression de sélection, allergénicité, brevet).
\item \textbf{Conclure par une position justifiée et nuancée} : par exemple, un usage encadré au cas par cas, dans le respect du principe de précaution, de la réglementation de biosécurité (loi 2003) et des droits des paysans (semences paysannes), en distinguant OGM pharmaceutiques (confinés, consensus large) et OGM agricoles (disséminés, débat vif).
\end{enumerate}
\end{methode}

\begin{exoresolu}[Analyse de données et avis argumenté]
\textbf{Énoncé.} Dans une zone cotonnière du Nord-Cameroun, on compare deux parcelles sur trois ans : la parcelle A est semée en coton conventionnel (8 traitements insecticides par campagne contre les lépidoptères), la parcelle B en coton Bt (variété Bollgard II, 2 traitements par campagne contre les insectes piqueurs-suceurs non ciblés par Bt, rendement supérieur de 20 \%). Par ailleurs, des pièges à pollen et des analyses PCR sur des plants de sorgho sauvage (\emph{Sorghum halepense}) en bordure de la parcelle B révèlent la présence du gène \emph{cry1Ac} chez 3 \% des plants analysés. À partir de ces données, discute des avantages et des risques de cette culture OGM, puis dégage une position justifiée.
\tcblower
\textbf{Solution détaillée.}
\begin{enumerate}
\item \emph{Faits scientifiques} : Le coton Bollgard II est un OGM obtenu par transgenèse (\emph{Agrobacterium tumefaciens} ou biolistique). Il exprime deux toxines Cry (Cry1Ac, Cry2Ab) de \emph{Bacillus thuringiensis}, actives sur les lépidoptères ravageurs (ver de la capsule, pyrale).
\item \emph{Avantages} : La parcelle B ne reçoit que 2 traitements insecticides contre 8 pour la parcelle A, soit une réduction de 75 \% des pulvérisations. Le rendement est supérieur de 20 \%. Cela signifie : moindre coût intrants, moindre pollution eaux/sols, moindre exposition de l'agriculteur aux pesticides (santé au travail), gain de temps.
\item \emph{Risques} : La détection du gène \emph{cry1Ac} dans le sorgho sauvage voisin (3 \% des plants) prouve un \textbf{flux de gènes} effectif par hybridation interspécifique (le coton et le sorgho ne s'hybrident pas naturellement, mais cet exemple suppose une espèce proche compatible ; \textbf{correction pédagogique} : pour le coton, le flux de gènes vers le coton sauvage \emph{Gossypium} spp. est le risque réel. Pour le sorgho, c'est un OGM sorgho qui poserait ce problème. L'énoncé teste la compréhension du mécanisme). De plus, l'expression continue des toxines Bt exerce une forte \textbf{pression de sélection} sur les populations de ravageurs : risque d'apparition de résistances (déjà observées sur \emph{Helicoverpa armigera} en Inde/Chine/Australie).
\item \emph{Position justifiée} : Cette culture présente un intérêt agronomique, sanitaire et économique réel pour le paysan. Cependant, sa dissémination doit être strictement encadrée par la loi de biosécurité 2003 : mise en place obligatoire de \textbf{zones refuges} (coton non-Bt, 20 \% surface) pour diluer les allèles de résistance, surveillance du flux de gènes vers espèces apparentées, évaluation socio-économique du coût des semences, et information transparente des producteurs. L'application du principe de précaution impose cette gestion intégrée des risques.
\end{enumerate}
\end{exoresolu}

\begin{aretenir}[L'essentiel de la section]
\begin{itemize}
\item Les OGM offrent des avantages réels et mesurés : rendements accrus, réduction forte des pesticides, médicaments recombinants sûrs et accessibles, lutte contre les carences (biofortification).
\item Ils comportent des risques identifiés et documentés : allergénicité potentielle (testée), effets long terme (surveillés), \textbf{flux de gènes} vers espèces sauvages apparentées (sorgho, riz, colza), \textbf{résistance des ravageurs} (pression de sélection Bt), \textbf{érosion de la biodiversité cultivée}, \textbf{dépendance économique} (brevets, hybrides F1).
\item L'évaluation se fait \textbf{au cas par cas} (gène $\times$ plante $\times$ milieu $\times$ usage), selon le \cle{principe de précaution} : évaluer rigoureusement \emph{avant} de diffuser.
\item Le \cle{Protocole de Cartagena} (2000) encadre les échanges transfrontaliers d'OGM (APRC) ; le Cameroun l'a ratifié et dispose d'une \textbf{loi de biosécurité (2003)} opérationnelle.
\item Un avis argumenté se construit en quatre temps : faits scientifiques $\rightarrow$ avantages $\rightarrow$ risques/limites $\rightarrow$ position justifiée (encadrement, précaution, choix de société).
\end{itemize}
\end{aretenir}

\begin{exercice}[Pour t'entraîner]
Une coopérative de planteurs de la région du Nord (bénoué) hésite à adopter un \textbf{sorgho transgénique tolérant l'herbicide glyphosate} (variété \emph{StrigAway} type, résistance par mutation EPSPS). Rédige un texte de quinze à vingt lignes présentant aux planteurs deux avantages et deux risques majeurs de cette technologie, puis propose des conditions concrètes d'encadrement de son utilisation. Tu t'appuieras obligatoirement sur les notions de \textbf{flux de gènes} (vers sorgho sauvage \emph{S. halepense}), de \textbf{pression de sélection} (adventices résistantes au glyphosate) et de \textbf{principe de précaution}.
\end{exercice}

\newpage

\coursec{Activité d'intégration}

\coursec{Sensibilisation sur la technique du génie génétique dans le cadre de l'amélioration des caractéristiques des organismes}

\courssub{Objectifs de l'activité}

À l'issue de cette activité d'intégration, tu seras capable de :
\begin{itemize}
\item décrire, dans l'ordre logique, les six étapes d'une transgénèse bactérienne en nommant à chaque étape l'outil du génie génétique mobilisé (enzyme de restriction, ligase, plasmide vecteur) ;
\item justifier le rôle d'un gène marqueur de résistance à un antibiotique dans la sélection des bactéries transformées ;
\item expliquer l'intérêt concret d'un organisme génétiquement modifié (OGM) dans une situation de santé publique camerounaise ;
\item rédiger un texte argumenté présentant des avantages et des risques du génie génétique, avec une position personnelle justifiée.
\end{itemize}

\courssub{Situation}

L'hôpital de district de Nkongsamba (région du Littoral) prend en charge environ \num{1200} patients diabétiques de type 1, qui ont besoin d'injections quotidiennes d'\cle{insuline}. Or l'insuline importée coûte en moyenne \num{15000} FCFA le flacon, et les ruptures de stock sont fréquentes : en 2023, le service a enregistré 4 ruptures totalisant 87 jours sans insuline disponible. Chaque rupture expose les malades à des comas hyperglycémiques parfois mortels.

Face à cette situation, le club scientifique du lycée de la ville, encadré par son professeur de SVT, propose un projet ambitieux : produire localement de l'\cle{insuline humaine} en utilisant des bactéries \emph{Escherichia coli} rendues \cle{transgéniques}, c'est-à-dire ayant reçu le gène humain de l'insuline. C'est d'ailleurs exactement ainsi que l'insuline « humaine » est produite industriellement dans le monde depuis 1982.

Le professeur remet à l'équipe un dossier documentaire contenant les informations suivantes :

\textbf{Document 1 — Le gène de l'insuline.} Le gène humain codant l'insuline est un fragment d'ADN de 1431 paires de bases, porteur à ses deux extrémités de séquences reconnues par l'enzyme de restriction \emph{EcoR}I. Une fois exprimé dans une bactérie, il dirige la synthèse d'une protéine qui, après maturation, est strictement identique à l'insuline humaine.

\textbf{Document 2 — Le plasmide vecteur.} Le plasmide pNK1 est une molécule d'ADN circulaire de 5200 paires de bases. Il possède :
\begin{itemize}
\item un \cle{site unique} de coupure pour \emph{EcoR}I, situé à l'intérieur du gène de résistance à la tétracycline ;
\item un gène de résistance à l'ampicilline (gène \cle{marqueur}), intact ;
\item une origine de réplication lui permettant de se multiplier dans \emph{E. coli}.
\end{itemize}

\textbf{Document 3 — Rendement attendu.} Une souche bactérienne transformée, cultivée en fermenteur de 100 litres, peut produire environ 5 grammes d'insuline par cycle de culture, soit l'équivalent de 500 flacons. Le coût de revient estimé est de \num{3000} FCFA par flacon, cinq fois moins cher que l'insuline importée.

\textbf{Document 4 — Cadre réglementaire.} Le Cameroun est signataire du Protocole de Cartagena sur la prévention des risques biotechnologiques : toute production d'OGM doit être encadrée (confinement des souches, traçabilité, autorisation éthique). Le comité d'éthique de l'hôpital doit donner son avis avant tout lancement du projet.

L'équipe d'élèves-biotechnologues doit préparer le dossier technique et éthique à présenter au comité.

\courssub{Consignes}

\begin{enumerate}
\item À partir des documents 1 et 2, reproduis et complète le schéma des \textbf{six étapes} de la transgénèse : pour chaque étape, indique l'opération réalisée et nomme l'outil du génie génétique utilisé (enzyme de restriction \emph{EcoR}I, ADN ligase, plasmide pNK1).
\item Explique pourquoi la même enzyme de restriction (\emph{EcoR}I) doit être utilisée pour couper le gène de l'insuline et le plasmide.
\item Le gène de résistance à la tétracycline est détruit par l'insertion du gène de l'insuline, mais le gène de résistance à l'ampicilline reste fonctionnel. Justifie alors le choix d'un milieu de culture contenant de l'\textbf{ampicilline} pour sélectionner les bactéries transformées. Que deviennent les bactéries non transformées sur ce milieu ?
\item Calcule le nombre de cycles de culture nécessaires pour couvrir les besoins annuels de l'hôpital, sachant qu'un patient utilise en moyenne 2 flacons par mois, et compare le coût annuel de la production locale à celui de l'importation.
\item Rédige un texte argumenté d'environ 10 lignes destiné au comité d'éthique de l'hôpital : tu y présenteras \textbf{deux avantages} du procédé (sanitaires, économiques ou scientifiques) et \textbf{un risque} (sanitaire, environnemental ou éthique), puis tu concluras par une position justifiée sur l'opportunité du projet, en tenant compte du document 4.
\end{enumerate}

\courssub{Ressources à mobiliser}

\begin{aretenir}[Ce qu'il faut avoir en tête]
\begin{itemize}
\item Le \cle{génie génétique} est l'ensemble des techniques permettant de modifier le patrimoine génétique d'un organisme en lui transférant un ou plusieurs gènes ; l'organisme obtenu est un \cle{OGM} (ou organisme transgénique si le gène vient d'une autre espèce).
\item Les \cle{enzymes de restriction} coupent l'ADN à des sites spécifiques (véritables « ciseaux moléculaires ») ; la même enzyme produit des extrémités complémentaires (« cohésives ») sur le gène et sur le vecteur.
\item L'\cle{ADN ligase} soude le fragment d'ADN étranger au vecteur ouvert : c'est la « colle moléculaire ».
\item Le \cle{plasmide} est un vecteur de transfert : petit ADN circulaire bactérien capable de pénétrer dans une bactérie et de s'y répliquer.
\item Le \cle{gène marqueur} (résistance à un antibiotique) permet de sélectionner sur milieu contenant l'antibiotique les seules bactéries transformées.
\item Les six étapes d'une transgénèse : isolement du gène, coupure du vecteur, insertion (recombinaison), transformation de la bactérie, sélection des transformants, multiplication (clonage) et expression du gène.
\item Tout débat sur les OGM oppose des avantages (santé, rendements, coûts) à des risques (dissémination, allergènes, questions éthiques) et exige une argumentation fondée sur des faits, pas sur des peurs.
\end{itemize}
\end{aretenir}

\courssub{Démarche et corrigé commenté}

\begin{exoresolu}[Mission : produire de l'insuline pour l'hôpital de district]
\textbf{Énoncé rappelé.} Compléter le schéma des six étapes de la transgénèse en nommant les outils, justifier le choix du milieu sélectif à l'ampicilline, calculer les besoins et coûts annuels, puis rédiger un avis argumenté pour le comité d'éthique.

\tcblower

\textbf{Question 1 — Les six étapes de la manipulation.}

\begin{popfigure}
\begin{tikzpicture}[
    node distance=1.2cm and 0.5cm,
    box/.style={draw=popdark, thick, rounded corners, fill=popcream, minimum width=3.5cm, minimum height=1.2cm, align=center, font=\small},
    arrow/.style={-Stealth, thick, popdark},
    tool/.style={draw=popblue, thick, rounded corners, fill=popblueL, minimum width=3cm, minimum height=0.8cm, align=center, font=\small\bfseries, text=popblue!80!black},
    label/.style={font=\footnotesize, text=popmuted, align=center}
]
% Nodes for steps
\node[box, align=center] (s1){\textbf{1. Isolement du gène}\\\emph{Gène insuline humain (1431 pb)}};
\node[tool, below=of s1] (t1) {Outil : \emph{EcoR}I (enzyme de restriction)};

\node[box, right=3.5cm of s1, align=center] (s2){\textbf{2. Ouverture du vecteur}\\\emph{Plasmide pNK1 circulaire}};
\node[tool, below=of s2] (t2) {Outil : \emph{EcoR}I (enzyme de restriction)};

\node[box, right=3.5cm of s2, align=center] (s3){\textbf{3. Insertion / Ligation}\\\emph{Plasmide recombinant pNK1-insuline}};
\node[tool, below=of s3] (t3) {Outil : ADN ligase};

\node[box, below=2.5cm of s1, align=center] (s4){\textbf{4. Transformation}\\\emph{E. coli compétentes + plasmide recombinant}};
\node[tool, below=of s4] (t4) {Outil : Choc thermique / Électroporation};

\node[box, right=3.5cm of s4, align=center] (s5){\textbf{5. Sélection}\\\emph{Milieu + Ampicilline → Colonies transformées}};
\node[tool, below=of s5] (t5) {Outil : Gène marqueur \textit{amp}$^R$};

\node[box, right=3.5cm of s5, align=center] (s6){\textbf{6. Multiplication \& Expression}\\\emph{Fermenteur 100 L → Insuline humaine}};
\node[tool, below=of s6] (t6) {Outil : Origine de réplication (ori)};

% Arrows horizontal top
\draw[arrow] (s1.east) -- node[above, label] {Coupure \emph{EcoR}I} (s2.west);
\draw[arrow] (s2.east) -- node[above, label] {Mise en contact + Soudure} (s3.west);
% Arrow down from s3 to s4
\draw[arrow] (s3.south) -- node[left, label] {Introduction} (s4.north);
% Arrows horizontal bottom
\draw[arrow] (s4.east) -- node[above, label] {Ensemencement} (s5.west);
\draw[arrow] (s5.east) -- node[above, label] {Inoculation} (s6.west);

% Vertical connections steps to tools
\foreach \s/\t in {s1/t1, s2/t2, s3/t3, s4/t4, s5/t5, s6/t6}
    \draw[arrow, poporange] (\s.south) -- (\t.north);

% Sticky ends detail (visual hint)
\node[label, above=0.2cm of s1] {Extrémités cohésives 5'-AATT-3'};
\node[label, above=0.2cm of s2] {Extrémités cohésives 5'-AATT-3' (site dans \textit{tet}$^R$)};
\node[label, above=0.2cm of s3] {\textit{tet}$^R$ inactivé, \textit{amp}$^R$ intact};
\end{tikzpicture}
\legende{Schéma des six étapes de la transgénèse pour la production d'insuline humaine dans \emph{E. coli} via le plasmide pNK1. Les outils du génie génétique sont indiqués en bleu sous chaque étape.}
\end{popfigure}

\begin{enumerate}
\item \textbf{Isolement du gène d'intérêt} : le gène humain de l'insuline (1431 pb) est extrait du génome humain (ou synthétisé) et excisé à l'aide de l'\cle{enzyme de restriction} \emph{EcoR}I, qui le détache en laissant des extrémités cohésives (surplombs 5'-AATT-3').
\item \textbf{Ouverture du vecteur} : le plasmide pNK1 est coupé par la \emph{même} enzyme \emph{EcoR}I, au niveau de son site unique situé dans le gène de résistance à la tétracycline (\textit{tet}$^R$). Le plasmide circulaire devient linéaire, avec des extrémités cohésives complémentaires de celles du gène.
\item \textbf{Insertion du gène dans le vecteur (Recombinaison)} : le gène de l'insuline et le plasmide ouvert sont mis en contact ; leurs extrémités cohésives s'apparient par liaisons hydrogène, puis l'\cle{ADN ligase} catalyse la formation de liaisons phosphodiester, soude covalentement les deux molécules. On obtient un \cle{plasmide recombinant} (pNK1-insuline) où le gène \textit{tet}$^R$ est inactivé par insertion, tandis que le gène \textit{amp}$^R$ reste fonctionnel.
\item \textbf{Transformation} : le plasmide recombinant est introduit dans des bactéries \emph{E. coli} rendues compétentes (par traitement au chlorure de calcium et choc thermique, ou par électroporation), rendant leur membrane perméable à l'ADN.
\item \textbf{Sélection des bactéries transformées} : la culture est étalée sur un milieu gélosé contenant de l'ampicilline ; seules les bactéries porteuses du plasmide (donc du gène de résistance \textit{amp}$^R$ intact) survivent et forment des colonies visibles.
\item \textbf{Multiplication (clonage) et expression} : les colonies sélectionnées sont cultivées en fermenteur (100 L) ; les bactéries se multiplient (amplification du plasmide recombinant grâce à son \cle{origine de réplication}) et expriment le gène humain (transcription $\rightarrow$ traduction $\rightarrow$ maturation), produisant l'insuline, qui est ensuite extraite, purifiée et conditionnée.
\end{enumerate}

\begin{attention}[Piège fréquent : sélection vs criblage]
Le milieu à l'ampicilline sélectionne les \textbf{bactéries transformées} (porteuses du plasmide, recombinant ou non). Il ne distingue pas les plasmides recombinants (avec gène insuline) des plasmides non recombinants (recircularisés sans insert). Pour isoler les \textbf{vrais recombinants}, on pratique un \emph{criblage} par réplicature sur milieu à la tétracycline : les colonies recombinants sont \textit{tet}$^S$ (sensibles, ne poussent pas), les non-recombinants sont \textit{tet}$^R$ (résistantes, poussent). C'est le principe de l'\emph{inactivation d'un gène marqueur} (ici \textit{tet}$^R$).
\end{attention}

\textbf{Question 2 — Pourquoi la même enzyme de restriction ?} L'enzyme de restriction \emph{EcoR}I reconnaît une séquence palindromique précise (\texttt{5'-GAATTC-3'}) et coupe toujours entre G et A, produisant des extrémités cohésives (\texttt{5'-AATT-3'}). En utilisant \emph{EcoR}I pour le gène et pour le plasmide, on obtient des extrémités \emph{complémentaires} : le gène peut alors s'insérer dans le plasmide de façon stable et orientée (bien que l'orientation soit aléatoire avec une seule enzyme). Avec deux enzymes différentes, les extrémités ne seraient pas compatibles (sauf cas rare d'extrémités compatibles par hasard) et l'ADN ligase ne pourrait rien souder efficacement.

\textbf{Question 3 — Rôle du milieu à l'ampicilline.} L'insertion du gène de l'insuline détruit le gène de résistance à la tétracycline (\textit{tet}$^R$), mais le gène de résistance à l'ampicilline (\textit{amp}$^R$) reste intact dans le plasmide recombinant. Sur un milieu contenant de l'ampicilline :
\begin{itemize}
\item les bactéries \textbf{non transformées} (sans plasmide) ne possèdent aucune résistance : elles sont lysées / inhibées par l'antibiotique ;
\item les bactéries \textbf{transformées} (porteuses du plasmide recombinant ou non recombinant) expriment la $\beta$-lactamase codée par \textit{amp}$^R$ : elles hydrolysent l'ampicilline, survivent et se multiplient en colonies.
\end{itemize}
Le milieu à l'ampicilline agit donc comme un \cle{filtre sélectif} : il élimine l'immense majorité des bactéries non transformées et ne laisse vivre que celles qui ont acquis le plasmide (et donc potentiellement le gène de l'insuline). C'est le rôle du \cle{gène marqueur}.

\textbf{Question 4 — Calcul des besoins et des coûts.}
\begin{align*}
\text{Besoins annuels en flacons} &= \num{1200}\ \text{patients} \times 2\ \text{flacons/mois} \times 12\ \text{mois} = \num{28800}\ \text{flacons} \\
\text{Nombre de cycles de culture} &= \frac{\num{28800}\ \text{flacons}}{500\ \text{flacons/cycle}} = 57{,}6 \approx 58\ \text{cycles/an} \\
\text{Coût annuel production locale} &= \num{28800} \times \num{3000} = \num{86400000}\ \text{FCFA} \\
\text{Coût annuel importation} &= \num{28800} \times \num{15000} = \num{432000000}\ \text{FCFA} \\
\text{Économie réalisée} &= \num{432000000} - \num{86400000} = \num{345600000}\ \text{FCFA} \quad (\text{soit } 80\% \text{ de réduction})
\end{align*}

\textbf{Question 5 — Texte argumenté pour le comité d'éthique (modèle).}

\emph{Monsieur le Président, notre projet consiste à produire localement de l'insuline humaine grâce à des bactéries transgéniques. Ce procédé présente deux avantages majeurs. D'abord, un avantage sanitaire : il garantirait un approvisionnement permanent en insuline, alors que l'hôpital a subi 87 jours de rupture de stock en 2023, mettant en danger la vie de 1200 diabétiques de type 1. Ensuite, un avantage économique : le flacon produit localement coûterait 3000 FCFA contre 15000 FCFA importé, soit une économie de plus de 345 millions de FCFA par an, réinvestissable dans d'autres services de santé. Nous reconnaissons cependant un risque environnemental : la dissémination accidentelle de bactéries transgéniques dans l'environnement, qui pourrait favoriser le transfert horizontal du gène de résistance à l'ampicilline vers des bactéries pathogènes. Ce risque peut être maîtrisé par un confinement strict des souches en fermenteur de niveau de sécurité L2, leur destruction systématique par stérilisation (autoclavage 121°C, 20 min) après usage, et le respect rigoureux du Protocole de Cartagena signé par le Cameroun. En conclusion, nous estimons que les bénéfices sanitaires et économiques du projet l'emportent largement sur ses risques, à condition d'appliquer rigoureusement les mesures de biosécurité. Nous sollicitons donc l'avis favorable du comité.}

\end{exoresolu}

\courssub{Bilan de l'activité}

Cette activité t'a permis de mettre en évidence plusieurs acquis essentiels de la leçon :
\begin{itemize}
\item la \textbf{logique séquentielle} d'une transgénèse : chaque étape correspond à un outil précis (enzyme de restriction pour couper, ligase pour souder, plasmide pour transporter, gène marqueur pour sélectionner), et l'ordre des étapes n'est pas interchangeable ;
\item le \textbf{rôle stratégique du gène marqueur} : sans lui, impossible de distinguer les rares bactéries transformées (fréquence $\sim 10^{-5}$ à $10^{-8}$) parmi des millions de bactéries non transformées ;
\item la \textbf{portée concrète des biotechnologies} : le génie génétique n'est pas une science abstraite, il répond à des problèmes de santé publique réels, comme l'accès à l'insuline dans un hôpital de district camerounais ;
\item la \textbf{démarche d'argumentation scientifique} : face aux OGM, le citoyen éclairé pèse des avantages et des risques documentés, propose des mesures de maîtrise (confinement, réglementation, traçabilité) et conclut par une position justifiée, plutôt que de céder à l'enthousiasme aveugle ou à la peur irraisonnée.
\end{itemize}

Tu as ainsi mobilisé les trois savoir-faire attendus au programme : décrire une manipulation de génie génétique, expliquer l'intérêt d'un OGM et discuter des avantages et risques de ces biotechnologies — des compétences directement évaluables au Probatoire, série D.

\newpage

\coursec{Méthodes et exercices résolus}

\coursec{Génie génétique et OGM : définitions, outils et applications}

\courssub{Définitions et enjeux}

\begin{definition}[Génie génétique]
Ensemble de techniques permettant d'isoler, de modifier, de transférer et d'exprimer un ou plusieurs gènes d'un organisme \cle{donneur} vers un organisme \cle{receveur} (souvent d'une espèce différente), en utilisant des \cle{enzymes de restriction}, des \cle{ADN ligases} et des \cle{vecteurs de transfert}. C'est la base de la \cle{transgenèse}.
\end{definition}

\begin{definition}[Organisme génétiquement modifié (OGM)]
Organisme (bactérie, levure, plante, animal) dont le patrimoine génétique a été modifié par l'introduction, via la transgenèse, d'un ou plusieurs \cle{transgènes} (gènes étrangers) afin de lui conférer un nouveau caractère (production d'une protéine, résistance à un ravageur, tolérance à un herbicide, etc.). Un mutant obtenu au hasard ou un hybride issu de croisements classiques n'est \emph{pas} un OGM au sens de la transgenèse.
\end{definition}

\begin{definition}[ADN recombinant]
Molécule d'ADN chimérique obtenue \emph{in vitro} par ligation d'un fragment d'ADN d'intérêt (ex. gène humain) dans un vecteur (ex. plasmide), grâce à l'action conjuguée d'une enzyme de restriction et de l'ADN ligase.
\end{definition}

\begin{savaistu}[Cadre légal au Cameroun]
La loi n° 2003/006 du 21 avril 2003 fixe le régime de sécurité en biotechnologie moderne au Cameroun. Elle impose l'évaluation des risques pour la santé humaine, l'environnement et la biodiversité avant toute dissémination volontaire d'OGM (essais au champ, mise sur le marché). L'Agence Nationale de Biosécurité (ANB) est l'autorité compétente pour l'instruction des dossiers.
\end{savaistu}

\begin{popfigure}
\begin{tikzpicture}[
    node distance=1.2cm and 1.5cm,
    box/.style={draw=popdark, thick, rounded corners, fill=popcream, minimum width=2.5cm, minimum height=1cm, align=center, font=\small},
    arrow/.style={-Stealth, thick, popblue},
    plus/.style={font=\Large, poporange}
]
% Nodes
\node[box, align=center] (donor){ADN donneur\\(ex. Humain)\\\cle{Gène d'intérêt}};
\node[box, right=of donor, align=center] (vector){ADN vecteur\\(Plasmide)\\\cle{Origine réplication}\\\cle{Gène marqueur}};
\node[box, below=of donor, align=center] (cut){Coupure par la même\\\cle{enzyme de restriction}\\(ex. \emph{Eco}RI)};
\node[box, below=of vector] (cut2) {};
\node[box, below=of cut] (lig) {Ligation par l'\cle{ADN ligase}};
\node[box, below=of lig, align=center] (rec){\cle{ADN recombinant}\\(Plasmide recombinant)};
\node[box, below=of rec, align=center] (transfo){Transformation\\\cle{Transgenèse}\\(Cellule hôte compétente)};
\node[box, below=of transfo, align=center] (clone){Clonage \& Expression\\(Fermenteur / Régénération)};
\node[box, below=of clone, align=center] (ogm){\cle{OGM}\\\cle{Produit d'intérêt}\\(Insuline, Toxine Bt, etc.)};

% Arrows
\draw[arrow] (donor) -- node[left] {Extraction} (cut);
\draw[arrow] (vector) -- (cut2);
\draw[arrow] (cut) -- node[left, align=center] {Fragments à\\ extrémités cohésives} (lig);
\draw[arrow] (cut2) -- (lig);
\draw[arrow] (lig) -- node[left, align=center] {Appariement +\\ Souder} (rec);
\draw[arrow] (rec) -- node[left] {Introduction} (transfo);
\draw[arrow] (transfo) -- node[left] {Multiplication} (clone);
\draw[arrow] (clone) -- node[left] {Récolte / Purification} (ogm);

% Plus signs between cut and lig
\node[plus] at ($(cut.east)!0.5!(cut2.west)$) {+};
\node[plus] at ($(cut.south)!0.5!(lig.north)$) {+};
\end{tikzpicture}
\legende{Principes généraux de la transgenèse : de l'ADN donneur à l'OGM producteur. Les étapes clés sont la coupure spécifique (enzyme de restriction), l'appariement des extrémités cohésives, la ligation (ADN ligase), la transformation de la cellule hôte et l'expression du transgène.}
\end{popfigure}

\courssub{Les outils du génie génétique}

\begin{methode}[Reconnaître et justifier chaque outil]
\begin{enumerate}
\item \textbf{Enzyme de restriction (endonucléase)} : rôle de « ciseaux moléculaires ». Elle reconnaît une \cle{séquence spécifique} (souvent un palindrome, ex. \texttt{GAATTC} pour \emph{Eco}RI) et coupe l'ADN double brin à cet endroit, générant le plus souvent des \cle{extrémités cohésives} (bouts collants). \textbf{Justification cruciale} : le \emph{même} enzyme doit couper l'ADN donneur et le vecteur pour garantir la complémentarité des extrémités.
\item \textbf{ADN ligase} : rôle de « colle moléculaire ». Elle catalyse la formation des liaisons \cle{phosphodiester} entre le squelette sucre-phosphate du gène inséré et celui du vecteur, stabilisant l'ADN recombinant.
\item \textbf{Vecteur de transfert} : plasmide (ADN circulaire bactérien extrachromosomique) ou virus désactivé. Rôles : \cle{transporter} le gène d'intérêt dans la cellule hôte, assurer sa \cle{réplication} (origine de réplication \emph{ori}) et permettre la \cle{sélection} des cellules transformées grâce à un \cle{gène marqueur} (ex. résistance à l'ampicilline).
\item \textbf{Cellule hôte} : bactérie (\emph{E. coli}), levure (\emph{S. cerevisiae}), cellule végétale (via \emph{Agrobacterium} ou biolistique) ; elle réplique le vecteur et exprime le transgène (transcription + traduction).
\end{enumerate}
\textbf{Réflexe Probatoire} : dans un tableau ou un schéma, associe chaque outil à son rôle en une phrase complète : « L'enzyme de restriction \emph{Eco}RI permet de couper l'ADN donneur et le plasmide au même site pour générer des extrémités cohésives complémentaires. »
\end{methode}

\begin{popfigure}
\begin{tikzpicture}[
    scale=1.2,
    plasmid/.style={circle, draw=popblue, thick, minimum size=3.5cm},
    feature/.style={draw=popdark, thick, fill=poporangeL, rounded corners, font=\tiny, inner sep=2pt},
    label/.style={font=\small, align=center}
]
% Plasmid circle
\draw[plasmid] (0,0) circle (1.75cm);
% Features
\node[feature, label=above:\texttt{ori}] at (0,1.5) {Origine de réplication};
\node[feature, label=right:\texttt{amp}$^R$] at (1.5,0) {Résistance ampicilline};
\node[feature, label=below:\texttt{MCS}] at (0,-1.5) {Site de clonage multiple (MCS)};
\node[feature, label=left:\texttt{Insuline}] at (-1.5,0) {Gène insuline humaine};
% Restriction sites in MCS
\foreach \angle/\enz in {225/EcoRI, 255/BamHI, 285/HindIII, 315/SalI, 345/PstI, 15/SmaI, 45/KpnI, 75/XbaI} {
    \draw[popmuted, thick] (\angle:1.75) -- (\angle:1.4);
    \node[label, rotate=\angle+90] at (\angle:1.2) {\enz};
}
% Arrow indicating insertion
\draw[popgreen, -Stealth, thick] (-1.5, -2.5) -- (-1.5, -1.8) node[midway, right, font=\tiny, popgreen] {Insertion};
\node[label, popgreen] at (-1.5, -2.8) {Plasmide recombinant pINS};
\end{tikzpicture}
\legende{Carte simplifiée d'un plasmide recombinant (type pBR322/pUC) portant le gène de l'insuline humaine. Le MCS (Multiple Cloning Site) concentre les sites de restriction uniques. L'origine de réplication (\texttt{ori}) assure la multiplication du plasmide dans la bactérie ; le gène \texttt{amp}$^R$ permet la sélection sur milieu antibiotique.}
\end{popfigure}

\begin{exoresolu}[QCM justifié sur les outils]
Pour chaque affirmation, indiquer si elle est vraie ou fausse, en justifiant.
\begin{enumerate}
\item L'enzyme de restriction \emph{Eco}RI coupe n'importe quelle séquence d'ADN.
\item La ligase est indispensable à la formation de l'ADN recombinant.
\item Le plasmide sert uniquement à couper l'ADN.
\item On peut couper le gène d'intérêt avec \emph{Eco}RI et le plasmide avec \emph{Bam}HI, puis les assembler directement.
\end{enumerate}
\tcblower
\begin{enumerate}
\item \textbf{Fausse.} Une enzyme de restriction est \cle{spécifique} d'une séquence de reconnaissance : \emph{Eco}RI ne coupe l'ADN qu'au niveau de la séquence palindromique \texttt{GAATTC} (entre G et A). Cette spécificité garantit que le gène est excisé de façon précise et reproductible.
\item \textbf{Vraie.} Après appariement des extrémités cohésives par liaisons hydrogène, des cassures (nicks) subsistent dans le squelette sucre-phosphate ; seule l'\cle{ADN ligase} restaure les liaisons phosphodiester covalentes et stabilise l'ADN recombinant. Sans ligase, le gène inséré serait perdu lors de la réplication.
\item \textbf{Fausse.} Le plasmide est un \cle{vecteur de transfert} : il transporte le gène d'intérêt dans la cellule hôte et assure sa réplication autonome grâce à son origine de réplication (\texttt{ori}). Ce sont les enzymes de restriction qui coupent l'ADN.
\item \textbf{Fausse.} Deux enzymes de restriction différentes génèrent des extrémités cohésives de séquences différentes (ex. \texttt{AATT} pour \emph{Eco}RI vs \texttt{GATC} pour \emph{Bam}HI), donc \cle{non complémentaires} : l'appariement entre le gène et le vecteur est impossible. Il faut utiliser le \emph{même} enzyme (ou deux enzymes produisant des extrémités compatibles, ex. \emph{Sau}3AI et \emph{Bam}HI) pour les deux ADN.
\end{enumerate}
\textbf{Conclusion} : la réussite d'une transgenèse repose sur la complémentarité des extrémités cohésives, d'où l'usage d'un unique enzyme de restriction (ou d'enzymes compatibles), puis sur la soudure par la ligase et le transport par un vecteur adapté.
\end{exoresolu}

\courssub{Étapes de l'obtention d'un OGM}

\begin{methode}[Analyser un schéma de fabrication d'un OGM]
Face à un schéma de transgenèse (document fréquent au Probatoire), procède toujours dans le même ordre :
\begin{enumerate}
\item \textbf{Identifier les deux ADN de départ} : l'ADN \cle{donneur} (porteur du gène d'intérêt, ex. gène de l'insuline humaine) et l'ADN \cle{vecteur} (plasmide bactérien, virus désactivé...).
\item \textbf{Repérer l'enzyme de restriction} utilisée : vérifier que le \emph{même} enzyme coupe les deux ADN, afin d'obtenir des \cle{extrémités cohésives} (bouts collants) complémentaires.
\item \textbf{Nommer l'étape de recombinaison} : l'insertion du gène dans le vecteur est assurée par l'ADN \cle{ligase} ; le produit obtenu est un \cle{ADN recombinant} (plasmide recombinant).
\item \textbf{Identifier l'étape de transformation} : introduction du vecteur recombinant dans la cellule hôte (bactérie, levure, cellule végétale) — c'est la \cle{transgenèse}. Mentionner la sélection sur gène marqueur.
\item \textbf{Conclure} : la cellule hôte multipliée (clonage) exprime le gène d'intérêt ; on obtient un \cle{OGM} qui produit la protéine recherchée ou présente le nouveau caractère.
\end{enumerate}
\textbf{Réflexes de rédaction} : nomme précisément chaque enzyme (endonucléase de restriction, ligase), chaque structure (plasmide, gène d'intérêt, ADN recombinant) et justifie chaque étape par son rôle. Ne dis jamais « on mélange les ADN » mais « les extrémités cohésives complémentaires s'apparient par liaisons hydrogène puis sont soudées par la ligase ».
\end{methode}

\begin{exoresolu}[Fabrication d'une bactérie productrice d'insuline humaine]
Le document ci-dessous résume une manipulation réalisée en laboratoire : un fragment d'ADN humain portant le gène de l'insuline et un plasmide de la bactérie \emph{Escherichia coli} sont traités par l'enzyme de restriction \emph{Eco}RI ; les fragments sont mélangés en présence d'ADN ligase ; le plasmide obtenu est introduit dans des bactéries qui, cultivées en fermenteur, sécrètent ensuite une protéine active sur la glycémie des diabétiques.
\begin{enumerate}
\item Décrire, en quatre étapes ordonnées, la manipulation réalisée.
\item Nommer l'organisme obtenu et préciser la nature du produit récolté.
\end{enumerate}
\tcblower
\textbf{1. Description des étapes.}
\begin{itemize}
\item \textbf{Étape 1 — Préparation et coupure des ADN.} On extrait d'une part l'ADN humain porteur du \cle{gène de l'insuline} (gène d'intérêt), d'autre part le \cle{plasmide} d'\emph{E. coli} qui servira de vecteur. L'enzyme de restriction \emph{Eco}RI, endonucléase qui reconnaît la séquence spécifique \texttt{GAATTC}, coupe les deux ADN au même site : le gène d'intérêt est excisé et le plasmide est linéarisé, avec des \cle{extrémités cohésives} complémentaires dans les deux cas (séquence \texttt{AATT} en simple brin).
\item \textbf{Étape 2 — Formation de l'ADN recombinant.} Les extrémités cohésives du gène humain et du plasmide ouvert s'apparient par complémentarité des bases (A-T) ; l'\cle{ADN ligase} soude ensuite les squelettes sucre-phosphate (liaisons phosphodiester). On obtient un \cle{plasmide recombinant} porteur du gène de l'insuline humaine.
\item \textbf{Étape 3 — Transformation (transgenèse).} Le plasmide recombinant est introduit dans des bactéries \emph{E. coli} rendues compétentes (choc thermique ou électroporation) ; les bactéries ayant intégré le plasmide sont sélectionnées sur milieu contenant l'antibiotique correspondant au gène marqueur (ex. ampicilline).
\item \textbf{Étape 4 — Clonage et expression.} Les bactéries transformées se multiplient en fermenteur : chaque division réplique le plasmide (clonage du gène) et la machinerie bactérienne (ARN polymérase, ribosomes) traduit le gène humain en \cle{proinsuline}, puis en \cle{insuline} mature (après clivage peptidique), sécrétée dans le milieu de culture.
\end{itemize}
\textbf{2. Organisme obtenu et produit.} La bactérie obtenue est un \cle{organisme génétiquement modifié} (OGM), car son patrimoine génétique a reçu un gène étranger (humain) par transgenèse. Le produit récolté et purifié est l'\cle{insuline humaine} (identique à l'insuline endogène, 51 acides aminés, deux chaînes A et B liées par ponts disulfures), hormone hypoglycémiante utilisée pour traiter les diabétiques de type 1 et certains de type 2. Avant 1982 (mise sur le marché de l'Humulin\textsuperscript{\textregistered}), on utilisait l'insuline extraite du pancréas de porc ou de bœuf, qui diffère par quelques acides aminés et provoquait parfois des réactions immunitaires (anticorps anti-insuline) ou des lipodystrophies.

\textbf{Conclusion} : grâce à une enzyme de restriction, une ligase et un plasmide vecteur, on a transféré un gène humain dans une bactérie devenue « usine cellulaire » à insuline : c'est le principe même du génie génétique.
\end{exoresolu}

\courssub{Applications du génie génétique}

\begin{methode}[Construire l'explication de l'intérêt d'un OGM]
Pour répondre à « Explique l'intérêt de cet OGM », structure ta réponse en trois temps :
\begin{enumerate}
\item \textbf{Rappeler le problème initial} (besoin de santé, contrainte agricole) : ex. pénurie/allergénicité de l'insuline animale ; destruction des récoltes de coton par les chenilles au Nord-Cameroun ; carence en vitamine A dans les zones rizicoles.
\item \textbf{Décrire la modification apportée} : quel gène, prélevé chez quel organisme donneur, a été transféré dans quel organisme receveur, et quel nouveau caractère (phénotype) il confère.
\item \textbf{Dégager le bénéfice concret} : produit obtenu en grande quantité, pur, moins coûteux, sécurisé ; rendement agricole préservé ; réduction des traitements phytosanitaires (protection santé producteurs, environnement) ; valeur nutritionnelle améliorée.
\end{enumerate}
\textbf{Réflexe} : conclus toujours en reliant le bénéfice au caractère nouveau apporté par le transgène — c'est ce lien de cause à effet qui est évalué.
\end{methode}

\begin{exoresolu}[Le coton Bt au service des producteurs du Nord-Cameroun]
La chenille \emph{Helicoverpa armigera} détruit chaque année une partie importante des récoltes de coton dans la région du Nord-Cameroun, malgré de coûteux traitements insecticides. Des chercheurs ont transféré dans des plants de coton un gène de la bactérie \emph{Bacillus thuringiensis} (Bt) codant une protéine toxique pour les chenilles mais inoffensive pour l'Homme et les mammifères. Les plants obtenus se passent de la quasi-totalité des traitements insecticides.

Expliquer l'intérêt de cet organisme génétiquement modifié en suivant une démarche argumentée.
\tcblower
\textbf{1. Problème initial.} La chenille \emph{Helicoverpa armigera} ( lépidoptère) est un ravageur majeur du cotonnier ; les pertes de rendement (jusqu'à 50-80 \% sans traitement) et le coût des insecticides (achat, pulvérisations répétées, intoxications aiguës des producteurs, pollution des sols et des nappes phréatiques) pèsent lourdement sur les exploitations familiales du Nord-Cameroun.

\textbf{2. Modification apportée.} Le gène \cle{cry1Ac} (ou \cle{cry2Ab}) codant la \cle{protéine Cry} (delta-endotoxine Bt), prélevé chez la bactérie \emph{Bacillus thuringiensis} souche \emph{kurstaki}, a été inséré dans le génome du cotonnier (\emph{Gossypium hirsutum}) par transgenèse (vecteur \emph{Agrobacterium tumefaciens} ou bombardement de particules, puis régénération in vitro). Le cotonnier transgénique produit lui-même, dans ses tissus (feuilles, capsules, pollen), la protoxine Cry : c'est un \cle{OGM} (coton Bt).

\textbf{3. Bénéfices concrets.}
\begin{itemize}
\item \textbf{Agronomique} : les chenilles qui consomment les organes aériens ingèrent la protoxine, activée en toxine dans leur tube digestif alcalin (pH $\approx$ 10) ; elle se lie à des récepteurs spécifiques de la membrane des cellules intestinales, provoquant lyse et mort. Le rendement cotonnier est préservé (gain de 20 à 30 \% en conditions réelles).
\item \textbf{Économique} : suppression de la plupart des traitements insecticides chimiques (souvent 6 à 8 passages/an), donc baisse drastique des coûts de production et de la main-d'œuvre.
\item \textbf{Sanitaire et environnemental} : moins d'exposition des producteurs (souvent sans EPI) aux pesticides chimiques (pyréthrinoïdes, organophosphorés) neurotoxiques ; moindre pollution des eaux de surface ; la toxine Bt, spécifique de l'ordre des Lépidoptères (et quelques Coléoptères/Diptères selon la souche), est sans effet sur l'Homme, les mammifères, les oiseaux et la plupart des auxiliaires (abeilles, coccinelles).
\end{itemize}
\textbf{Conclusion} : l'intérêt du coton Bt est de conférer à la plante une \cle{résistance insecticide héréditaire et constitutive}, ce qui sécurise les récoltes tout en réduisant drastiquement l'usage d'insecticides chimiques — avantage à la fois économique, sanitaire et écologique.
\end{exoresolu}

\begin{methode}[Lire et interpréter des données de production (courbes, tableaux)]
\begin{enumerate}
\item \textbf{Lire les axes} : grandeur, unité, échelle (ex. quantité d'insuline en mg.L$^{-1}$ en fonction du temps en h).
\item \textbf{Identifier les conditions} : souche transformée (avec gène d'intérêt) \emph{vs} souche \cle{témoin} (sans gène, plasmide vide) — le témoin est indispensable pour attribuer l'effet au transgène.
\item \textbf{Dégager les faits chiffrés} : cite des valeurs précises (ex. « à 20 h, la souche transgénique produit \SI{45}{mg.L^{-1}} contre \SI{0}{mg.L^{-1}} pour le témoin »).
\item \textbf{Interpréter biologiquement} : relie l'écart observé à l'expression du gène transféré (transcription/traduction) et à la croissance de la population cellulaire ; conclus sur l'efficacité de la transgenèse et la faisabilité industrielle.
\end{enumerate}
\textbf{Réflexe} : une interprétation sans donnée chiffrée citée \emph{et} sans comparaison au témoin est une interprétation non étayée — elle perd des points.
\end{methode}

\begin{exoresolu}[Production d'insuline par deux souches bactériennes]
On cultive en fermenteurs identiques (même milieu, température, agitation, aération) deux souches d'\emph{E. coli} : la souche T (témoin, plasmide vide pUC18) et la souche M (plasmide recombinant pINS portant le gène de l'insuline humaine sous promoteur \emph{lac} ou \emph{tac}). On mesure la concentration en insuline du surnageant de culture :

\begin{center}
\begin{tabularx}{\linewidth}{|l|X|X|X|X|}
\hline
\rowcolor{popblueL} Temps de culture (h) & 0 & 10 & 20 & 30 \\
\hline
Souche T (mg.L$^{-1}$) & 0 & 0 & 0 & 0 \\
\hline
Souche M (mg.L$^{-1}$) & 0 & 15 & 45 & 60 \\
\hline
\end{tabularx}
\end{center}

\begin{enumerate}
\item Comparer l'évolution de la production d'insuline chez les deux souches.
\item Expliquer les résultats obtenus.
\end{enumerate}
\tcblower
\textbf{1. Comparaison.} Chez la souche témoin T, la concentration en insuline reste nulle (\SI{0}{mg.L^{-1}}) à tous les temps de culture. Chez la souche modifiée M, elle augmente régulièrement : \SI{15}{mg.L^{-1}} à 10 h, \SI{45}{mg.L^{-1}} à 20 h, puis \SI{60}{mg.L^{-1}} à 30 h, avec un ralentissement net en fin de culture (gain de \SI{15}{mg.L^{-1}} entre 20 h et 30 h contre \SI{30}{mg.L^{-1}} entre 10 h et 20 h).

\textbf{2. Explication.} La souche T ne possède pas le gène de l'insuline : \emph{E. coli} ne produit naturellement aucune insuline, d'où une concentration nulle — elle sert de \cle{témoin} prouvant que l'insuline mesurée chez M ne vient ni du milieu de culture ni du métabolisme bactérien normal (pas de fuite, pas de contamination). La souche M porte le \cle{plasmide recombinant} contenant le gène humain de l'insuline sous contrôle d'un promoteur fort ; lors de la croissance exponentielle (phase log), le gène est transcrit en ARNm puis traduit par les ribosomes bactériens en préproinsuline, traitée en insuline active, sécrétée dans le milieu. L'accumulation suit la cinétique de croissance bactérienne. Le ralentissement après 20 h correspond à l'entrée en phase stationnaire (épuisement du carbone/azote, accumulation de métabolites inhibiteurs, limitation en oxygène), limitant la synthèse protéique.

\textbf{Conclusion} : seule la souche ayant subi la transgenèse produit de l'insuline, ce qui démontre que le gène humain transféré est bien \cle{exprimé} (transcrit et traduit fonctionnellement) chez l'OGM bactérien ; la production industrielle peut donc reposer sur la culture à grande échelle (fermenteurs de \SI{10000}{L} et plus) de ces bactéries, suivie de purification chromatographique.
\end{exoresolu}

\courssub{Enjeux et débats autour des OGM}

\begin{methode}[Construire une argumentation équilibrée sur les OGM]
\begin{enumerate}
\item \textbf{Poser le cadre} : le génie génétique permet de transférer un caractère déterminé, rapidement et en franchissant la barrière des espèces (transgenèse), ce que la sélection classique (croisements intra-spécifiques) ne permet pas.
\item \textbf{Présenter les avantages} avec un exemple précis chacun :
      \begin{itemize}
      \item \textbf{Santé} : production de médicaments recombinants (insuline, hormone de croissance, facteurs de coagulation, vaccins à sous-unités comme anti-HBs) — pureté, innocuité, coût, disponibilité.
      \item \textbf{Agriculture} : résistance aux ravageurs (coton Bt, maïs Bt, aubergine Bt au Bangladesh), tolérance aux herbicides (soja RR, maïs RR — gestion simplifiée des adventices), tolérance aux stress abiotiques (sécheresse, salinité — ex. riz \emph{SUB1A} pour submersion), enrichissement nutritionnel (riz doré \emph{Golden Rice} : gènes \emph{psy} + \emph{crtI} $\rightarrow$ $\beta$-carotène/provitamine A dans l'endosperme).
      \item \textbf{Environnement} : réduction des pesticides (Bt), phytoremédiation (peupliers séquestrant mercure), bio-carburants (algues modifiées).
      \end{itemize}
\item \textbf{Présenter les risques ou craintes légitimes} avec le même soin :
      \begin{itemize}
      \item \textbf{Dissémination du transgène} (\cle{flux de gènes}) par le pollen vers des plantes sauvages apparentées ou des variétés conventionnelles voisines $\rightarrow$ adventices résistantes (« super-weeds »), contamination de filières « sans OGM ».
      \item \textbf{Apparition de ravageurs résistants} à la toxine Bt (pression de sélection forte) $\rightarrow$ nécessité de \cle{réfuges} (zones non-Bt) pour maintenir des allèles de sensibilité.
      \item \textbf{Risques allergéniques/toxicologiques} potentiels (nouvelle protéine) $\rightarrow$ évaluation cas par cas (comparaison séquence/allergènes connus, digestibilité, tests in vivo).
      \item \textbf{Enjeux socio-économiques} : brevetage du vivant, dépendance des paysans envers les semenciers (semences F1 non reproductibles, contrats), éthique (frontière des espèces).
      \end{itemize}
\item \textbf{Conclure de façon nuancée} : nécessité d'une évaluation \emph{au cas par cas}, de la \cle{biosécurité} (confinement, essais contrôlés, surveillance post-commercialisation), de la réglementation stricte et de l'information/participation du public — au Cameroun, la loi n° 2003/006 du 21 avril 2003 et son décret d'application 2005/0772/PM fixent ce cadre.
\end{enumerate}
\textbf{Réflexe Probatoire} : « Discuter » exige les deux points de vue \emph{et} une position argumentée ; un exposé à sens unique (tout pour ou tout contre) est sanctionné. Ne confonds pas génie génétique (transfert ciblé d'un gène connu) et sélection classique (brassage allélique sur tout le génome).
\end{methode}

\begin{exoresolu}[Sujet de discussion argumentée]
Dans le cadre d'un débat scientifique au lycée, un élève affirme : « Les OGM sont la solution à tous les problèmes agricoles et sanitaires, il faut les généraliser sans restriction. »

En t'appuyant sur tes connaissances et sur des exemples précis, discute cette affirmation (environ une vingtaine de lignes).
\tcblower
\textbf{Éléments de corrigé.}

L'affirmation de l'élève mérite d'être discutée car, si le génie génétique apporte des bénéfices réels et mesurables, il comporte aussi des risques et des limites qui imposent un encadrement rigoureux et une évaluation au cas par cas.

\textbf{Des avantages indéniables et documentés.} Sur le plan sanitaire, la production d'\cle{insuline humaine} par des bactéries transgéniques (\emph{E. coli} ou levure \emph{S. cerevisiae}) a mis fin aux rejets immunitaires liés à l'insuline animale (porcine/bovine) et garantit un approvisionnement illimité, standardisé et moins coûteux pour les millions de diabétiques, y compris dans les hôpitaux camerounais. D'autres protéines thérapeutiques (hormone de croissance, interférons, vaccins recombinants) suivent le même modèle. Sur le plan agricole, le \cle{coton Bt} (gène \emph{cry1Ac} de \emph{Bacillus thuringiensis}) résiste aux chenilles (\emph{Helicoverpa}) sans insecticides, préservant les rendements, la santé des producteurs du Nord-Cameroun et la biodiversité auxiliaire ; le \cle{riz doré} (Golden Rice), enrichi en provitamine A ($\beta$-carotène) via deux gènes (\emph{psy} du maïs, \emph{crtI} de bactérie), vise à lutter contre la cécité infantile liée à la carence en vitamine A en Asie. Le génie génétique agit en une génération et cible un caractère précis, là où la sélection classique demande 10 à 15 ans de croisements et de sélection.

\textbf{Mais des risques réels à maîtriser, pas à ignorer.} Le transgène peut se disséminer par le pollen vers des espèces sauvages voisines ou des cultures conventionnelles (\cle{flux de gènes}), créant des adventices résistantes aux herbicides (ex. colza tolérant au glyphosate devenu adventice) ou contaminant les filières « bio »/« sans OGM ». L'usage massif et continu du coton Bt ou maïs Bt exerce une pression de sélection intense qui favorise l'apparition de \cle{populations de ravageurs résistantes} à la toxine Cry (cas documentés pour \emph{Helicoverpa zea}, \emph{Busseola fusca}) ; la stratégie des \cle{réfuges} (zones non-Bt) est indispensable mais pas toujours respectée. Des risques \cle{allergéniques}, bien que rares et systématiquement évalués (comparaison bioinformatique, tests de digestibilité gastrique), ne peuvent être exclus \emph{a priori} pour toute nouvelle protéine. Enfin, le brevetage des semences OGM par quelques firmes multinationales peut créer une dépendance économique forte des petits producteurs (interdiction de ressemer, prix des semences, contrats léonins), posant un problème éthique de souveraineté alimentaire.

\textbf{Position argumentée.} Les OGM ne sont ni une solution universelle miracle ni un danger absolu : chaque OGM (événement de transformation précis : gène $\times$ variété $\times$ site d'insertion) doit être évalué \emph{au cas par cas}, avec des essais confinés, un plan de gestion de la résistance (réfuges), un contrôle de la dissémination (distance d'isolement, stérilité mâle) et une réglementation de biosécurité transparente (au Cameroun, loi n° 2003/006 du 21 avril 2003, décret 2005/0772/PM, Agence Nationale de Biosécurité). Généraliser « sans restriction », comme le propose l'élève, serait donc scientifiquement irresponsable (ignorance des risques écologiques évolutifs) et socialement inacceptable (mépris du principe de précaution et des choix de société).

\textbf{Conclusion} : le génie génétique est un outil puissant dont l'usage raisonné, réglementé, traçable et associé à l'information du public, permet de concilier progrès sanitaire, sécurité alimentaire, protection de l'environnement et équité socio-économique.
\end{exoresolu}

\courssub{Complément : Amélioration variétale classique et génie génétique}

\begin{methode}[Exploiter un croisement dans le cadre de l'amélioration des caractéristiques]
L'amélioration variétale classique (croisements dirigés + sélection) précède souvent ou complète le génie génétique ; le Probatoire demande de la distinguer (brassage allélique vs transgenèse) et de l'exploiter formellement.
\begin{enumerate}
\item \textbf{Définir les allèles} : note clairement l'allèle dominant (majuscule) et l'allèle récessif (minuscule), ex. $R$ = résistant (dominant), $s$ = sensible (récessif).
\item \textbf{Écrire les génotypes des parents} et les gamètes produits (un parent homozygote $R/\!/R$ ne produit que des gamètes $R$).
\item \textbf{Construire l'échiquier de croisement} (tableau à double entrée : gamètes d'un parent en ligne, de l'autre en colonne) et remplir chaque case avec le génotype et le phénotype.
\item \textbf{Dénombrer les phénotypes} et exprimer les proportions (fractions ou pourcentages).
\item \textbf{Conclure sur l'intérêt agronomique} : quelle génération faut-il sélectionner pour fixer le caractère utile ? (les individus homozygotes $R/\!/R$ sont stables à la descendance, contrairement aux hétérozygotes $R/\!/s$ qui ségrèguent).
\end{enumerate}
\textbf{Réflexe} : vérifie toujours que la somme des proportions vaut 1 (ou 100 \%). Utilise la notation $R/\!/s$ pour l'hétérozygote (barre de fraction pour séparer les allèles sur chromosomes homologues).
\end{methode}

\begin{exoresolu}[Croisement de plants de maïs pour l'amélioration variétale (Ouest-Cameroun)]
Chez le maïs (\emph{Zea mays}), l'allèle $R$ conférant la résistance à la striure (maladie virale transmise par la cicadelle \emph{Cicadulina mbila}, répandue dans l'Ouest-Cameroun) domine l'allèle $s$ de sensibilité. On croise un plant résistant de lignée pure avec un plant sensible, puis on croise entre eux les individus de la première génération (F$_1$).
\begin{enumerate}
\item Déterminer le génotype et le phénotype de la génération F$_1$.
\item À l'aide d'un échiquier de croisement, déterminer les proportions phénotypiques de la génération F$_2$.
\item Indiquer quels individus de F$_2$ le sélectionneur doit conserver pour obtenir une lignée stable de maïs résistant.
\end{enumerate}
\tcblower
\textbf{1. Génération F$_1$.} Le parent résistant de lignée pure est homozygote $R/\!/R$ (gamètes : 100 \% $R$) ; le parent sensible est homozygote $s/\!/s$ (gamètes : 100 \% $s$). Tous les individus F$_1$ reçoivent $R$ d'un parent et $s$ de l'autre : génotype $R/\!/s$ (hétérozygote). Comme $R$ domine $s$, le phénotype de toute la F$_1$ est \textbf{[résistant]}.

\textbf{2. Génération F$_2$ (croisement F$_1$ $\times$ F$_1$, soit $R/\!/s \times R/\!/s$).} Chaque parent F$_1$ produit deux types de gamètes en proportions égales (ségrégation méiotique) : $\dfrac{1}{2} R$ et $\dfrac{1}{2} s$.

\begin{center}
\begin{tabular}{|c|c|c|}
\hline
\rowcolor{popblueL} Gamètes $\downarrow$ / $\rightarrow$ & $R$ ($1/2$) & $s$ ($1/2$) \\
\hline
$R$ ($1/2$) & $R/\!/R$ \textbf{[résistant]} & $R/\!/s$ \textbf{[résistant]} \\
\hline
$s$ ($1/2$) & $R/\!/s$ \textbf{[résistant]} & $s/\!/s$ \textbf{[sensible]} \\
\hline
\end{tabular}
\end{center}

Proportions génotypiques : $\dfrac{1}{4} R/\!/R$ ; $\dfrac{2}{4} R/\!/s$ ; $\dfrac{1}{4} s/\!/s$.
Proportions \textbf{phénotypiques} : $\dfrac{3}{4}$ de plants résistants (75 \%) et $\dfrac{1}{4}$ de plants sensibles (25 \%). Vérification : $\dfrac{3}{4} + \dfrac{1}{4} = 1$.

\textbf{3. Choix du sélectionneur.} Parmi les plants résistants de F$_2$ (phénotype identique), seuls les homozygotes $R/\!/R$ (soit $\dfrac{1}{3}$ des résistants) produiront une descendance 100 \% résistante (lignée pure stable) ; les hétérozygotes $R/\!/s$ (2/3 des résistants) re-ségregueront 25 \% de sensibles à la génération suivante. Le sélectionneur conserve donc les plants $R/\!/R$, identifiés par des \cle{croisements-tests} (croisement avec un sensible $s/\!/s$ : un plant $R/\!/R$ ne donne que des descendants résistants ; un plant $R/\!/s$ donne 50 \% résistants / 50 \% sensibles).

\textbf{Conclusion} : l'amélioration variétale par croisements exploite les lois de la ségrégation (Mendel) pour fixer un caractère utile à l'état homozygote ; c'est un processus long (plusieurs générations d'autofécondation et sélection) limité au patrimoine génétique de l'espèce (ou espèces interfécondes). Le génie génétique, lui, aurait permis d'introduire directement un gène de résistance (ex. gène viral \emph{coat protein}, gène \emph{R} d'une espèce sauvage non croisable), en une seule transformation, franchissant la barrière des espèces.
\end{exoresolu}

\begin{attention}[Les 5 erreurs qui coûtent des points au Probatoire]
\begin{enumerate}
\item \textbf{Confondre enzyme de restriction et ligase.} La restriction \emph{coupe} l'ADN à une séquence spécifique ; la ligase \emph{soude} les fragments. Écrire « la ligase coupe le gène » ou « l'enzyme de restriction colle les ADN » annule le point de l'étape concernée.
\item \textbf{Oublier que le même enzyme de restriction doit couper les deux ADN.} Sans extrémités cohésives complémentaires, l'appariement et la ligation sont impossibles : c'est la justification la plus souvent attendue et la plus souvent omise dans les schémas.
\item \textbf{Définir l'OGM de façon vague.} « Un organisme modifié » ne suffit pas : il faut préciser qu'il a reçu, par \cle{transgenèse}, un ou plusieurs gènes \emph{étrangers} (issus d'une autre espèce, ou synthétiques) — un simple mutant spontané, un polyploïde ou un hybride de sélection classique n'est pas un OGM au sens réglementaire.
\item \textbf{Interpréter un tableau ou une courbe sans citer de valeurs chiffrées} et sans comparer à la souche \cle{témoin}. Sans témoin (plasmide vide), on ne peut pas attribuer la production observée au gène transféré (risque de protéine endogène ou contamination).
\item \textbf{Traiter le débat sur les OGM à sens unique.} « Discuter » impose de présenter avantages \emph{et} risques avec des exemples précis (insuline, coton Bt, riz doré, flux de gènes, résistance ravageurs, brevetage), puis une conclusion nuancée évoquant la biosécurité, la réglementation (loi camerounaise) et le principe de précaution. Ne confonds pas non plus génie génétique (transfert ciblé d'un gène connu) et sélection classique (croisements sur plusieurs générations, brassage de tout le génome).
\end{enumerate}
\end{attention}

\newpage

\coursec{Exercices}

\courssub{Je vérifie mes connaissances}

\begin{exercice}[QCM : les outils du génie génétique]
Pour chaque question, une seule réponse est exacte. Recopie la lettre de la bonne réponse.
\begin{enumerate}
\item Une enzyme de restriction est :
\begin{enumerate}
\item une enzyme qui colle deux fragments d'ADN ;
\item une enzyme qui coupe l'ADN à des sites spécifiques ;
\item une enzyme qui recopie l'ADN en ARN ;
\item une protéine qui transporte le gène dans la cellule.
\end{enumerate}
\item Le rôle de l'ADN ligase est de :
\begin{enumerate}
\item couper le plasmide ;
\item amplifier le gène d'intérêt ;
\item souder le gène d'intérêt au vecteur ;
\item détruire les bactéries non transformées.
\end{enumerate}
\item Un plasmide est :
\begin{enumerate}
\item un chromosome bactérien ;
\item une petite molécule d'ADN circulaire utilisée comme vecteur ;
\item une enzyme de restriction ;
\item une hormone protéique.
\end{enumerate}
\item Un OGM est un organisme :
\begin{enumerate}
\item issu d'une sélection naturelle ;
\item dont le patrimoine génétique a été modifié par transgénèse ;
\item qui a subi une mutation spontanée ;
\item stérilisé par irradiation.
\end{enumerate}
\end{enumerate}
\end{exercice}

\begin{exercice}[Vrai ou faux justifié]
Pour chacune des affirmations suivantes, indique si elle est vraie ou fausse, puis justifie ta réponse en une ou deux phrases.
\begin{enumerate}
\item Toutes les enzymes de restriction coupent l'ADN au même endroit.
\item Pour obtenir un ADN recombinant, il faut couper le gène d'intérêt et le vecteur avec la même enzyme de restriction.
\item Une bactérie transgénique peut produire une protéine humaine.
\item La transgénèse est une technique utilisée depuis l'Antiquité par les agriculteurs.
\end{enumerate}
\end{exercice}

\begin{exercice}[Définitions]
Définis en deux ou trois lignes chacun des termes suivants : \cle{génie génétique}, \cle{transgénèse}, \cle{vecteur de transfert}, \cle{ADN recombinant}.
\end{exercice}

\begin{exercice}[Sites de coupure]
On donne le site de reconnaissance de l'enzyme de restriction \emph{Eco}RI : elle reconnaît la séquence 5'-GAATTC-3' et coupe entre le G et le A.
\begin{enumerate}
\item Combien de sites de coupure cette enzyme possède-t-elle dans la séquence suivante ?

5'-\texttt{ATGC\textbf{GAATTC}CGTA\textbf{GAATTC}GGCA}-3'
\item Combien de fragments d'ADN obtient-on après digestion complète de cette molécule \emph{linéaire} ?
\end{enumerate}
\end{exercice}

\courssub{Je m'entraîne}

\begin{exercice}[Remettre les étapes dans l'ordre]
Voici, dans le désordre, les étapes de l'obtention d'une bactérie produisant de l'insuline humaine :
\begin{enumerate}
\item[a.] introduction du plasmide recombinant dans la bactérie ;
\item[b.] extraction du gène de l'insuline à partir de cellules humaines (obtention de l'ADNc par transcription inverse) ;
\item[c.] culture des bactéries transformées et récolte de l'insuline ;
\item[d.] soudure du gène au plasmide ouvert grâce à une ligase ;
\item[e.] ouverture du plasmide par une enzyme de restriction ;
\item[f.] sélection des bactéries ayant intégré le plasmide.
\end{enumerate}
Remets ces étapes dans l'ordre chronologique et nomme, pour chacune, l'enzyme ou l'outil utilisé lorsqu'il y en a un.
\end{exercice}

\begin{exercice}[La même enzyme pour les deux molécules]
On donne deux séquences d'ADN :

Gène d'intérêt : 5'-\texttt{CCG\textbf{GAATTC}TTG}-3'

Plasmide : 5'-\texttt{AAT\textbf{GAATTC}CCG}-3'
\begin{enumerate}
\item Repère le site de coupure de l'enzyme \emph{Eco}RI (séquence GAATTC, coupure entre G et A) dans chaque molécule.
\item Explique pourquoi il est indispensable d'utiliser la \emph{même} enzyme de restriction pour couper le gène d'intérêt et le plasmide. Que se passerait-il si l'on utilisait deux enzymes différentes ?
\item Nomme l'enzyme qui permet ensuite de réunir les deux fragments.
\end{enumerate}
\end{exercice}

\begin{exercice}[Insuline bactérienne et insuline porcine]
Avant 1982, les diabétiques étaient traités avec de l'insuline extraite du pancréas de porc ou de bœuf. Aujourd'hui, l'insuline est produite par des bactéries transgéniques.
\begin{enumerate}
\item Rappelle le principe de la production d'insuline humaine par des bactéries.
\item Explique pourquoi l'insuline produite par les bactéries transgéniques est mieux tolérée par les diabétiques que l'insuline extraite du pancréas de porc (différence d'un acide aminé en position B30).
\item Cite un autre avantage de cette production pour la société (quantité, coût, disponibilité, risque infectieux).
\end{enumerate}
\end{exercice}

\begin{exercice}[Analyse d'un schéma de transgénèse]
On te donne un schéma de transgénèse comportant les étiquettes suivantes, placées de façon désordonnée : « plasmide recombinant », « gène d'intérêt », « enzyme de restriction », « ligase », « bactérie transformée », « protéine d'intérêt ».
\begin{enumerate}
\item Replace chaque étiquette à l'étape qui lui correspond en décrivant le schéma corrigé sous forme d'un texte ordonné.
\item Pour chaque étape, précise s'il s'agit d'une étape de coupure, de soudure, de transfert ou d'expression du gène.
\end{enumerate}
\end{exercice}

\courssub{Je me prépare au Probatoire}

\begin{exercice}[Le maïs Bt en Afrique du Sud]
\textbf{Document.} Depuis 1998, l'Afrique du Sud cultive du maïs Bt. Ce maïs a reçu, par transgénèse, un gène de la bactérie \emph{Bacillus thuringiensis} (Bt). Ce gène code une protéine toxique pour les chenilles de la pyrale, insecte ravageur qui détruit une partie importante des récoltes de maïs. Les agriculteurs qui cultivent le maïs Bt utilisent moins d'insecticides chimiques et obtiennent des rendements plus élevés. Cependant, certains scientifiques craignent l'apparition d'insectes résistants à la toxine Bt et la dissémination du transgène vers les variétés locales par pollinisation croisée.
\begin{enumerate}
\item Dégage le problème agricole posé dans ce document.
\item Indique le gène transféré, l'organisme donneur, l'organisme receveur et le caractère nouveau obtenu.
\item Décris, en quatre étapes numérotées, comment on a pu obtenir ce maïs transgénique (outils enzymatiques et vecteur attendus).
\item Relève dans le document deux avantages du maïs Bt pour l'agriculteur.
\item Relève deux risques évoqués et propose une mesure permettant de limiter chacun d'eux.
\end{enumerate}
\end{exercice}

\begin{exercice}[Production d'insuline à l'échelle industrielle]
Une entreprise pharmaceutique souhaite produire de l'insuline humaine en grande quantité à l'aide de la bactérie \emph{Escherichia coli}. On dispose de l'ADNc du gène humain de l'insuline (sans introns) et d'un plasmide portant un gène de résistance à un antibiotique. Les deux molécules possèdent un site de coupure pour la même enzyme de restriction.
\begin{enumerate}
\item Définis les termes : enzyme de restriction, vecteur, organisme génétiquement modifié.
\item Décris, sous forme d'un protocole en cinq étapes numérotées, la fabrication de la souche bactérienne productrice d'insuline, en nommant à chaque étape l'outil utilisé.
\item Explique le rôle du gène de résistance à l'antibiotique porté par le plasmide : comment permet-il de sélectionner les bactéries transformées ?
\item La bactérie produit une insuline identique à l'insuline humaine. Justifie cette identité en rappelant une propriété fondamentale du code génétique.
\item Cite deux avantages de cette technique par rapport à l'extraction d'insuline à partir de pancréas d'animaux d'élevage.
\end{enumerate}
\end{exercice}

\begin{exercice}[Débat argumenté : les OGM au Cameroun]
Le Cameroun fait face à des défis agricoles importants : pertes dues aux ravageurs, sécheresses, malnutrition. Certains proposent d'introduire des cultures OGM (maïs résistant aux insectes, manioc enrichi en vitamines) ; d'autres s'y opposent au nom du principe de précaution.
\begin{enumerate}
\item Rappelle la définition d'un OGM et le principe général de sa fabrication.
\item Présente, en un paragraphe structuré, deux avantages que pourraient apporter les cultures OGM à l'agriculture camerounaise.
\item Présente, en un second paragraphe, deux risques (sanitaires ou environnementaux) associés à ces cultures.
\item Rédige une conclusion argumentée d'une dizaine de lignes dans laquelle tu défends une position personnelle (favorable, opposée ou conditionnelle) en t'appuyant sur le principe de précaution et sur les connaissances scientifiques du cours.
\end{enumerate}
\end{exercice}

\courssub{Corrigés des exercices du Probatoire}



\newpage

\coursec{Corrigés des exercices}

\begin{corrige}[Exercice 1 — QCM : les outils du génie génétique]
\begin{enumerate}
\item \textbf{Réponse b.} Une enzyme de restriction est une endonucléase qui reconnaît une séquence spécifique de l'ADN (site de reconnaissance, souvent palindromique, par exemple GAATTC pour \emph{Eco}RI) et \textbf{coupe la molécule d'ADN} au niveau de ce site. C'est le « ciseau moléculaire » du génie génétique. Les réponses a (ligase), c (transcriptase) et d (vecteur) correspondent à d'autres acteurs de la manipulation.
\item \textbf{Réponse c.} L'\cle{ADN ligase} catalyse la formation de liaisons phosphodiesters : elle \textbf{soude} le gène d'intérêt au vecteur ouvert, formant ainsi l'ADN recombinant. C'est la « colle moléculaire ».
\item \textbf{Réponse b.} Un plasmide est une \textbf{petite molécule d'ADN circulaire double brin}, extra-chromosomique, présente naturellement chez les bactéries et capable de se répliquer de façon autonome. On l'utilise comme \cle{vecteur de transfert} pour introduire un gène étranger dans une cellule hôte.
\item \textbf{Réponse b.} Un OGM est un organisme dont le \textbf{patrimoine génétique a été modifié artificiellement par transgénèse} (insertion d'un ou plusieurs gènes étrangers), modification qui ne peut pas s'obtenir naturellement par croisement. La sélection naturelle (a) et la mutation spontanée (c) ne relèvent pas du génie génétique.
\end{enumerate}
\begin{attention}[Points de vigilance]
Ne pas confondre enzyme de restriction (coupe) et ligase (soude) : c'est l'inversion la plus fréquente au Probatoire. De même, le plasmide n'est \emph{pas} le chromosome bactérien : c'est une molécule d'ADN \emph{supplémentaire}, circulaire et autonome.
\end{attention}
\end{corrige}

\begin{corrige}[Exercice 2 — Vrai ou faux justifié]
\begin{enumerate}
\item \textbf{FAUX.} Chaque enzyme de restriction reconnaît une \textbf{séquence spécifique} qui lui est propre : \emph{Eco}RI reconnaît GAATTC, \emph{Hin}dIII reconnaît AAGCTT, \emph{Bam}HI reconnaît GGATCC, etc. Il existe des centaines d'enzymes de restriction, chacune avec son propre site de coupure.
\item \textbf{VRAI.} La même enzyme de restriction génère des \textbf{extrémités complémentaires} (bouts cohésifs ou « bouts collants ») identiques sur le gène d'intérêt et sur le vecteur. Ces extrémités complémentaires s'apparient par complémentarité des bases, ce qui permet ensuite à la ligase de souder les deux molécules. Avec deux enzymes différentes, les extrémités seraient incompatibles et la soudure impossible.
\item \textbf{VRAI.} Grâce à l'\textbf{universalité du code génétique}, un gène humain inséré dans une bactérie (via un plasmide recombinant) est transcrit et traduit par la machinerie bactérienne en une protéine identique à la protéine humaine. C'est le principe de la production d'insuline humaine par \emph{Escherichia coli} depuis 1982.
\item \textbf{FAUX.} Les agriculteurs pratiquent depuis l'Antiquité la \textbf{sélection} (choix des meilleurs individus reproducteurs) et les \textbf{croisements}, qui sont des techniques de génétique classique, sans modification directe de l'ADN. La \cle{transgénèse} (transfert d'un gène d'une espèce à une autre par génie génétique) n'existe que depuis les années 1970 (premières bactéries recombinantes : Cohen et Boyer, 1973).
\end{enumerate}
\end{corrige}

\begin{corrige}[Exercice 3 — Définitions]
\begin{itemize}
\item \cle{Génie génétique} : ensemble des techniques de manipulation \emph{in vitro} de l'ADN (coupure, soudure, transfert) permettant d'isoler un gène, de le modifier ou de le transférer d'un organisme à un autre, afin de conférer à l'organisme receveur un caractère nouveau.
\item \cle{Transgénèse} : technique du génie génétique consistant à introduire dans le génome d'un organisme receveur un ou plusieurs gènes étrangers (\cle{transgènes}), issus d'un organisme donneur, de sorte que ce gène soit exprimé et transmis à la descendance. L'organisme obtenu est dit « transgénique ».
\item \cle{Vecteur de transfert} : molécule d'ADN (le plus souvent un plasmide bactérien, parfois un virus) capable de pénétrer dans une cellule hôte et de s'y répliquer, utilisée comme « véhicule » pour transporter le gène d'intérêt jusqu'à l'intérieur de la cellule receveuse.
\item \cle{ADN recombinant} : molécule d'ADN artificielle obtenue par la \textbf{soudure} (grâce à une ligase) de fragments d'ADN d'origines différentes, par exemple un gène humain inséré dans un plasmide bactérien. C'est la molécule clé de toute manipulation de génie génétique.
\end{itemize}
\end{corrige}

\begin{corrige}[Exercice 4 — Sites de coupure]
\begin{enumerate}
\item La séquence proposée est : 5'-ATGC\textbf{GAATTC}CGTA\textbf{GAATTC}GGCA-3'. On repère la séquence de reconnaissance GAATTC : elle apparaît \textbf{deux fois} (en gras dans l'énoncé). L'enzyme \emph{Eco}RI possède donc \textbf{2 sites de coupure} dans cette molécule.
\item Une molécule d'ADN \textbf{linéaire} coupée en $n$ points donne $n+1$ fragments. Ici, $n = 2$ coupures, donc on obtient \textbf{3 fragments} :
\begin{center}
5'-ATGCG $\vert$ AATTCCGTAG $\vert$ AATTCGGCA-3'
\end{center}
soit un fragment de 5 nucléotides (ATGCG), un fragment intermédiaire (AATTCCGTAG) et un fragment terminal (AATTCGGCA).
\end{enumerate}
\begin{popfigure}
\begin{tikzpicture}[scale=0.9, every node/.style={font=\small}]
% DNA backbone
\draw[popdark, thick] (0,0) -- (12,0) node[midway, below=4pt] {ADN linéaire (24 pb)};
% Ticks and labels
\foreach \x/\label in {0/1, 5/6, 11/12, 12/13, 17/18, 23/24} {
    \draw[popdark] (\x*0.5, 0.1) -- (\x*0.5, -0.1) node[below=2pt] {\label};
}
% EcoRI Sites
% Site 1: between 5 and 6 (index 5-6) -> position 5.5 bp -> x=2.75
% Site 2: between 17 and 18 (index 17-18) -> position 17.5 bp -> x=8.75
\draw[poporange, very thick] (2.75, 0.3) -- (2.75, 1.2) node[above] {\emph{Eco}RI \#1};
\draw[poporange, very thick] (8.75, 0.3) -- (8.75, 1.2) node[above] {\emph{Eco}RI \#2};
% Recognition sequence boxes
\node[draw=poporange, fill=poporangeL, rounded corners, inner sep=2pt] at (2.75, 1.6) {GAATTC};
\node[draw=poporange, fill=poporangeL, rounded corners, inner sep=2pt] at (8.75, 1.6) {GAATTC};
% Fragments brackets
\draw[popblue, decorate, decoration={brace, amplitude=5pt}, thick] (0, -0.5) -- (2.75, -0.5) node[midway, below=6pt] {Fragment 1 (5 pb)};
\draw[popblue, decorate, decoration={brace, amplitude=5pt}, thick] (2.75, -0.5) -- (8.75, -0.5) node[midway, below=6pt] {Fragment 2 (10 pb)};
\draw[popblue, decorate, decoration={brace, amplitude=5pt}, thick] (8.75, -0.5) -- (12, -0.5) node[midway, below=6pt] {Fragment 3 (9 pb)};
\end{tikzpicture}
\legende{Carte de restriction d'un ADN linéaire de 24 paires de bases par l'enzyme \emph{Eco}RI. Deux sites de reconnaissance (GAATTC) génèrent trois fragments après digestion.}
\end{popfigure}
\begin{attention}[Point de vigilance]
La règle « $n$ coupures $\rightarrow$ $n+1$ fragments » ne vaut que pour une molécule \textbf{linéaire}. Pour un ADN \textbf{circulaire} (comme un plasmide), $n$ coupures donnent $n$ fragments. Cette distinction est un piège classique du Probatoire.
\end{attention}
\end{corrige}

\begin{corrige}[Exercice 5 — Remettre les étapes dans l'ordre]
Ordre chronologique correct : \textbf{b $\rightarrow$ e $\rightarrow$ d $\rightarrow$ a $\rightarrow$ f $\rightarrow$ c}.
\begin{enumerate}
\item \textbf{b. Extraction du gène de l'insuline} à partir de cellules humaines ; comme le gène humain contient des introns que la bactérie ne sait pas éliminer, on obtient en pratique l'\textbf{ADNc} (ADN complémentaire, sans introns) par \textbf{transcription inverse} de l'ARNm, grâce à l'enzyme \textbf{transcriptase inverse}.
\item \textbf{e. Ouverture du plasmide} : le vecteur (plasmide) est coupé par une \textbf{enzyme de restriction} (la même que celle utilisée pour préparer le gène).
\item \textbf{d. Soudure} du gène au plasmide ouvert par l'\textbf{ADN ligase} $\rightarrow$ formation du \textbf{plasmide recombinant} (ADN recombinant).
\item \textbf{a. Introduction du plasmide recombinant dans la bactérie} : c'est l'étape de \textbf{transformation} (transfert), réalisée par choc thermique en présence de CaCl$_2$ ou par électroporation.
\item \textbf{f. Sélection des bactéries transformées} : culture sur milieu contenant un \textbf{antibiotique} ; seules les bactéries ayant intégré le plasmide (porteur d'un gène de résistance à cet antibiotique, le \textbf{marqueur de sélection}) survivent.
\item \textbf{c. Culture des bactéries transformées et récolte de l'insuline} : multiplication en fermenteurs, expression du gène, puis extraction et purification de l'insuline humaine.
\end{enumerate}
\begin{methode}[Astuce de mémorisation]
Retiens l'ordre logique en 4 verbes : \textbf{Couper} (enzyme de restriction) $\rightarrow$ \textbf{Coller} (ligase) $\rightarrow$ \textbf{Transférer} (transformation) $\rightarrow$ \textbf{Sélectionner/Cultiver} (marqueur de sélection, fermenteurs). L'obtention du gène (ADNc) précède toujours la coupure du vecteur.
\end{methode}
\end{corrige}

\begin{corrige}[Exercice 6 — La même enzyme pour les deux molécules]
\begin{enumerate}
\item Site de coupure d'\emph{Eco}RI (coupure entre G et A de GAATTC) :
\begin{itemize}
\item Gène d'intérêt : 5'-CCG$\vert$\textbf{G$\vert$AATTC}TTG-3' $\rightarrow$ coupure entre le G et le A du site GAATTC, soit un fragment 5'-CCGG-3' et un fragment 5'-AATTCTTG-3'.
\item Plasmide : 5'-AAT\textbf{G$\vert$AATTC}CCG-3' $\rightarrow$ coupure entre le G et le A du site GAATTC, soit 5'-AATG-3' et 5'-AATTCCCG-3'.
\end{itemize}
Dans les deux cas, la coupure génère l'extrémité cohésive \textbf{5'-AATT-3'} (bout collant).
\begin{popfigure}
\begin{tikzpicture}[scale=0.85, every node/.style={font=\small, font=\ttfamily}]
% Gene fragment top
\node[anchor=east] at (0, 1.5) {Gène :};
\node[draw, fill=popblueL, rounded corners, inner sep=2pt] at (2.5, 1.5) {5'-CCGG};
\node[draw, fill=poporangeL, rounded corners, inner sep=2pt] at (4.5, 1.5) {AATT};
\node[draw, fill=popblueL, rounded corners, inner sep=2pt] at (6.5, 1.5) {CTTG-3'};
% Plasmid fragment bottom
\node[anchor=east] at (0, 0) {Plasmide :};
\node[draw, fill=popgreenL, rounded corners, inner sep=2pt] at (2.5, 0) {5'-AATG};
\node[draw, fill=poporangeL, rounded corners, inner sep=2pt] at (4.5, 0) {AATT};
\node[draw, fill=popgreenL, rounded corners, inner sep=2pt] at (6.5, 0) {CCCG-3'};
% Complementary bases
\draw[popdark, dashed] (3.5, 1.3) -- (3.5, 0.7);
\draw[popdark, dashed] (5.5, 1.3) -- (5.5, 0.7);
\node[popdark] at (4.5, 1.9) {Extrémités cohésives complémentaires (bouts collants)};
\node[popdark] at (4.5, -0.5) {Appariement par ponts hydrogène (A=T)};
% Ligase action
\draw[->, poppurple, thick] (7.5, 0.75) -- (9, 0.75) node[midway, above] {ADN Ligase};
\node[draw, fill=poppurpleL, rounded corners, inner sep=4pt, align=center] at (10.5, 0.75) {ADN Recombinant\\5'-CCGG-AATT-CTTG-3'\\3'-GGCC-TTAA-GAAC-5'};
\end{tikzpicture}
\legende{Visualisation de la coupure par \emph{Eco}RI générant des bouts collants complémentaires (AATT) sur le gène d'intérêt et le plasmide, permettant leur assemblage par l'ADN ligase.}
\end{popfigure}
\item Il est indispensable d'utiliser la \textbf{même enzyme de restriction} car elle produit sur les deux molécules des \textbf{extrémités cohésives complémentaires} (ici AATT), capables de s'apparier par complémentarité des bases (A avec T). Cet appariement positionne correctement le gène dans le vecteur avant la soudure. Si l'on utilisait \textbf{deux enzymes différentes}, les extrémités produites seraient de séquences différentes, \textbf{non complémentaires} : le gène et le vecteur ne pourraient pas s'apparier, et la formation de l'ADN recombinant serait impossible.
\item L'enzyme qui réunit (soude) ensuite les deux fragments est l'\textbf{ADN ligase}, qui restaure les liaisons phosphodiesters du squelette sucre-phosphate.
\end{enumerate}
\end{corrige}

\begin{corrige}[Exercice 7 — Insuline bactérienne et insuline porcine]
\begin{enumerate}
\item \textbf{Principe :} on isole le gène humain de l'insuline (en pratique son ADNc obtenu par transcription inverse de l'ARNm), on l'insère dans un plasmide grâce à une enzyme de restriction puis une ligase (ADN recombinant), on introduit ce plasmide dans la bactérie \emph{Escherichia coli} (transformation), on sélectionne les bactéries transformées et on les cultive en fermenteurs. Les bactéries expriment le gène humain et produisent de l'insuline, que l'on récolte et purifie.
\item \textbf{Meilleure tolérance :} l'insuline de porc diffère de l'insuline humaine par \textbf{un acide aminé} (position B30 : \textbf{thréonine chez l'homme}, \textbf{alanine chez le porc}). Cette différence de structure primaire peut être reconnue comme étrangère par le système immunitaire du patient, provoquant la production d'\textbf{anticorps anti-insuline} (réactions immunitaires, allergies, diminution d'efficacité du traitement). L'insuline produite par les bactéries transgéniques possède une séquence d'acides aminés \textbf{strictement identique} à l'insuline humaine (universalité du code génétique) : elle n'est pas reconnue comme étrangère et est donc \textbf{parfaitement tolérée}.
\item \textbf{Autres avantages pour la société} (un seul demandé) :
\begin{itemize}
\item production en \textbf{grande quantité} et continue (fermenteurs), indépendante des abattages d'animaux $\rightarrow$ disponibilité assurée pour tous les diabétiques ;
\item \textbf{coût réduit} à grande échelle ;
\item absence de \textbf{risque de contamination} par des agents infectieux d'origine animale (virus, prions) présents dans les pancréas.
\end{itemize}
\end{enumerate}
\end{corrige}

\begin{corrige}[Exercice 8 — Analyse d'un schéma de transgénèse]
\begin{enumerate}
\item \textbf{Schéma corrigé décrit sous forme de texte ordonné :}
\begin{enumerate}
\item L'\textbf{enzyme de restriction} coupe l'ADN du donneur pour libérer le \textbf{gène d'intérêt}, et ouvre le plasmide au même site.
\item La \textbf{ligase} soude le gène d'intérêt dans le plasmide ouvert : on obtient le \textbf{plasmide recombinant}.
\item Le plasmide recombinant est introduit dans une bactérie : on obtient une \textbf{bactérie transformée}.
\item La bactérie transformée, cultivée, exprime le gène et fabrique la \textbf{protéine d'intérêt}.
\end{enumerate}
\begin{popfigure}
\begin{tikzpicture}[node distance=1.2cm and 1.5cm, >=Stealth, font=\small,
    box/.style={draw, thick, rounded corners, minimum height=1.2cm, minimum width=2.8cm, align=center, fill=popcream},
    arrow/.style={->, thick, popdark}
]
\node[box, fill=popblueL, align=center] (start){ADN Donneur\\+ Plasmide};
\node[box, right=of start, align=center] (cut){Enzyme de restriction\\\textbf{Coupure}};
\node[box, right=of cut, fill=poporangeL, align=center] (lig){ADN Ligase\\\textbf{Soudure}};
\node[box, right=of lig, fill=popgreenL, align=center] (transf){Introduction dans la bactérie\\\textbf{Transfert (Transformation)}};
\node[box, right=of transf, fill=poppurpleL, align=center] (expr){Culture en fermenteur\\\textbf{Expression du gène}};
\node[box, right=of expr, fill=poppinkL, align=center] (end){Protéine d'intérêt\\ (ex: Insuline)};

\draw[arrow] (start) -- (cut) node[midway, above] {Cible : Gène d'intérêt};
\draw[arrow] (cut) -- (lig) node[midway, above] {Fragments à bouts collants};
\draw[arrow] (lig) -- (transf) node[midway, above] {Plasmide recombinant};
\draw[arrow] (transf) -- (expr) node[midway, above] {Bactérie transformée};
\draw[arrow] (expr) -- (end) node[midway, above] {Récolte \& Purification};
\end{tikzpicture}
\legende{Flux chronologique des étapes de la transgénèse bactérienne (production d'insuline). Chaque étape est conditionnée par la réussite de la précédente.}
\end{popfigure}
\item \textbf{Nature de chaque étape :}
\begin{center}
\begin{tabularx}{\linewidth}{|l|X|}\hline
\rowcolor{popblueL} \textbf{Étape} & \textbf{Nature} \\ \hline
Enzyme de restriction : libération du gène et ouverture du plasmide & \textbf{Coupure} \\ \hline
Ligase : formation du plasmide recombinant & \textbf{Soudure} \\ \hline
Introduction dans la bactérie (bactérie transformée) & \textbf{Transfert} \\ \hline
Production de la protéine d'intérêt & \textbf{Expression du gène} \\ \hline
\end{tabularx}
\end{center}
\end{enumerate}
\begin{attention}[Point de vigilance]
Dans un schéma de transgénèse, vérifie toujours que la \textbf{coupure précède la soudure}, que la \textbf{soudure précède le transfert}, et que l'\textbf{expression} est la dernière étape. Une étiquette « ligase » placée avant « enzyme de restriction » est une erreur logique rédhibitoire.
\end{attention}
\end{corrige}

\begin{corrige}[Exercice 9 — Le maïs Bt en Afrique du Sud]
\begin{enumerate}
\item \textbf{Problème agricole :} les chenilles de la pyrale, insecte ravageur, détruisent une partie importante des récoltes de maïs, ce qui réduit les rendements et oblige les agriculteurs à utiliser massivement des insecticides chimiques coûteux et polluants.
\item \textbf{Éléments de la transgénèse :}
\begin{itemize}
\item Gène transféré : le gène codant la \textbf{toxine Bt} (gène \emph{cry}).
\item Organisme donneur : la bactérie \emph{Bacillus thuringiensis}.
\item Organisme receveur : le maïs (\emph{Zea mays}).
\item Caractère nouveau obtenu : la \textbf{résistance aux chenilles de la pyrale} (le maïs produit lui-même la protéine insecticide).
\end{itemize}
\item \textbf{Obtention du maïs transgénique en 4 étapes :}
\begin{enumerate}
\item \textbf{Isolement et coupure :} on extrait le gène \emph{cry} de la bactérie et on ouvre le vecteur (plasmide) avec la \textbf{même enzyme de restriction}, afin d'obtenir des extrémités complémentaires.
\item \textbf{Soudure :} on insère le gène \emph{cry} dans le plasmide ouvert grâce à l'\textbf{ADN ligase} $\rightarrow$ obtention d'un \textbf{ADN recombinant}.
\item \textbf{Transfert :} on introduit le gène recombinant dans des cellules de maïs (via une bactérie vecteur de type \emph{Agrobacterium tumefaciens} ou par bombardement de particules).
\item \textbf{Sélection et régénération :} on sélectionne les cellules de maïs ayant intégré le transgène (marqueur de sélection) et on régénère des plants entiers de maïs Bt, qui expriment la toxine et transmettent le caractère à leur descendance.
\end{enumerate}
\item \textbf{Deux avantages pour l'agriculteur (relevés dans le document) :}
\begin{itemize}
\item utilisation \textbf{réduite d'insecticides chimiques} (économies, moindre toxicité pour l'homme et l'environnement) ;
\item \textbf{rendements plus élevés} (moins de pertes dues au ravageur).
\end{itemize}
\item \textbf{Deux risques et mesures de limitation :}
\begin{itemize}
\item \textbf{Risque 1 :} apparition d'insectes \textbf{résistants à la toxine Bt}. \emph{Mesure :} semer des \textbf{parcelles refuges} de maïs non Bt à proximité, afin de maintenir des insectes sensibles qui s'accouplent avec les éventuels résistants et diluent l'allèle de résistance.
\item \textbf{Risque 2 :} \textbf{dissémination du transgène} vers les variétés locales par pollinisation croisée (flux de gènes). \emph{Mesure :} instaurer des \textbf{distances d'isolement} (zones tampons) entre champs OGM et champs conventionnels, ou décaler les périodes de floraison.
\end{itemize}
\end{enumerate}
\textbf{Barème indicatif (10 points) :} problème dégagé (1 pt) ; quatre éléments de la transgénèse correctement identifiés (2 pts) ; quatre étapes ordonnées avec outils nommés (4 pts) ; deux avantages relevés (1 pt) ; deux risques + mesures cohérentes (2 pts).
\begin{attention}[Points de vigilance]
Exigences du correcteur : nommer \textbf{précisément} les outils (enzyme de restriction, ligase, vecteur, marqueur de sélection) et non « on met le gène dans le maïs » ; distinguer clairement \emph{donneur} (fournit le gène) et \emph{receveur} (reçoit le gène) ; pour les mesures, chaque mesure proposée doit répondre au risque cité.
\end{attention}
\end{corrige}

\begin{corrige}[Exercice 10 — Production d'insuline à l'échelle industrielle]
\begin{enumerate}
\item \textbf{Définitions :}
\begin{itemize}
\item \cle{Enzyme de restriction} : enzyme qui reconnaît une séquence spécifique de l'ADN et le coupe au niveau de ce site ; c'est le « ciseau moléculaire » du génie génétique.
\item \cle{Vecteur} : molécule d'ADN (ici un plasmide) capable de pénétrer dans une cellule hôte et de s'y répliquer, utilisée pour transporter un gène étranger.
\item \cle{Organisme génétiquement modifié (OGM)} : organisme dont le patrimoine génétique a été modifié par l'insertion d'un gène étranger (transgénèse), modification impossible à obtenir par croisements naturels.
\end{itemize}
\item \textbf{Protocole en 5 étapes :}
\begin{enumerate}
\item \textbf{Préparation du gène :} on récupère l'ADNc du gène de l'insuline et on le découpe avec l'\textbf{enzyme de restriction} (extrémités cohésives).
\item \textbf{Ouverture du vecteur :} on coupe le plasmide (porteur du gène de résistance à l'antibiotique) avec la \textbf{même enzyme de restriction}.
\item \textbf{Soudure :} on mélange gène et plasmide ouvert en présence d'\textbf{ADN ligase} $\rightarrow$ obtention du \textbf{plasmide recombinant}.
\item \textbf{Transformation :} on introduit le plasmide recombinant dans les bactéries \emph{E. coli} (transfert).
\item \textbf{Sélection et culture :} on étale les bactéries sur milieu contenant l'\textbf{antibiotique} ; seules les bactéries transformées survivent ; on les cultive en fermenteurs et on récolte l'insuline.
\end{enumerate}
\item \textbf{Rôle du gène de résistance (marqueur de sélection) :} le gène de résistance permet aux bactéries porteuses du plasmide de détruire l'antibiotique. Sur un milieu de culture contenant cet antibiotique, \textbf{seules les bactéries transformées} (ayant intégré le plasmide) \textbf{survivent et forment des colonies} ; les bactéries non transformées meurent. Le marqueur permet donc de \textbf{sélectionner} facilement les bactéries d'intérêt parmi des millions de bactéries.
\item \textbf{Justification de l'identité de l'insuline :} le \textbf{code génétique est universel} : un même codon correspond au même acide aminé chez l'être humain et chez la bactérie. La machinerie de traduction bactérienne lit donc le gène humain exactement comme le ferait une cellule humaine, et produit une protéine dont la \textbf{séquence d'acides aminés est strictement identique} à celle de l'insuline humaine.
\item \textbf{Deux avantages par rapport à l'extraction animale :}
\begin{itemize}
\item insuline \textbf{identique à l'insuline humaine} $\rightarrow$ aucune réaction immunitaire ni allergie (contrairement à l'insuline de porc ou de bœuf) ;
\item production \textbf{illimitée, standardisée et moins coûteuse}, indépendante de l'abattage des animaux, et \textbf{sans risque de transmission d'agents infectieux} d'origine animale.
\end{itemize}
\end{enumerate}
\textbf{Barème indicatif (10 points) :} trois définitions exactes (1,5 pt) ; cinq étapes ordonnées avec outils nommés (4 pts) ; rôle du marqueur de sélection expliqué (1,5 pt) ; universalité du code génétique citée et exploitée (1,5 pt) ; deux avantages pertinents (1,5 pt).
\begin{attention}[Points de vigilance]
La justification de la question 4 doit obligatoirement citer l'\textbf{universalité du code génétique} : c'est la propriété fondamentale attendue. Pour la question 3, il ne suffit pas de dire « le gène sert à la sélection » : il faut expliquer le mécanisme (survie des seules bactéries transformées sur milieu + antibiotique). Enfin, l'énoncé précise que l'ADNc est \textbf{sans introns} : sache justifier pourquoi (la bactérie ne réalise pas l'épissage).
\end{attention}
\end{corrige}

\begin{corrige}[Exercice 11 — Débat argumenté : les OGM au Cameroun]
\begin{enumerate}
\item \textbf{Définition et principe :} un OGM est un organisme dont le patrimoine génétique a été modifié par l'insertion d'un ou plusieurs gènes étrangers (transgénèse), grâce aux techniques du génie génétique. Principe général de fabrication : on \textbf{isole} un gène d'intérêt chez un organisme donneur, on l'\textbf{insère} dans un vecteur (enzyme de restriction puis ligase $\rightarrow$ ADN recombinant), on le \textbf{transfère} dans les cellules de l'organisme receveur, puis on \textbf{sélectionne} et on régénère l'organisme modifié qui exprime le nouveau caractère.
\item \textbf{Paragraphe « avantages » (exemple de rédaction) :} Les cultures OGM pourraient apporter des réponses concrètes aux défis agricoles du Cameroun. D'une part, des variétés résistantes aux ravageurs (maïs Bt contre les foreurs de tiges, niébé résistant aux insectes) réduiraient fortement les pertes de récoltes et diminueraient l'usage d'insecticides chimiques, coûteux et dangereux pour la santé des producteurs. D'autre part, des cultures biofortifiées (manioc enrichi en provitamine A, riz enrichi en fer) contribueraient à lutter contre la malnutrition et les carences, fréquentes dans les zones rurales. On peut ajouter l'intérêt de variétés tolérantes à la sécheresse pour les régions septentrionales confrontées à l'irrégularité des pluies.
\item \textbf{Paragraphe « risques » (exemple de rédaction) :} Ces cultures comportent cependant des risques qu'il serait imprudent de nier. Sur le plan environnemental, le transgène peut se disséminer par pollinisation croisée vers les variétés locales ou les espèces sauvages apparentées (flux de gènes), menaçant la biodiversité cultivée, et l'usage prolongé peut sélectionner des insectes ou adventices résistants. Sur le plan sanitaire et socio-économique, des incertitudes persistent sur les effets à long terme (allergénicité éventuelle), et surtout les semences OGM, brevetées et non ressemables, créeraient une dépendance des paysans camerounais vis-à-vis des firmes semencières étrangères, menaçant la souveraineté alimentaire du pays.
\item \textbf{Conclusion argumentée (exemple de rédaction, position conditionnelle) :} Au terme de cette analyse, je défends une position conditionnelle, fondée sur le principe de précaution. Ce principe ne signifie pas l'interdiction systématique, mais l'obligation d'une évaluation rigoureuse des risques avant toute dissémination. Les OGM ne sont ni une solution miracle ni un danger absolu : ce sont des outils dont la valeur dépend de l'usage qui en est fait. Pour le Cameroun, l'intérêt est réel en matière de sécurité alimentaire et de lutte contre la malnutrition. Mais trois conditions me semblent indispensables : d'abord, chaque OGM doit faire l'objet d'une évaluation scientifique indépendante, au cas par cas, conformément à la loi camerounaise sur la biosécurité ; ensuite, les droits des agriculteurs doivent être protégés (coexistence des filières OGM et non OGM, préservation des semences paysannes) ; enfin, la priorité devrait aller aux OGM d'intérêt public développés par la recherche nationale plutôt qu'aux variétés purement commerciales. À ces conditions, le génie génétique peut devenir un allié du développement agricole camerounais ; sans elles, les risques l'emporteraient sur les bénéfices.
\end{enumerate}
\textbf{Barème indicatif (10 points) :} définition exacte de l'OGM + principe de fabrication (2 pts) ; deux avantages pertinents et contextualisés (2 pts) ; deux risques correctement identifiés (2 pts) ; conclusion structurée, position clairement annoncée, argumentée avec le principe de précaution et des connaissances scientifiques (4 pts).
\begin{attention}[Points de vigilance]
Dans un sujet de débat argumenté, le correcteur attend : (1) une \textbf{position clairement annoncée} (favorable, opposée ou conditionnelle — toutes sont acceptables si elles sont argumentées) ; (2) des arguments \textbf{scientifiques} (pas seulement des opinions) ; (3) la mention explicite du \textbf{principe de précaution} correctement défini (évaluer avant de disséminer, et non « interdire par peur ») ; (4) une conclusion cohérente avec les paragraphes précédents. Évite les affirmations non étayées du type « les OGM rendent stérile » ou « les OGM vont nourrir toute l'Afrique ».
\end{attention}
\end{corrige}

\newpage

\coursec{Fiche bilan}

\begin{popfigure}
\begin{tikzpicture}[
  mindmap,
  grow cyclic,
  every node/.style={concept, text=white, font=\small},
  root concept/.append style={concept color=popink, font=\small\bfseries},
  level 1/.append style={level distance=3.4cm, sibling angle=60},
  level 2/.append style={level distance=2.2cm, font=\scriptsize},
]
\node[root concept]{Génie génétique et OGM}
  child[concept color=popblue]{ node {Définitions}
    child { node {OGM : ADN modifié} }
    child { node {Transgénèse} }
  }
  child[concept color=poporange]{ node {Outils}
    child { node {Enzymes de restriction} }
    child { node {Ligase} }
    child { node {Vecteurs : plasmides} }
  }
  child[concept color=popgreen]{ node {Étapes d'obtention}
    child { node {Isolement du gène} }
    child { node {Insertion dans le vecteur} }
    child { node {Transformation + sélection} }
  }
  child[concept color=poppurple]{ node {Applications}
    child { node {Insuline humaine} }
    child { node {Amélioration variétale} }
  }
  child[concept color=poppink]{ node {Avantages}
    child { node {Rendements, résistances} }
    child { node {Médicaments moins chers} }
  }
  child[concept color=popgold]{ node {Risques et débats}
    child { node {Santé, environnement} }
    child { node {Éthique, réglementation} }
  };
\end{tikzpicture}
\legende{Carte mentale de la leçon : le génie génétique, ses outils, ses étapes, ses applications et ses enjeux.}
\end{popfigure}

\begin{bilan}
\textbf{Définitions clés}
\begin{itemize}
  \item \cle{Génie génétique} : ensemble des techniques permettant de modifier l'information génétique (l'ADN) d'un organisme de manière ciblée, en laboratoire.
  \item \cle{OGM} (organisme génétiquement modifié) : organisme dont le patrimoine génétique a été modifié par le génie génétique, souvent par insertion d'un \cle{transgène} (gène étranger) : c'est la \cle{transgénèse}.
  \item \cle{Enzyme de restriction} : enzyme (d'origine bactérienne) qui coupe l'ADN à des sites précis, comme des « ciseaux moléculaires ».
  \item \cle{Ligase} : enzyme qui soude (recolle) des fragments d'ADN entre eux.
  \item \cle{Vecteur de transfert} : support qui transporte le gène d'intérêt dans la cellule hôte (ex. : \cle{plasmide} bactérien, petite molécule d'ADN circulaire).
\end{itemize}

\textbf{Connaissances à maîtriser par c\oe ur}
\begin{center}
\begin{tabularx}{\linewidth}{|l|X|}
\hline
\rowcolor{popblueL} \textbf{Notion} & \textbf{À retenir} \\
\hline
Outils & Restriction (couper) -- Ligase (coller) -- Vecteur (transporter) \\
\hline
Étapes d'obtention d'un OGM & 1. Isolement du gène d'intérêt ; 2. Insertion dans le vecteur (ADN recombinant) ; 3. Introduction dans la cellule hôte (transformation) ; 4. Sélection et multiplication des cellules transformées ; 5. Vérification de l'expression du gène \\
\hline
Insuline humaine & Gène humain de l'insuline inséré dans une bactérie (\emph{E. coli}) qui produit alors l'hormone en grande quantité pour les diabétiques \\
\hline
Amélioration variétale & Plantes résistantes aux insectes (maïs Bt), aux herbicides, à meilleure valeur nutritionnelle (riz doré enrichi en provitamine A) \\
\hline
\end{tabularx}
\end{center}

\textbf{Méthodes express}
\begin{itemize}
  \item \emph{Décrire une manipulation sur schéma} : suivre l'ordre gène $\rightarrow$ vecteur $\rightarrow$ cellule hôte $\rightarrow$ sélection $\rightarrow$ expression, en nommant l'enzyme utilisée à chaque étape.
  \item \emph{Expliquer l'intérêt d'un OGM} : identifier le caractère ajouté, le bénéfice obtenu, et le comparer à la méthode classique (sélection, extraction animale).
  \item \emph{Discuter avantages/risques} : toujours présenter les deux points de vue (santé, environnement, économie, éthique) avant de conclure de façon nuancée.
\end{itemize}
\end{bilan}

\begin{aretenir}[Checklist avant le Probatoire]
\begin{itemize}
  \item[$\square$] Je sais définir le génie génétique, un OGM et la transgénèse.
  \item[$\square$] Je connais le rôle des enzymes de restriction (couper l'ADN à des sites précis).
  \item[$\square$] Je connais le rôle de la ligase (souder les fragments d'ADN).
  \item[$\square$] Je sais ce qu'est un vecteur de transfert et je cite le plasmide comme exemple.
  \item[$\square$] Je sais citer dans l'ordre les cinq étapes de l'obtention d'un OGM.
  \item[$\square$] Je sais décrire un schéma de manipulation de génie génétique étape par étape.
  \item[$\square$] Je sais expliquer comment une bactérie produit de l'insuline humaine.
  \item[$\square$] Je cite au moins deux exemples d'amélioration variétale par transgénèse.
  \item[$\square$] Je sais présenter au moins deux avantages et deux risques des OGM.
  \item[$\square$] Je sais construire une conclusion argumentée et nuancée sur les enjeux des OGM.
\end{itemize}
\end{aretenir}

\findecours

\end{document}
